% Mengeninterne Syntax und Semantik. Keine neuen Mengenaxiome.
% Alle Hilfssaetze stehen als Proofparts beim Existenzsatz.
\newcommand{\BXLVIIIIf}[3]{\mathsf{if}(#1,#2,#3)}
\subsection*{Formelcodes und axiomatische Charakterisierung}
Wie bei den induktiven Mengen in
\FormulaRefAuto{InductiveSetDef}[def] stehen die kennzeichnenden Axiome
eines Strukturbegriffs einzeln in dessen Definition. Die Definition
dient als Einfuehrungsregel: Ihre Axiome werden durch getrennte
Beweiszeilen nachgewiesen. Die anschliessenden Zugriffsaxiome liefern
aus dem Strukturpraedikat jeweils genau eine dieser Bedingungen.
Keine dieser Regeln behauptet die Existenz einer solchen Struktur.
In verbleibenden logischen Formeln werden mehrgliedrige Konjunktionen
rechts geklammert.
Die Mengen \(D,E\) sind feste Parameter dieses Abschnitts; \(M=(D,E)\).
\FormulaDefDeltaK[Kurzschreibweise des Fallterms]
{\BXLVIIIIf{P}{U}{V}\coloneq\IfThenElse{P}{U}{V}}
{B48IfNotation}{\DeltaRow{Aussage}{P}\DeltaRow{Mengen}{U,V}}


Die folgenden Bedingungen charakterisieren eine gegebene Relation.
Sie behaupten noch nicht deren Existenz. Als Codebereich dienen die
bereits konstruierten Klammerungsbaeume aus Band~27.
Gebundene Tupel sind die uebliche Abkuerzung fuer die eindeutige
Produktzerlegung; es werden keine neuen Arten von Variablen eingefuehrt.
Die Strukturbegriffe sind Praedikate ueber Mengen; die uebrigen
Kurzzeichen bezeichnen Mengen und Mengenoperationen.

\FormulaDefDeltaK[Marken und Codebereich]
{\begin{gathered}
 t_{\rm e}\coloneq0,\quad t_{\rm m}\coloneq\Succ(t_{\rm e}),\quad
 t_{\rm n}\coloneq\Succ(t_{\rm m}),\\
 t_{\rm c}\coloneq\Succ(t_{\rm n}),\quad
 t_{\rm q}\coloneq\Succ(t_{\rm c}),\quad
 t_{\rm b}\coloneq\Succ(t_{\rm q}),\\
 \Sigma\coloneq\mathbb N\times(\mathbb N\times\mathbb N),\qquad
 \mathcal T\coloneq\BracketTreeSet{\Sigma},\\
 \mathsf l(c)\coloneq\LeafTree{\Sigma}{c},\qquad
 \mathsf n(p,q)\coloneq\NodeTree{\Sigma}{p}{q}
\end{gathered}}
{B48CodeCarrier}{\DeltaRow{Mengen}{c\in\Sigma,\quad p,q\in\mathcal T}}

\FormulaDefDeltaK[Formelkonstruktoren]
{\begin{gathered}
 \begin{aligned}
 \mathsf e(i,j)&\coloneq\mathsf l((t_{\rm e},(i,j))),&
 \mathsf m(i,j)&\coloneq\mathsf l((t_{\rm m},(i,j))),\\
 \mathsf b&\coloneq\mathsf l((t_{\rm b},(0,0))),&
 \mathsf n_0&\coloneq\mathsf l((t_{\rm n},(0,0))),\\
 \mathsf c_0&\coloneq\mathsf l((t_{\rm c},(0,0))),&
 \mathsf q_i&\coloneq\mathsf l((t_{\rm q},(i,0))),\\
 \mathsf{neg}(p)&\coloneq\mathsf n(\mathsf n_0,p),&
 \mathsf{and}(p,q)&\coloneq\mathsf n(\mathsf n(\mathsf c_0,p),q),\\
 \mathsf{ex}_i(p)&\coloneq\mathsf n(\mathsf q_i,p)
 \end{aligned}\\
 \begin{aligned}
 \Sigma_{\rm at}\coloneq\{(k,(i,j))\in\Sigma\mid{}&
 k=t_{\rm e}\lor k=t_{\rm m}\\
 &{}\lor(k=t_{\rm b}\land i=0\land j=0)\}.
 \end{aligned}
\end{gathered}}
{B48FormulaConstructors}{\DeltaRow{Indizes}{i,j\in\mathbb N}
 \DeltaRow{Baumcodes}{p,q\in\mathcal T}}

\FormulaDefDeltaK[Formelabgeschlossene Menge]
{\mathsf{Cl}(X)}{B48FormulaClosure}{
 \DeltaRow{Mengen}{X,p,c,q,i}
 \DeltaRow{\textbf{Axiome}}{}
 \DeltaRow{Baumcodes}{X\subseteq\mathcal T}
 \DeltaRow{Atome}{\forall c\in\Sigma_{\rm at}\;(\mathsf l(c)\in X)}
 \DeltaRow{Negation}{\forall p\in X\;(\mathsf{neg}(p)\in X)}
 \DeltaRow{Konjunktion}{\begin{aligned}[t]&\forall p\in X\,\forall q\in X\\&\quad(\mathsf{and}(p,q)\in X)\end{aligned}}
 \DeltaRow{Existenzquantor}{\begin{aligned}[t]&\forall i\in\mathbb N\,\forall p\in X\\&\quad(\mathsf{ex}_i(p)\in X)\end{aligned}}
 \DeltaRow{\textbf{Neues Symbol}}{}
 \DeltaPrem{Formelabschluss}{\mathsf{Cl}(X)}
}

\FormulaAxiomDeltaK[Traeger einer formelabgeschlossenen Menge]
{\mathsf{Cl}(X)\vdash X\subseteq\mathcal T}
{B48ClDomainAxiom}{\DeltaRow{Mengen}{X,p,c,q,i}}

\FormulaAxiomDeltaK[Atomabschluss]
{\mathsf{Cl}(X)\vdash \forall c\in\Sigma_{\rm at}\;(\mathsf l(c)\in X)}
{B48ClAtomsAxiom}{\DeltaRow{Mengen}{X,p,c,q,i}}

\FormulaAxiomDeltaK[Abschluss unter Negation]
{\mathsf{Cl}(X)\vdash \forall p\in X\;(\mathsf{neg}(p)\in X)}
{B48ClNegAxiom}{\DeltaRow{Mengen}{X,p,c,q,i}}

\FormulaAxiomDeltaK[Abschluss unter Konjunktion]
{\begin{gathered}\mathsf{Cl}(X)\vdash{}\\
 \forall p\in X\,\forall q\in X\;(\mathsf{and}(p,q)\in X)\end{gathered}}
{B48ClAndAxiom}{\DeltaRow{Mengen}{X,p,c,q,i}}

\FormulaAxiomDeltaK[Abschluss unter Existenzquantifikation]
{\begin{gathered}\mathsf{Cl}(X)\vdash{}\\
 \forall i\in\mathbb N\,\forall p\in X\;(\mathsf{ex}_i(p)\in X)\end{gathered}}
{B48ClExistsAxiom}{\DeltaRow{Mengen}{X,p,c,q,i}}

\FormulaDefDeltaK[Formelmenge]
{\mathcal F\coloneq
 \{p\in\mathcal T\mid\forall X(\mathsf{Cl}(X)\to p\in X)\}}
{B48FormulaSet}{\DeltaRow{Mengen}{p,X}}

Dies ist eine beschraenkte Aussonderung aus einer bereits vorhandenen
Menge, keine rekursive Existenzbehauptung.
Die Codes \(\mathsf e(i,j)\), \(\mathsf m(i,j)\), \(\mathsf b\),
\(\mathsf{neg}(p)\), \(\mathsf{and}(p,q)\) und \(\mathsf{ex}_i(p)\)
codieren jeweils \(v_i=v_j\), \(v_i\in v_j\), \(\bot\), \(\neg p\),
\(p\land q\) und \(\exists v_i\,p\).
Innerhalb von \(\models\) werden spaeter die lesbaren Formeln statt
dieser Codes geschrieben. \(\mathsf n_0,\mathsf c_0,\mathsf q_i\)
sind Hilfsbaeume, nicht zusaetzliche atomare Formeln.

\FormulaDefDeltaK[Mengenstruktur]
{\mathsf{Str}(M)}{B48SetStructureDef}{
 \DeltaRow{Mengen}{M,D,E}
 \DeltaRow{\textbf{Axiome}}{}
 \DeltaRow{Strukturpaar}{M=(D,E)}
 \DeltaRow{Grundmenge}{D\neq\varnothing}
 \DeltaRow{Relation}{E\subseteq D\times D}
 \DeltaRow{\textbf{Neues Symbol}}{}
 \DeltaPrem{Mengenstruktur}{\mathsf{Str}(M)}
}

\FormulaAxiomDeltaK[Paar einer Mengenstruktur]
{\mathsf{Str}(M)\vdash M=(D,E)}
{B48StrPairAxiom}{\DeltaRow{Mengen}{M,D,E}}

\FormulaAxiomDeltaK[Nichtleere Grundmenge]
{\mathsf{Str}(M)\vdash D\neq\varnothing}
{B48StrDomainAxiom}{\DeltaRow{Mengen}{M,D,E}}

\FormulaAxiomDeltaK[Interpretation der Elementrelation]
{\mathsf{Str}(M)\vdash E\subseteq D\times D}
{B48StrRelationAxiom}{\DeltaRow{Mengen}{M,D,E}}

\FormulaDefDeltaK[Belegungen und Belegungswechsel]
{\begin{aligned}
 A_D&\coloneq\{s\in\powerset(\mathbb N\times D)\mid s\colon\mathbb N\to D\},\\
 s[i\mapsto a]&\coloneq
 \{(j,b)\in\mathbb N\times D\mid
       b=\BXLVIIIIf{j=i}{a}{s(j)}\}.
\end{aligned}}
{B48Assignments}{\DeltaRow{Mengen}{D,s,i,a,j,b}}

\FormulaDefDeltaK[Atomare Wertemengen und Quantorenoperator]
{\begin{aligned}
 B_M((k,(i,j)))\coloneq{}&
 \mathsf{if}\bigl(k=t_{\rm e},\{s\in A_D\mid s(i)=s(j)\},\\
 &\qquad\BXLVIIIIf{k=t_{\rm m}}
     {\{s\in A_D\mid(s(i),s(j))\in E\}}{\varnothing}\bigr),\\
 Q_i(U)\coloneq{}&\{s\in A_D\mid\exists a\in D\;(s[i\mapsto a]\in U)\},\\
 Q(t,U)\coloneq{}&\{s\in A_D\mid
       \exists i\in\mathbb N\;(t=\mathsf q_i\land
                              \exists a\in D\;(s[i\mapsto a]\in U))\},\\
 R_S(p)\coloneq{}&\{s\in A_D\mid(p,s)\in S\}.
\end{aligned}}
{B48SemanticTerms}{\DeltaRow{Mengen}{M=(D,E),\ (k,(i,j))\in\Sigma,\ t,U,S,p}}

\FormulaDefDeltaK[Atomklausel der Erfuellung]
{C_{\rm at}(M,S)\coloneq \forall c\in\Sigma_{\rm at}\;(R_S(\mathsf l(c))=B_M(c))}
{B48SatisfactionAtomsDef}{\DeltaRow{Mengen}{M=(D,E),S,p,q,c,i}}

\FormulaDefDeltaK[Negationsklausel der Erfuellung]
{C_{\rm n}(M,S)\coloneq \forall p\in\mathcal F\;(R_S(\mathsf{neg}(p))=A_D\setminus R_S(p))}
{B48SatisfactionNegDef}{\DeltaRow{Mengen}{M=(D,E),S,p,q,c,i}}

\FormulaDefDeltaK[Konjunktionsklausel der Erfuellung]
{C_{\rm c}(M,S)\coloneq \begin{aligned}[t]&\forall p\in\mathcal F\,\forall q\in\mathcal F\\&\quad(R_S(\mathsf{and}(p,q))=R_S(p)\cap R_S(q))\end{aligned}}
{B48SatisfactionAndDef}{\DeltaRow{Mengen}{M=(D,E),S,p,q,c,i}}

\FormulaDefDeltaK[Quantorenklausel der Erfuellung]
{C_{\rm q}(M,S)\coloneq \begin{aligned}[t]&\forall i\in\mathbb N\,\forall p\in\mathcal F\\&\quad(R_S(\mathsf{ex}_i(p))=Q_i(R_S(p)))\end{aligned}}
{B48SatisfactionExistsDef}{\DeltaRow{Mengen}{M=(D,E),S,p,q,c,i}}

\FormulaDefDeltaK[Erfuellungsrelation]
{\mathsf{Erf}(M,S)}{B48SatisfactionStructure}{
 \DeltaRow{Mengen}{M=(D,E),S}
 \DeltaRow{\textbf{Axiome}}{}
 \DeltaRow{Struktur}{\mathsf{Str}(M)}
 \DeltaRow{Relationstraeger}{S\subseteq\mathcal F\times A_D}
 \DeltaRow{Atome}{C_{\rm at}(M,S)}[\FormulaRefAuto{B48SatisfactionAtomsDef}[def]]
 \DeltaRow{Negation}{C_{\rm n}(M,S)}[\FormulaRefAuto{B48SatisfactionNegDef}[def]]
 \DeltaRow{Konjunktion}{C_{\rm c}(M,S)}[\FormulaRefAuto{B48SatisfactionAndDef}[def]]
 \DeltaRow{Existenzquantor}{C_{\rm q}(M,S)}[\FormulaRefAuto{B48SatisfactionExistsDef}[def]]
 \DeltaRow{\textbf{Neues Symbol}}{}
 \DeltaPrem{Erfuellung}{\mathsf{Erf}(M,S)}
}

\FormulaAxiomDeltaK[Struktur einer Erfuellungsrelation]
{\mathsf{Erf}(M,S)\vdash \mathsf{Str}(M)}
{B48ErfStructureAxiom}{\DeltaRow{Mengen}{M=(D,E),S}}

\FormulaAxiomDeltaK[Traeger einer Erfuellungsrelation]
{\mathsf{Erf}(M,S)\vdash S\subseteq\mathcal F\times A_D}
{B48ErfSubsetAxiom}{\DeltaRow{Mengen}{M=(D,E),S}}

\FormulaAxiomDeltaK[Atomklausel der Erfuellung]
{\begin{gathered}\mathsf{Erf}(M,S)\vdash{}\\
 \forall c\in\Sigma_{\rm at}\;(R_S(\mathsf l(c))=B_M(c))\end{gathered}}
{B48ErfAtomsAxiom}{\DeltaRow{Mengen}{M=(D,E),S}}

\FormulaAxiomDeltaK[Negationsklausel der Erfuellung]
{\begin{gathered}\mathsf{Erf}(M,S)\vdash{}\\
 \forall p\in\mathcal F\;(R_S(\mathsf{neg}(p))=A_D\setminus R_S(p))\end{gathered}}
{B48ErfNegAxiom}{\DeltaRow{Mengen}{M=(D,E),S}}

\FormulaAxiomDeltaK[Konjunktionsklausel der Erfuellung]
{\begin{gathered}\mathsf{Erf}(M,S)\vdash{}\\
 \forall p\in\mathcal F\,\forall q\in\mathcal F\;(R_S(\mathsf{and}(p,q))=R_S(p)\cap R_S(q))\end{gathered}}
{B48ErfAndAxiom}{\DeltaRow{Mengen}{M=(D,E),S}}

\FormulaAxiomDeltaK[Quantorenklausel der Erfuellung]
{\begin{gathered}\mathsf{Erf}(M,S)\vdash{}\\
 \forall i\in\mathbb N\,\forall p\in\mathcal F\;(R_S(\mathsf{ex}_i(p))=Q_i(R_S(p)))\end{gathered}}
{B48ErfExistsAxiom}{\DeltaRow{Mengen}{M=(D,E),S}}

Die vier Klauselpraedikate wurden jeweils durch eine eigene Formel
definiert. Zusammen mit den Bedingungen an Struktur und Relation bilden
sie die sechs Axiome des Begriffs \(\mathsf{Erf}(M,S)\).
Aus dem Nachweis dieser sechs Bedingungen wird der Begriff eingefuehrt;
aus dem Begriff wird jede Bedingung durch ihr eigenes Zugriffsaxiom
gewonnen. Die Existenz einer solchen Relation folgt erst aus dem
nachstehenden Konstruktionsbeweis.

\FormulaDefDeltaK[Totaler Auswertungsoperator auf Hilfsbaeumen]
{\begin{aligned}
 J(t)&\coloneq\exists p\in\mathcal T\;(t=\mathsf n(\mathsf c_0,p)),\\
 H_{2,M}(t,U,V)&\coloneq\BXLVIIIIf{J(t)}{U\cap V}{Q(t,V)},\\
 H_{1,M}(t,U,V)&\coloneq\BXLVIIIIf{t=\mathsf c_0}{V}{H_{2,M}(t,U,V)},\\
 H_M(t,U,V)&\coloneq\BXLVIIIIf{t=\mathsf n_0}{A_D\setminus V}{H_{1,M}(t,U,V)},\\
 Y_D&\coloneq\mathcal T\times\powerset(A_D),\\
 f_M&\coloneq\{(c,z)\in\Sigma\times Y_D\mid z=(\mathsf l(c),B_M(c))\},\\
 g_M&\coloneq
 \{(((p,U),(q,V)),z)\in(Y_D\times Y_D)\times Y_D\mid\\
 &\hspace{42mm}z=(\mathsf n(p,q),H_M(p,U,V))\}.
\end{aligned}}
{B48EvaluationOperators}{\DeltaRow{Mengen}{M=(D,E),t,U,V,c,z,p,q}}

Hier wird auch auf Hilfsbaeumen ausgewertet. Der erste Bestandteil eines
Auswertungswerts speichert den Baumcode selbst. Dadurch kann der zweite
Bestandteil die drei logischen Konstruktoren unterscheiden.
Diese Hilfsauswertung erweitert nicht die Objektlogik.

\FormulaDefDeltaK[Rekursive Auswertung]
{\mathsf{Rec}(e,M)}{B48RecursiveEvaluationDef}{
 \DeltaRow{Mengen}{M=(D,E),e,p,q,c}
 \DeltaRow{\textbf{Axiome}}{}
 \DeltaRow{Auswertung}{e\colon\mathcal T\to Y_D}
 \DeltaRow{Blaetter}{\forall c\in\Sigma\;(e(\mathsf l(c))=f_M(c))}
 \DeltaRow{Knoten}{\begin{aligned}[t]&\forall p\in\mathcal T\,\forall q\in\mathcal T\\&\quad(e(\mathsf n(p,q))=g_M(e(p),e(q)))\end{aligned}}
 \DeltaRow{\textbf{Neues Symbol}}{}
 \DeltaPrem{Auswertung}{\mathsf{Rec}(e,M)}
}

\FormulaAxiomDeltaK[Typ einer rekursiven Auswertung]
{\mathsf{Rec}(e,M)\vdash e\colon\mathcal T\to Y_D}
{B48RecTypeAxiom}{\DeltaRow{Mengen}{M=(D,E),e,p,q,c}}

\FormulaAxiomDeltaK[Blattklausel der rekursiven Auswertung]
{\mathsf{Rec}(e,M)\vdash \forall c\in\Sigma\;(e(\mathsf l(c))=f_M(c))}
{B48RecLeafAxiom}{\DeltaRow{Mengen}{M=(D,E),e,p,q,c}}

\FormulaAxiomDeltaK[Knotenklausel der rekursiven Auswertung]
{\begin{gathered}\mathsf{Rec}(e,M)\vdash{}\\
 \forall p\in\mathcal T\,\forall q\in\mathcal T\;(e(\mathsf n(p,q))=g_M(e(p),e(q)))\end{gathered}}
{B48RecNodeAxiom}{\DeltaRow{Mengen}{M=(D,E),e,p,q,c}}

\FormulaDefDeltaK[Auslesen und Kandidat]
{\begin{aligned}
 V_A(z,p)&\coloneq\{s\in A\mid
           \exists U\in\powerset(A)\;(z=(p,U)\land s\in U)\},\\
 \beta_e(p)&\coloneq V_{A_D}(e(p),p),\\
 S_e&\coloneq\{(p,s)\in\mathcal F\times A_D\mid s\in\beta_e(p)\}.
\end{aligned}}
{B48EvaluationCandidate}{\DeltaRow{Mengen}{M=(D,E),e,p,z,A,s,U}}

\FormulaDefDeltaK[Gleichheitstraeger zweier Relationen]
{K(S,T)\coloneq\{p\in\mathcal F\mid R_S(p)=R_T(p)\}}
{B48EqualityCarrier}{\DeltaRow{Mengen}{S,T,p}}

\FormulaThmDeltaK[Existenz und Eindeutigkeit der Erfuellungsrelation]
{\mathsf{Str}(M)\vdash\exists!S\;\mathsf{Erf}(M,S)}
{B48SatisfactionExists}{\DeltaRow{Struktur}{M=(D,E)}
 \DeltaRow{Mengen}{S,T,e,p,q,c,i,j,U,V,X,z,s,a}
 \DeltaRow{Praedikate}{P}}
\noindent
Die folgenden Proofparts bilden den Beweis. Jeder Hilfssatz wird vor
seiner Verwendung hergeleitet. In Zeugenbeweisen sind die gewaelten
Mengen frisch fuer die ausserhalb des jeweiligen Teilbeweises offenen
Voraussetzungen.
\begin{tabproofsplitwide}
\par\medskip\noindent\textbf{A. Konkrete Marken und Baumcodes}\par
\proofpartwideindR[Typisierung der Marke e]{t_{\rm e}\in\mathbb N}{t_{\rm e}\in\mathbb N}[B48TagType1]
 \proofstepwidestar[]{0\in\mathbb N}{\FormulaRefAuto{PeanoZero}[ax]}
 \proofstepwidestar[]{t_{\rm e}\in\mathbb N}{\FormulaRefAuto{B48CodeCarrier}[def]{1}}
\closeproofpartwideindR

\proofpartwideindR[Typisierung der Marke m]{t_{\rm m}\in\mathbb N}{t_{\rm m}\in\mathbb N}[B48TagType2]
 \proofstepwidestar[]{t_{\rm e}\in\mathbb N}{\FormulaRefAuto{B48TagType1}[thm]}
 \proofstepwidestar[]{\Succ(t_{\rm e})\in\mathbb N}{\FormulaRefAuto{PeanoSuccClosure}[ax]{1}}
 \proofstepwidestar[]{t_{\rm m}\in\mathbb N}{\FormulaRefAuto{B48CodeCarrier}[def]{2}}
\closeproofpartwideindR

\proofpartwideindR[Typisierung der Marke n]{t_{\rm n}\in\mathbb N}{t_{\rm n}\in\mathbb N}[B48TagType3]
 \proofstepwidestar[]{t_{\rm m}\in\mathbb N}{\FormulaRefAuto{B48TagType2}[thm]}
 \proofstepwidestar[]{\Succ(t_{\rm m})\in\mathbb N}{\FormulaRefAuto{PeanoSuccClosure}[ax]{1}}
 \proofstepwidestar[]{t_{\rm n}\in\mathbb N}{\FormulaRefAuto{B48CodeCarrier}[def]{2}}
\closeproofpartwideindR

\proofpartwideindR[Typisierung der Marke c]{t_{\rm c}\in\mathbb N}{t_{\rm c}\in\mathbb N}[B48TagType4]
 \proofstepwidestar[]{t_{\rm n}\in\mathbb N}{\FormulaRefAuto{B48TagType3}[thm]}
 \proofstepwidestar[]{\Succ(t_{\rm n})\in\mathbb N}{\FormulaRefAuto{PeanoSuccClosure}[ax]{1}}
 \proofstepwidestar[]{t_{\rm c}\in\mathbb N}{\FormulaRefAuto{B48CodeCarrier}[def]{2}}
\closeproofpartwideindR

\proofpartwideindR[Typisierung der Marke q]{t_{\rm q}\in\mathbb N}{t_{\rm q}\in\mathbb N}[B48TagType5]
 \proofstepwidestar[]{t_{\rm c}\in\mathbb N}{\FormulaRefAuto{B48TagType4}[thm]}
 \proofstepwidestar[]{\Succ(t_{\rm c})\in\mathbb N}{\FormulaRefAuto{PeanoSuccClosure}[ax]{1}}
 \proofstepwidestar[]{t_{\rm q}\in\mathbb N}{\FormulaRefAuto{B48CodeCarrier}[def]{2}}
\closeproofpartwideindR

\proofpartwideindR[Typisierung der Marke b]{t_{\rm b}\in\mathbb N}{t_{\rm b}\in\mathbb N}[B48TagType6]
 \proofstepwidestar[]{t_{\rm q}\in\mathbb N}{\FormulaRefAuto{B48TagType5}[thm]}
 \proofstepwidestar[]{\Succ(t_{\rm q})\in\mathbb N}{\FormulaRefAuto{PeanoSuccClosure}[ax]{1}}
 \proofstepwidestar[]{t_{\rm b}\in\mathbb N}{\FormulaRefAuto{B48CodeCarrier}[def]{2}}
\closeproofpartwideindR

\proofpartwideindR[Trennung der Marken e und m]{t_{\rm e}\neq t_{\rm m}}{t_{\rm e}\neq t_{\rm m}}[B48TagNe01]
 \proofstepwidestar[]{t_{\rm e}\in\mathbb N}{\FormulaRefAuto{B48TagType1}[thm]}
 \proofstepwidestar[]{0\neq\Succ(t_{\rm e})}{\FormulaRefAuto{PeanoZeroNeSuccessor}[thm]{1}}
 \proofstepwidestar[]{t_{\rm e}\neq t_{\rm m}}{\FormulaRefAuto{B48CodeCarrier}[def]{2}}
\closeproofpartwideindR

\proofpartwideindR[Trennung der Marken e und n]{t_{\rm e}\neq t_{\rm n}}{t_{\rm e}\neq t_{\rm n}}[B48TagNe02]
 \proofstepwidestar[]{t_{\rm m}\in\mathbb N}{\FormulaRefAuto{B48TagType2}[thm]}
 \proofstepwidestar[]{0\neq\Succ(t_{\rm m})}{\FormulaRefAuto{PeanoZeroNeSuccessor}[thm]{1}}
 \proofstepwidestar[]{t_{\rm e}\neq t_{\rm n}}{\FormulaRefAuto{B48CodeCarrier}[def]{2}}
\closeproofpartwideindR

\proofpartwideindR[Trennung der Marken e und c]{t_{\rm e}\neq t_{\rm c}}{t_{\rm e}\neq t_{\rm c}}[B48TagNe03]
 \proofstepwidestar[]{t_{\rm n}\in\mathbb N}{\FormulaRefAuto{B48TagType3}[thm]}
 \proofstepwidestar[]{0\neq\Succ(t_{\rm n})}{\FormulaRefAuto{PeanoZeroNeSuccessor}[thm]{1}}
 \proofstepwidestar[]{t_{\rm e}\neq t_{\rm c}}{\FormulaRefAuto{B48CodeCarrier}[def]{2}}
\closeproofpartwideindR

\proofpartwideindR[Trennung der Marken e und q]{t_{\rm e}\neq t_{\rm q}}{t_{\rm e}\neq t_{\rm q}}[B48TagNe04]
 \proofstepwidestar[]{t_{\rm c}\in\mathbb N}{\FormulaRefAuto{B48TagType4}[thm]}
 \proofstepwidestar[]{0\neq\Succ(t_{\rm c})}{\FormulaRefAuto{PeanoZeroNeSuccessor}[thm]{1}}
 \proofstepwidestar[]{t_{\rm e}\neq t_{\rm q}}{\FormulaRefAuto{B48CodeCarrier}[def]{2}}
\closeproofpartwideindR

\proofpartwideindR[Trennung der Marken e und b]{t_{\rm e}\neq t_{\rm b}}{t_{\rm e}\neq t_{\rm b}}[B48TagNe05]
 \proofstepwidestar[]{t_{\rm q}\in\mathbb N}{\FormulaRefAuto{B48TagType5}[thm]}
 \proofstepwidestar[]{0\neq\Succ(t_{\rm q})}{\FormulaRefAuto{PeanoZeroNeSuccessor}[thm]{1}}
 \proofstepwidestar[]{t_{\rm e}\neq t_{\rm b}}{\FormulaRefAuto{B48CodeCarrier}[def]{2}}
\closeproofpartwideindR

\proofpartwideindR[Trennung der Marken m und n]{t_{\rm m}\neq t_{\rm n}}{t_{\rm m}\neq t_{\rm n}}[B48TagNe12]
 \proofstepwidestar[]{t_{\rm e}\in\mathbb N}{\FormulaRefAuto{B48TagType1}[thm]}
 \proofstepwidestar[]{t_{\rm m}\in\mathbb N}{\FormulaRefAuto{B48TagType2}[thm]}
 \proofstepwidestar[]{t_{\rm e}\neq t_{\rm m}}{\FormulaRefAuto{B48TagNe01}[thm]}
 \proofstepwidestar[]{\Succ(t_{\rm e})\neq\Succ(t_{\rm m})}{\FormulaRefAuto{PeanoSuccPreservesNe}[thm]{1,2,3}}
 \proofstepwidestar[]{t_{\rm m}\neq t_{\rm n}}{\FormulaRefAuto{B48CodeCarrier}[def]{4}}
\closeproofpartwideindR

\proofpartwideindR[Trennung der Marken m und c]{t_{\rm m}\neq t_{\rm c}}{t_{\rm m}\neq t_{\rm c}}[B48TagNe13]
 \proofstepwidestar[]{t_{\rm e}\in\mathbb N}{\FormulaRefAuto{B48TagType1}[thm]}
 \proofstepwidestar[]{t_{\rm n}\in\mathbb N}{\FormulaRefAuto{B48TagType3}[thm]}
 \proofstepwidestar[]{t_{\rm e}\neq t_{\rm n}}{\FormulaRefAuto{B48TagNe02}[thm]}
 \proofstepwidestar[]{\Succ(t_{\rm e})\neq\Succ(t_{\rm n})}{\FormulaRefAuto{PeanoSuccPreservesNe}[thm]{1,2,3}}
 \proofstepwidestar[]{t_{\rm m}\neq t_{\rm c}}{\FormulaRefAuto{B48CodeCarrier}[def]{4}}
\closeproofpartwideindR

\proofpartwideindR[Trennung der Marken m und q]{t_{\rm m}\neq t_{\rm q}}{t_{\rm m}\neq t_{\rm q}}[B48TagNe14]
 \proofstepwidestar[]{t_{\rm e}\in\mathbb N}{\FormulaRefAuto{B48TagType1}[thm]}
 \proofstepwidestar[]{t_{\rm c}\in\mathbb N}{\FormulaRefAuto{B48TagType4}[thm]}
 \proofstepwidestar[]{t_{\rm e}\neq t_{\rm c}}{\FormulaRefAuto{B48TagNe03}[thm]}
 \proofstepwidestar[]{\Succ(t_{\rm e})\neq\Succ(t_{\rm c})}{\FormulaRefAuto{PeanoSuccPreservesNe}[thm]{1,2,3}}
 \proofstepwidestar[]{t_{\rm m}\neq t_{\rm q}}{\FormulaRefAuto{B48CodeCarrier}[def]{4}}
\closeproofpartwideindR

\proofpartwideindR[Trennung der Marken m und b]{t_{\rm m}\neq t_{\rm b}}{t_{\rm m}\neq t_{\rm b}}[B48TagNe15]
 \proofstepwidestar[]{t_{\rm e}\in\mathbb N}{\FormulaRefAuto{B48TagType1}[thm]}
 \proofstepwidestar[]{t_{\rm q}\in\mathbb N}{\FormulaRefAuto{B48TagType5}[thm]}
 \proofstepwidestar[]{t_{\rm e}\neq t_{\rm q}}{\FormulaRefAuto{B48TagNe04}[thm]}
 \proofstepwidestar[]{\Succ(t_{\rm e})\neq\Succ(t_{\rm q})}{\FormulaRefAuto{PeanoSuccPreservesNe}[thm]{1,2,3}}
 \proofstepwidestar[]{t_{\rm m}\neq t_{\rm b}}{\FormulaRefAuto{B48CodeCarrier}[def]{4}}
\closeproofpartwideindR

\proofpartwideindR[Trennung der Marken n und c]{t_{\rm n}\neq t_{\rm c}}{t_{\rm n}\neq t_{\rm c}}[B48TagNe23]
 \proofstepwidestar[]{t_{\rm m}\in\mathbb N}{\FormulaRefAuto{B48TagType2}[thm]}
 \proofstepwidestar[]{t_{\rm n}\in\mathbb N}{\FormulaRefAuto{B48TagType3}[thm]}
 \proofstepwidestar[]{t_{\rm m}\neq t_{\rm n}}{\FormulaRefAuto{B48TagNe12}[thm]}
 \proofstepwidestar[]{\Succ(t_{\rm m})\neq\Succ(t_{\rm n})}{\FormulaRefAuto{PeanoSuccPreservesNe}[thm]{1,2,3}}
 \proofstepwidestar[]{t_{\rm n}\neq t_{\rm c}}{\FormulaRefAuto{B48CodeCarrier}[def]{4}}
\closeproofpartwideindR

\proofpartwideindR[Trennung der Marken n und q]{t_{\rm n}\neq t_{\rm q}}{t_{\rm n}\neq t_{\rm q}}[B48TagNe24]
 \proofstepwidestar[]{t_{\rm m}\in\mathbb N}{\FormulaRefAuto{B48TagType2}[thm]}
 \proofstepwidestar[]{t_{\rm c}\in\mathbb N}{\FormulaRefAuto{B48TagType4}[thm]}
 \proofstepwidestar[]{t_{\rm m}\neq t_{\rm c}}{\FormulaRefAuto{B48TagNe13}[thm]}
 \proofstepwidestar[]{\Succ(t_{\rm m})\neq\Succ(t_{\rm c})}{\FormulaRefAuto{PeanoSuccPreservesNe}[thm]{1,2,3}}
 \proofstepwidestar[]{t_{\rm n}\neq t_{\rm q}}{\FormulaRefAuto{B48CodeCarrier}[def]{4}}
\closeproofpartwideindR

\proofpartwideindR[Trennung der Marken n und b]{t_{\rm n}\neq t_{\rm b}}{t_{\rm n}\neq t_{\rm b}}[B48TagNe25]
 \proofstepwidestar[]{t_{\rm m}\in\mathbb N}{\FormulaRefAuto{B48TagType2}[thm]}
 \proofstepwidestar[]{t_{\rm q}\in\mathbb N}{\FormulaRefAuto{B48TagType5}[thm]}
 \proofstepwidestar[]{t_{\rm m}\neq t_{\rm q}}{\FormulaRefAuto{B48TagNe14}[thm]}
 \proofstepwidestar[]{\Succ(t_{\rm m})\neq\Succ(t_{\rm q})}{\FormulaRefAuto{PeanoSuccPreservesNe}[thm]{1,2,3}}
 \proofstepwidestar[]{t_{\rm n}\neq t_{\rm b}}{\FormulaRefAuto{B48CodeCarrier}[def]{4}}
\closeproofpartwideindR

\proofpartwideindR[Trennung der Marken c und q]{t_{\rm c}\neq t_{\rm q}}{t_{\rm c}\neq t_{\rm q}}[B48TagNe34]
 \proofstepwidestar[]{t_{\rm n}\in\mathbb N}{\FormulaRefAuto{B48TagType3}[thm]}
 \proofstepwidestar[]{t_{\rm c}\in\mathbb N}{\FormulaRefAuto{B48TagType4}[thm]}
 \proofstepwidestar[]{t_{\rm n}\neq t_{\rm c}}{\FormulaRefAuto{B48TagNe23}[thm]}
 \proofstepwidestar[]{\Succ(t_{\rm n})\neq\Succ(t_{\rm c})}{\FormulaRefAuto{PeanoSuccPreservesNe}[thm]{1,2,3}}
 \proofstepwidestar[]{t_{\rm c}\neq t_{\rm q}}{\FormulaRefAuto{B48CodeCarrier}[def]{4}}
\closeproofpartwideindR

\proofpartwideindR[Trennung der Marken c und b]{t_{\rm c}\neq t_{\rm b}}{t_{\rm c}\neq t_{\rm b}}[B48TagNe35]
 \proofstepwidestar[]{t_{\rm n}\in\mathbb N}{\FormulaRefAuto{B48TagType3}[thm]}
 \proofstepwidestar[]{t_{\rm q}\in\mathbb N}{\FormulaRefAuto{B48TagType5}[thm]}
 \proofstepwidestar[]{t_{\rm n}\neq t_{\rm q}}{\FormulaRefAuto{B48TagNe24}[thm]}
 \proofstepwidestar[]{\Succ(t_{\rm n})\neq\Succ(t_{\rm q})}{\FormulaRefAuto{PeanoSuccPreservesNe}[thm]{1,2,3}}
 \proofstepwidestar[]{t_{\rm c}\neq t_{\rm b}}{\FormulaRefAuto{B48CodeCarrier}[def]{4}}
\closeproofpartwideindR

\proofpartwideindR[Trennung der Marken q und b]{t_{\rm q}\neq t_{\rm b}}{t_{\rm q}\neq t_{\rm b}}[B48TagNe45]
 \proofstepwidestar[]{t_{\rm c}\in\mathbb N}{\FormulaRefAuto{B48TagType4}[thm]}
 \proofstepwidestar[]{t_{\rm q}\in\mathbb N}{\FormulaRefAuto{B48TagType5}[thm]}
 \proofstepwidestar[]{t_{\rm c}\neq t_{\rm q}}{\FormulaRefAuto{B48TagNe34}[thm]}
 \proofstepwidestar[]{\Succ(t_{\rm c})\neq\Succ(t_{\rm q})}{\FormulaRefAuto{PeanoSuccPreservesNe}[thm]{1,2,3}}
 \proofstepwidestar[]{t_{\rm q}\neq t_{\rm b}}{\FormulaRefAuto{B48CodeCarrier}[def]{4}}
\closeproofpartwideindR

\proofpartwideindR[Typisierung eines Blattzeichens]{k\in\mathbb N,\ i\in\mathbb N,\ j\in\mathbb N\vdash(k,(i,j))\in\Sigma}{k\in\mathbb N,\ i\in\mathbb N,\ j\in\mathbb N\vdash(k,(i,j))\in\Sigma}[B48LabelType]
 \proofstepwidestar[1]{k\in\mathbb N}{\rA}
 \proofstepwidestar[2]{i\in\mathbb N}{\rA}
 \proofstepwidestar[3]{j\in\mathbb N}{\rA}
 \proofstepwidestar[2,3]{(i,j)\in\mathbb N\times\mathbb N}{\FormulaRefAuto{(a,b)\in A\times B\eqvdash a\in A\land b\in B}[thm]{\rAI{2,3}}}
 \proofstepwidestar[1,2,3]{(k,(i,j))\in\mathbb N\times(\mathbb N\times\mathbb N)}{\FormulaRefAuto{(a,b)\in A\times B\eqvdash a\in A\land b\in B}[thm]{\rAI{1,4}}}
 \proofstepwidestar[1,2,3]{(k,(i,j))\in\Sigma}{\FormulaRefAuto{B48CodeCarrier}[def]{5}}
\closeproofpartwideindR

\proofpartwideindR[Blattcodes sind Baumcodes]{c\in\Sigma\vdash\mathsf l(c)\in\mathcal T}{c\in\Sigma\vdash\mathsf l(c)\in\mathcal T}[B48LeafType]
 \proofstepwidestar[1]{c\in\Sigma}{\rA}
 \proofstepwidestar[1]{\LeafTree{\Sigma}{c}\in\BracketTreeSet{\Sigma}}{\FormulaRefAuto{TreeLeafClosure}[thm]{1}}
 \proofstepwidestar[1]{\mathsf l(c)\in\mathcal T}{\FormulaRefAuto{B48CodeCarrier}[def]{2}}
\closeproofpartwideindR

\proofpartwideindR[Knotencodes sind Baumcodes]{p\in\mathcal T,\ q\in\mathcal T\vdash\mathsf n(p,q)\in\mathcal T}{p\in\mathcal T,\ q\in\mathcal T\vdash\mathsf n(p,q)\in\mathcal T}[B48NodeType]
 \proofstepwidestar[1]{p\in\mathcal T}{\rA}
 \proofstepwidestar[2]{q\in\mathcal T}{\rA}
 \proofstepwidestar[1,2]{\NodeTree{\Sigma}{p}{q}\in\BracketTreeSet{\Sigma}}{\FormulaRefAuto{TreeNodeClosure}[thm]{1,2}}
 \proofstepwidestar[1,2]{\mathsf n(p,q)\in\mathcal T}{\FormulaRefAuto{B48CodeCarrier}[def]{3}}
\closeproofpartwideindR

\proofpartwideindR[Baumcode der Negationsmarke]{\mathsf n_0\in\mathcal T}{\mathsf n_0\in\mathcal T}[B48NegRootType]
 \proofstepwidestar[]{t_{\rm n}\in\mathbb N}{\FormulaRefAuto{B48TagType3}[thm]}
 \proofstepwidestar[]{0\in\mathbb N}{\FormulaRefAuto{PeanoZero}[ax]}
 \proofstepwidestar[]{(t_{\rm n},(0,0))\in\Sigma}{\FormulaRefAuto{B48LabelType}[thm]{1,2,2}}
 \proofstepwidestar[]{\mathsf l((t_{\rm n},(0,0)))\in\mathcal T}{\FormulaRefAuto{B48LeafType}[thm]{3}}
 \proofstepwidestar[]{\mathsf n_0\in\mathcal T}{\FormulaRefAuto{B48FormulaConstructors}[def]{4}}
\closeproofpartwideindR

\proofpartwideindR[Baumcode der Konjunktionsmarke]{\mathsf c_0\in\mathcal T}{\mathsf c_0\in\mathcal T}[B48AndRootType]
 \proofstepwidestar[]{t_{\rm c}\in\mathbb N}{\FormulaRefAuto{B48TagType4}[thm]}
 \proofstepwidestar[]{0\in\mathbb N}{\FormulaRefAuto{PeanoZero}[ax]}
 \proofstepwidestar[]{(t_{\rm c},(0,0))\in\Sigma}{\FormulaRefAuto{B48LabelType}[thm]{1,2,2}}
 \proofstepwidestar[]{\mathsf l((t_{\rm c},(0,0)))\in\mathcal T}{\FormulaRefAuto{B48LeafType}[thm]{3}}
 \proofstepwidestar[]{\mathsf c_0\in\mathcal T}{\FormulaRefAuto{B48FormulaConstructors}[def]{4}}
\closeproofpartwideindR

\proofpartwideindR[Baumcode einer Quantorenmarke]{i\in\mathbb N\vdash\mathsf q_i\in\mathcal T}{i\in\mathbb N\vdash\mathsf q_i\in\mathcal T}[B48ExistsRootType]
 \proofstepwidestar[1]{i\in\mathbb N}{\rA}
 \proofstepwidestar[]{t_{\rm q}\in\mathbb N}{\FormulaRefAuto{B48TagType5}[thm]}
 \proofstepwidestar[]{0\in\mathbb N}{\FormulaRefAuto{PeanoZero}[ax]}
 \proofstepwidestar[1]{(t_{\rm q},(i,0))\in\Sigma}{\FormulaRefAuto{B48LabelType}[thm]{2,1,3}}
 \proofstepwidestar[1]{\mathsf l((t_{\rm q},(i,0)))\in\mathcal T}{\FormulaRefAuto{B48LeafType}[thm]{4}}
 \proofstepwidestar[1]{\mathsf q_i\in\mathcal T}{\FormulaRefAuto{B48FormulaConstructors}[def]{5}}
\closeproofpartwideindR

\proofpartwideindR[Typisierung des Negationscodes]{p\in\mathcal T\vdash\mathsf{neg}(p)\in\mathcal T}{p\in\mathcal T\vdash\mathsf{neg}(p)\in\mathcal T}[B48NegCodeType]
 \proofstepwidestar[1]{p\in\mathcal T}{\rA}
 \proofstepwidestar[]{\mathsf n_0\in\mathcal T}{\FormulaRefAuto{B48NegRootType}[thm]}
 \proofstepwidestar[1]{\mathsf n(\mathsf n_0,p)\in\mathcal T}{\FormulaRefAuto{B48NodeType}[thm]{2,1}}
 \proofstepwidestar[1]{\mathsf{neg}(p)\in\mathcal T}{\FormulaRefAuto{B48FormulaConstructors}[def]{3}}
\closeproofpartwideindR

\proofpartwideindR[Typisierung des Konjunktionscodes]{p\in\mathcal T,\ q\in\mathcal T\vdash\mathsf{and}(p,q)\in\mathcal T}{p\in\mathcal T,\ q\in\mathcal T\vdash\mathsf{and}(p,q)\in\mathcal T}[B48AndCodeType]
 \proofstepwidestar[1]{p\in\mathcal T}{\rA}
 \proofstepwidestar[2]{q\in\mathcal T}{\rA}
 \proofstepwidestar[]{\mathsf c_0\in\mathcal T}{\FormulaRefAuto{B48AndRootType}[thm]}
 \proofstepwidestar[1]{\mathsf n(\mathsf c_0,p)\in\mathcal T}{\FormulaRefAuto{B48NodeType}[thm]{3,1}}
 \proofstepwidestar[1,2]{\mathsf n(\mathsf n(\mathsf c_0,p),q)\in\mathcal T}{\FormulaRefAuto{B48NodeType}[thm]{4,2}}
 \proofstepwidestar[1,2]{\mathsf{and}(p,q)\in\mathcal T}{\FormulaRefAuto{B48FormulaConstructors}[def]{5}}
\closeproofpartwideindR

\proofpartwideindR[Typisierung des Existenzcodes]{i\in\mathbb N,\ p\in\mathcal T\vdash\mathsf{ex}_i(p)\in\mathcal T}{i\in\mathbb N,\ p\in\mathcal T\vdash\mathsf{ex}_i(p)\in\mathcal T}[B48ExistsCodeType]
 \proofstepwidestar[1]{i\in\mathbb N}{\rA}
 \proofstepwidestar[2]{p\in\mathcal T}{\rA}
 \proofstepwidestar[1]{\mathsf q_i\in\mathcal T}{\FormulaRefAuto{B48ExistsRootType}[thm]{1}}
 \proofstepwidestar[1,2]{\mathsf n(\mathsf q_i,p)\in\mathcal T}{\FormulaRefAuto{B48NodeType}[thm]{3,2}}
 \proofstepwidestar[1,2]{\mathsf{ex}_i(p)\in\mathcal T}{\FormulaRefAuto{B48FormulaConstructors}[def]{4}}
\closeproofpartwideindR

\par\medskip\noindent\textbf{B. Belegungsmenge und Formelabschluss}\par
\proofpartwideindR[Wahrer Zweig des Fallterms]{P\vdash\BXLVIIIIf{P}{U}{V}=U}{P\vdash\BXLVIIIIf{P}{U}{V}=U}[B48IfTrue]
 \proofstepwidestar[1]{P}{\rA}
 \proofstepwidestar[1]{\IfThenElse{P}{U}{V}=U}{\FormulaRefAuto{P\vdash\IfThenElse{P}{a}{b}=a}[thm]{1}}
 \proofstepwidestar[1]{\BXLVIIIIf{P}{U}{V}=U}{\FormulaRefAuto{B48IfNotation}[def]{2}}
\closeproofpartwideindR

\proofpartwideindR[Falscher Zweig des Fallterms]{\neg P\vdash\BXLVIIIIf{P}{U}{V}=V}{\neg P\vdash\BXLVIIIIf{P}{U}{V}=V}[B48IfFalse]
 \proofstepwidestar[1]{\neg P}{\rA}
 \proofstepwidestar[1]{\IfThenElse{P}{U}{V}=V}{\FormulaRefAuto{\neg P\vdash\IfThenElse{P}{a}{b}=b}[thm]{1}}
 \proofstepwidestar[1]{\BXLVIIIIf{P}{U}{V}=V}{\FormulaRefAuto{B48IfNotation}[def]{2}}
\closeproofpartwideindR

\proofpartwideindR[Typisierung beider Fallzweige]{U\in X,\ V\in X\vdash\BXLVIIIIf{P}{U}{V}\in X}{U\in X,\ V\in X\vdash\BXLVIIIIf{P}{U}{V}\in X}[B48IfTyped]
 \proofstepwidestar[1]{U\in X}{\rA}
 \proofstepwidestar[2]{V\in X}{\rA}
 \proofstepwidestar[]{P\lor\neg P}{\FormulaRefAuto{P\lor\neg P}[thm]}
 \proofstepwidestar[4]{P}{\rA}
 \proofstepwidestar[4]{\BXLVIIIIf{P}{U}{V}=U}{\FormulaRefAuto{B48IfTrue}[thm]{4}}
 \proofstepwidestar[1,4]{\BXLVIIIIf{P}{U}{V}\in X}{\FormulaRefAuto{a=b,\ b\in C\vdash a\in C}[thm]{5,1}}
 \proofstepwidestar[7]{\neg P}{\rA}
 \proofstepwidestar[7]{\BXLVIIIIf{P}{U}{V}=V}{\FormulaRefAuto{B48IfFalse}[thm]{7}}
 \proofstepwidestar[2,7]{\BXLVIIIIf{P}{U}{V}\in X}{\FormulaRefAuto{a=b,\ b\in C\vdash a\in C}[thm]{8,2}}
 \proofstepwidestar[1,2]{\BXLVIIIIf{P}{U}{V}\in X}{\rOE{3,4:6,7:9}}
\closeproofpartwideindR

\proofpartwideindR[Beschraenkte Praedikate liefern Wertemengen]{\{x\in A\mid P(x)\}\in\powerset(A)}{\{x\in A\mid P(x)\}\in\powerset(A)}[B48SeparationPower]
 \proofstepwidestar[]{\{x\in A\mid P(x)\}\subseteq A}{\FormulaRefAuto{\{x\in A\mid P(x)\}\subseteq A}[thm]}
 \proofstepwidestar[]{\{x\in A\mid P(x)\}\in\powerset(A)}{\FormulaRefAuto{x\in\powerset(A)\eqvdash x\subseteq A}[thm]{1}}
\closeproofpartwideindR

\proofpartwideindR[Elementkriterium der Belegungsmenge]{s\in A_D\eqvdash s\colon\mathbb N\to D}{s\in A_D\eqvdash s\colon\mathbb N\to D}[B48AssignmentCriterion]
 \proofstepwidestar[1]{s\in A_D}{\rA}
 \proofstepwidestar[1]{s\in\powerset(\mathbb N\times D)\land s\colon\mathbb N\to D}{\FormulaRefAuto{B48Assignments}[def]{1}}
 \proofstepwidestar[1]{s\colon\mathbb N\to D}{\rAEb{2}}
 \proofstepwidestar[]{s\in A_D\to s\colon\mathbb N\to D}{\rRI{1,3}}
 \proofstepwidestar[5]{s\colon\mathbb N\to D}{\rA}
 \proofstepwidestar[5]{s\subseteq\mathbb N\times D}{\FormulaRefAuto{FunctionGraphSubsetProduct}[thm]{5}}
 \proofstepwidestar[5]{s\in\powerset(\mathbb N\times D)}{\FormulaRefAuto{x\in\powerset(A)\eqvdash x\subseteq A}[thm]{6}}
 \proofstepwidestar[5]{s\in A_D}{\FormulaRefAuto{B48Assignments}[def]{\rAI{7,5}}}
 \proofstepwidestar[]{(s\colon\mathbb N\to D)\to s\in A_D}{\rRI{5,8}}
 \proofstepwidestar[]{s\in A_D\leftrightarrow s\colon\mathbb N\to D}{\rLRI{4,9}}
\closeproofpartwideindR

\proofpartwideindR[Ein Belegungswechsel bleibt eine Belegung]{s\in A_D,\ i\in\mathbb N,\ a\in D\vdash s[i\mapsto a]\in A_D}{s\in A_D,\ i\in\mathbb N,\ a\in D\vdash s[i\mapsto a]\in A_D}[B48AssignmentUpdate]
 \proofstepwidestar[1]{s\in A_D}{\rA}
 \proofstepwidestar[2]{i\in\mathbb N}{\rA}
 \proofstepwidestar[3]{a\in D}{\rA}
 \proofstepwidestar[1]{s\colon\mathbb N\to D}{\FormulaRefAuto{B48AssignmentCriterion}[thm]{1}}
 \proofstepwidestar[5]{j\in\mathbb N}{\rA}
 \proofstepwidestar[1,5]{s(j)\in D}{\FormulaRefAuto{F\colon A\to B,\ x\in A\vdash F(x)\in B}[thm]{4,5}}
 \proofstepwidestar[1,3,5]{\BXLVIIIIf{j=i}{a}{s(j)}\in D}{\FormulaRefAuto{B48IfTyped}[thm]{3,6}}
 \proofstepwidestar[1,3]{\forall j\in\mathbb N\;(\BXLVIIIIf{j=i}{a}{s(j)}\in D)}{\rUI{\rRI{5,7}}}
 \proofstepwidestar[1,3]{\{(j,b)\in\mathbb N\times D\mid b=\BXLVIIIIf{j=i}{a}{s(j)}\}\colon\mathbb N\to D}{\FormulaRefAuto{TypedTermGraphFunction}[thm]{8}}
 \proofstepwidestar[1,3]{s[i\mapsto a]\colon\mathbb N\to D}{\FormulaRefAuto{B48Assignments}[def]{9}}
 \proofstepwidestar[1,3]{s[i\mapsto a]\in A_D}{\FormulaRefAuto{B48AssignmentCriterion}[thm]{10}}
\closeproofpartwideindR

\proofpartwideindR[Abschlussklausel 1]{\mathsf{Cl}(X)\vdash X\subseteq\mathcal T}{\mathsf{Cl}(X)\vdash X\subseteq\mathcal T}[B48ClDomain]
 \proofstepwidestar[1]{\mathsf{Cl}(X)}{\rA}
 \proofstepwidestar[1]{X\subseteq\mathcal T}{\FormulaRefAuto{B48ClDomainAxiom}[ax]{1}}
\closeproofpartwideindR

\proofpartwideindR[Abschlussklausel 2]{\mathsf{Cl}(X)\vdash \forall c\in\Sigma_{\rm at}\;(\mathsf l(c)\in X)}{\mathsf{Cl}(X)\vdash \forall c\in\Sigma_{\rm at}\;(\mathsf l(c)\in X)}[B48ClAtoms]
 \proofstepwidestar[1]{\mathsf{Cl}(X)}{\rA}
 \proofstepwidestar[1]{\forall c\in\Sigma_{\rm at}\;(\mathsf l(c)\in X)}{\FormulaRefAuto{B48ClAtomsAxiom}[ax]{1}}
\closeproofpartwideindR

\proofpartwideindR[Abschlussklausel 3]{\mathsf{Cl}(X)\vdash \forall p\in X\;(\mathsf{neg}(p)\in X)}{\mathsf{Cl}(X)\vdash \forall p\in X\;(\mathsf{neg}(p)\in X)}[B48ClNeg]
 \proofstepwidestar[1]{\mathsf{Cl}(X)}{\rA}
 \proofstepwidestar[1]{\forall p\in X\;(\mathsf{neg}(p)\in X)}{\FormulaRefAuto{B48ClNegAxiom}[ax]{1}}
\closeproofpartwideindR

\proofpartwideindR[Abschlussklausel 4]{\mathsf{Cl}(X)\vdash \forall p\in X\,\forall q\in X\;(\mathsf{and}(p,q)\in X)}{\mathsf{Cl}(X)\vdash \forall p\in X\,\forall q\in X\;(\mathsf{and}(p,q)\in X)}[B48ClAnd]
 \proofstepwidestar[1]{\mathsf{Cl}(X)}{\rA}
 \proofstepwidestar[1]{\forall p\in X\,\forall q\in X\;(\mathsf{and}(p,q)\in X)}{\FormulaRefAuto{B48ClAndAxiom}[ax]{1}}
\closeproofpartwideindR

\proofpartwideindR[Abschlussklausel 5]{\mathsf{Cl}(X)\vdash \forall i\in\mathbb N\,\forall p\in X\;(\mathsf{ex}_i(p)\in X)}{\mathsf{Cl}(X)\vdash \forall i\in\mathbb N\,\forall p\in X\;(\mathsf{ex}_i(p)\in X)}[B48ClExists]
 \proofstepwidestar[1]{\mathsf{Cl}(X)}{\rA}
 \proofstepwidestar[1]{\forall i\in\mathbb N\,\forall p\in X\;(\mathsf{ex}_i(p)\in X)}{\FormulaRefAuto{B48ClExistsAxiom}[ax]{1}}
\closeproofpartwideindR

\proofpartwideindR[Elementkriterium fuer Formelcodes]{p\in\mathcal F\eqvdash p\in\mathcal T\land\forall X(\mathsf{Cl}(X)\to p\in X)}{p\in\mathcal F\eqvdash p\in\mathcal T\land\forall X(\mathsf{Cl}(X)\to p\in X)}[B48FormulaMember]
 \proofstepwidestar[]{p\in\mathcal F\leftrightarrow p\in\mathcal T\land\forall X(\mathsf{Cl}(X)\to p\in X)}{\FormulaRefAuto{B48FormulaSet}[def]}
\closeproofpartwideindR

\proofpartwideindR[Jeder Formelcode ist ein Baumcode]{p\in\mathcal F\vdash p\in\mathcal T}{p\in\mathcal F\vdash p\in\mathcal T}[B48FormulaInTree]
 \proofstepwidestar[1]{p\in\mathcal F}{\rA}
 \proofstepwidestar[1]{p\in\mathcal T\land\forall X(\mathsf{Cl}(X)\to p\in X)}{\FormulaRefAuto{B48FormulaMember}[thm]{1}}
 \proofstepwidestar[1]{p\in\mathcal T}{\rAEa{2}}
\closeproofpartwideindR

\proofpartwideindR[Minimalitaet der Formelcodes]{\mathsf{Cl}(X),\ p\in\mathcal F\vdash p\in X}{\mathsf{Cl}(X),\ p\in\mathcal F\vdash p\in X}[B48FormulaMinimalMember]
 \proofstepwidestar[1]{\mathsf{Cl}(X)}{\rA}
 \proofstepwidestar[2]{p\in\mathcal F}{\rA}
 \proofstepwidestar[2]{p\in\mathcal T\land\forall Y(\mathsf{Cl}(Y)\to p\in Y)}{\FormulaRefAuto{B48FormulaMember}[thm]{2}}
 \proofstepwidestar[2]{\forall Y(\mathsf{Cl}(Y)\to p\in Y)}{\rAEb{3}}
 \proofstepwidestar[2]{\mathsf{Cl}(X)\to p\in X}{\rUE{4}}
 \proofstepwidestar[1,2]{p\in X}{\rRE{5,1}}
\closeproofpartwideindR

\proofpartwideindR[Atomare Zeichen liegen im Alphabet]{c\in\Sigma_{\rm at}\vdash c\in\Sigma}{c\in\Sigma_{\rm at}\vdash c\in\Sigma}[B48AtomicLabelType]
 \proofstepwidestar[1]{c\in\Sigma_{\rm at}}{\rA}
 \proofstepwidestar[]{\Sigma_{\rm at}\subseteq\Sigma}{\FormulaRefAuto{\{x\in A\mid P(x)\}\subseteq A}[thm]}
 \proofstepwidestar[1]{c\in\Sigma}{\FormulaRefAuto{A\subseteq B,\ x\in A\vdash x\in B}[thm]{2,1}}
\closeproofpartwideindR

\proofpartwideindR[Atomare Formelcodes]{c\in\Sigma_{\rm at}\vdash\mathsf l(c)\in\mathcal F}{c\in\Sigma_{\rm at}\vdash\mathsf l(c)\in\mathcal F}[B48FormulaAtom]
 \proofstepwidestar[1]{c\in\Sigma_{\rm at}}{\rA}
 \proofstepwidestar[1]{c\in\Sigma}{\FormulaRefAuto{B48AtomicLabelType}[thm]{1}}
 \proofstepwidestar[1]{\mathsf l(c)\in\mathcal T}{\FormulaRefAuto{B48LeafType}[thm]{2}}
 \proofstepwidestar[4]{\mathsf{Cl}(X)}{\rA}
 \proofstepwidestar[4]{\forall d\in\Sigma_{\rm at}\;(\mathsf l(d)\in X)}{\FormulaRefAuto{B48ClAtoms}[thm]{4}}
 \proofstepwidestar[1,4]{\mathsf l(c)\in X}{\rRE{\rUE{5},1}}
 \proofstepwidestar[1]{\forall X(\mathsf{Cl}(X)\to\mathsf l(c)\in X)}{\rUI{\rRI{4,6}}}
 \proofstepwidestar[1]{\mathsf l(c)\in\mathcal F}{\FormulaRefAuto{B48FormulaMember}[thm]{\rAI{3,7}}}
\closeproofpartwideindR

\proofpartwideindR[Abschluss unter Negationscodes]{p\in\mathcal F\vdash\mathsf{neg}(p)\in\mathcal F}{p\in\mathcal F\vdash\mathsf{neg}(p)\in\mathcal F}[B48FormulaNeg]
 \proofstepwidestar[1]{p\in\mathcal F}{\rA}
 \proofstepwidestar[1]{p\in\mathcal T}{\FormulaRefAuto{B48FormulaInTree}[thm]{1}}
 \proofstepwidestar[1]{\mathsf{neg}(p)\in\mathcal T}{\FormulaRefAuto{B48NegCodeType}[thm]{2}}
 \proofstepwidestar[4]{\mathsf{Cl}(X)}{\rA}
 \proofstepwidestar[1,4]{p\in X}{\FormulaRefAuto{B48FormulaMinimalMember}[thm]{4,1}}
 \proofstepwidestar[4]{\forall q\in X\;(\mathsf{neg}(q)\in X)}{\FormulaRefAuto{B48ClNeg}[thm]{4}}
 \proofstepwidestar[1,4]{\mathsf{neg}(p)\in X}{\rRE{\rUE{6},5}}
 \proofstepwidestar[1]{\forall X(\mathsf{Cl}(X)\to\mathsf{neg}(p)\in X)}{\rUI{\rRI{4,7}}}
 \proofstepwidestar[1]{\mathsf{neg}(p)\in\mathcal F}{\FormulaRefAuto{B48FormulaMember}[thm]{\rAI{3,8}}}
\closeproofpartwideindR

\proofpartwideindR[Abschluss unter Konjunktionscodes]{p\in\mathcal F,\ q\in\mathcal F\vdash\mathsf{and}(p,q)\in\mathcal F}{p\in\mathcal F,\ q\in\mathcal F\vdash\mathsf{and}(p,q)\in\mathcal F}[B48FormulaAnd]
 \proofstepwidestar[1]{p\in\mathcal F}{\rA}
 \proofstepwidestar[2]{q\in\mathcal F}{\rA}
 \proofstepwidestar[1]{p\in\mathcal T}{\FormulaRefAuto{B48FormulaInTree}[thm]{1}}
 \proofstepwidestar[2]{q\in\mathcal T}{\FormulaRefAuto{B48FormulaInTree}[thm]{2}}
 \proofstepwidestar[1,2]{\mathsf{and}(p,q)\in\mathcal T}{\FormulaRefAuto{B48AndCodeType}[thm]{3,4}}
 \proofstepwidestar[6]{\mathsf{Cl}(X)}{\rA}
 \proofstepwidestar[1,6]{p\in X}{\FormulaRefAuto{B48FormulaMinimalMember}[thm]{6,1}}
 \proofstepwidestar[2,6]{q\in X}{\FormulaRefAuto{B48FormulaMinimalMember}[thm]{6,2}}
 \proofstepwidestar[6]{\forall u\in X\,\forall v\in X\;(\mathsf{and}(u,v)\in X)}{\FormulaRefAuto{B48ClAnd}[thm]{6}}
 \proofstepwidestar[1,2,6]{\mathsf{and}(p,q)\in X}{\rRE{\rUE{\rRE{\rUE{9},7}},8}}
 \proofstepwidestar[1,2]{\forall X(\mathsf{Cl}(X)\to\mathsf{and}(p,q)\in X)}{\rUI{\rRI{6,10}}}
 \proofstepwidestar[1,2]{\mathsf{and}(p,q)\in\mathcal F}{\FormulaRefAuto{B48FormulaMember}[thm]{\rAI{5,11}}}
\closeproofpartwideindR

\proofpartwideindR[Abschluss unter Existenzcodes]{i\in\mathbb N,\ p\in\mathcal F\vdash\mathsf{ex}_i(p)\in\mathcal F}{i\in\mathbb N,\ p\in\mathcal F\vdash\mathsf{ex}_i(p)\in\mathcal F}[B48FormulaExists]
 \proofstepwidestar[1]{i\in\mathbb N}{\rA}
 \proofstepwidestar[2]{p\in\mathcal F}{\rA}
 \proofstepwidestar[2]{p\in\mathcal T}{\FormulaRefAuto{B48FormulaInTree}[thm]{2}}
 \proofstepwidestar[1,2]{\mathsf{ex}_i(p)\in\mathcal T}{\FormulaRefAuto{B48ExistsCodeType}[thm]{1,3}}
 \proofstepwidestar[5]{\mathsf{Cl}(X)}{\rA}
 \proofstepwidestar[2,5]{p\in X}{\FormulaRefAuto{B48FormulaMinimalMember}[thm]{5,2}}
 \proofstepwidestar[5]{\forall j\in\mathbb N\,\forall q\in X\;(\mathsf{ex}_j(q)\in X)}{\FormulaRefAuto{B48ClExists}[thm]{5}}
 \proofstepwidestar[1,2,5]{\mathsf{ex}_i(p)\in X}{\rRE{\rUE{\rRE{\rUE{7},1}},6}}
 \proofstepwidestar[1,2]{\forall X(\mathsf{Cl}(X)\to\mathsf{ex}_i(p)\in X)}{\rUI{\rRI{5,8}}}
 \proofstepwidestar[1,2]{\mathsf{ex}_i(p)\in\mathcal F}{\FormulaRefAuto{B48FormulaMember}[thm]{\rAI{4,9}}}
\closeproofpartwideindR

\proofpartwideindR[Induktion auf der definierten Formelmenge]{\mathsf{Cl}(\{p\in\mathcal F\mid P(p)\})\vdash\forall p\in\mathcal F\;P(p)}{\mathsf{Cl}(\{p\in\mathcal F\mid P(p)\})\vdash\forall p\in\mathcal F\;P(p)}[B48FormulaInduction]
 \proofstepwidestar[1]{\mathsf{Cl}(\{p\in\mathcal F\mid P(p)\})}{\rA}
 \proofstepwidestar[2]{q\in\mathcal F}{\rA}
 \proofstepwidestar[1,2]{q\in\{p\in\mathcal F\mid P(p)\}}{\FormulaRefAuto{B48FormulaMinimalMember}[thm]{1,2}}
 \proofstepwidestar[1,2]{q\in\mathcal F\land P(q)}{\FormulaRefAuto{x\in\{u\in A\mid P(u)\}\eqvdash x\in A\land P(x)}[thm]{3}}
 \proofstepwidestar[1,2]{P(q)}{\rAEb{4}}
 \proofstepwidestar[1]{\forall q\in\mathcal F\;P(q)}{\rUI{\rRI{2,5}}}
\closeproofpartwideindR

\par\medskip\noindent\textbf{C. Unterscheidung der Auswertungsfaelle}\par
\proofpartwideindR[Baumcodes sind nichtleere Woerter]{p\in\mathcal T\vdash p\in\NonemptyWords{\TreeAlphabet{\Sigma}}}{p\in\mathcal T\vdash p\in\NonemptyWords{\TreeAlphabet{\Sigma}}}[B48TreeWord]
 \proofstepwidestar[1]{p\in\mathcal T}{\rA}
 \proofstepwidestar[1]{p\in\NonemptyWords{\TreeAlphabet{\Sigma}}\land\exists n\in\mathbb N\;(p\in\TreeStage{\Sigma}{n})}{\FormulaRefAuto{FullPlanarBinaryTreeSetDef}[def]{1}}
 \proofstepwidestar[1]{p\in\NonemptyWords{\TreeAlphabet{\Sigma}}}{\rAEa{2}}
\closeproofpartwideindR

\proofpartwideindR[Injektivitaet der Blattcodes]{c\in\Sigma,\ d\in\Sigma,\ \mathsf l(c)=\mathsf l(d)\vdash c=d}{c\in\Sigma,\ d\in\Sigma,\ \mathsf l(c)=\mathsf l(d)\vdash c=d}[B48LeafInjective]
 \proofstepwidestar[1]{c\in\Sigma}{\rA}
 \proofstepwidestar[2]{d\in\Sigma}{\rA}
 \proofstepwidestar[3]{\mathsf l(c)=\mathsf l(d)}{\rA}
 \proofstepwidestar[1,2]{\mathsf l(c)=\mathsf l(d)\leftrightarrow c=d}{\FormulaRefAuto{TreeLeafConstructorInjective}[thm]{1,2}}
 \proofstepwidestar[1,2,3]{c=d}{\rLREa{4,3}}
\closeproofpartwideindR

\proofpartwideindR[Verschiedene Marken liefern verschiedene Blattcodes]{(k,a)\in\Sigma,\ (l,b)\in\Sigma,\ k\neq l\vdash\mathsf l((k,a))\neq\mathsf l((l,b))}{(k,a)\in\Sigma,\ (l,b)\in\Sigma,\ k\neq l\vdash\mathsf l((k,a))\neq\mathsf l((l,b))}[B48TaggedLeafNe]
 \proofstepwidestar[1]{(k,a)\in\Sigma}{\rA}
 \proofstepwidestar[2]{(l,b)\in\Sigma}{\rA}
 \proofstepwidestar[3]{k\neq l}{\rA}
 \proofstepwidestar[4]{\mathsf l((k,a))=\mathsf l((l,b))}{\rA}
 \proofstepwidestar[1,2,4]{(k,a)=(l,b)}{\FormulaRefAuto{B48LeafInjective}[thm]{1,2,4}}
 \proofstepwidestar[1,2,4]{k=l\land a=b}{\FormulaRefAuto{(a,b)=(c,d)\eqvdash a=c\land b=d}[thm]{5}}
 \proofstepwidestar[1,2,4]{k=l}{\rAEa{6}}
 \proofstepwidestar[1,2,3,4]{\bot}{\rBI{7,3}}
 \proofstepwidestar[1,2,3]{\mathsf l((k,a))\neq\mathsf l((l,b))}{\rCI{4,8}}
\closeproofpartwideindR

\proofpartwideindR[Blatt und Knoten sind verschieden]{c\in\Sigma,\ p\in\mathcal T,\ q\in\mathcal T\vdash\mathsf l(c)\neq\mathsf n(p,q)}{c\in\Sigma,\ p\in\mathcal T,\ q\in\mathcal T\vdash\mathsf l(c)\neq\mathsf n(p,q)}[B48LeafNotNode]
 \proofstepwidestar[1]{c\in\Sigma}{\rA}
 \proofstepwidestar[2]{p\in\mathcal T}{\rA}
 \proofstepwidestar[3]{q\in\mathcal T}{\rA}
 \proofstepwidestar[2]{p\in\NonemptyWords{\TreeAlphabet{\Sigma}}}{\FormulaRefAuto{B48TreeWord}[thm]{2}}
 \proofstepwidestar[3]{q\in\NonemptyWords{\TreeAlphabet{\Sigma}}}{\FormulaRefAuto{B48TreeWord}[thm]{3}}
 \proofstepwidestar[1,2,3]{\mathsf l(c)\neq\mathsf n(p,q)}{\FormulaRefAuto{TreeConstructorsDisjoint}[thm]{1,4,5}}
\closeproofpartwideindR

\proofpartwideindR[Trennung der Konjunktionsmarke]{\mathsf c_0\neq\mathsf n_0}{\mathsf c_0\neq\mathsf n_0}[B48AndNeNeg]
 \proofstepwidestar[]{0\in\mathbb N}{\FormulaRefAuto{PeanoZero}[ax]}
 \proofstepwidestar[]{t_{\rm c}\in\mathbb N}{\FormulaRefAuto{B48TagType4}[thm]}
 \proofstepwidestar[]{t_{\rm n}\in\mathbb N}{\FormulaRefAuto{B48TagType3}[thm]}
 \proofstepwidestar[]{(t_{\rm c},(0,0))\in\Sigma}{\FormulaRefAuto{B48LabelType}[thm]{2,1,1}}
 \proofstepwidestar[]{(t_{\rm n},(0,0))\in\Sigma}{\FormulaRefAuto{B48LabelType}[thm]{3,1,1}}
 \proofstepwidestar[]{t_{\rm n}\neq t_{\rm c}}{\FormulaRefAuto{B48TagNe23}[thm]}
 \proofstepwidestar[]{t_{\rm c}\neq t_{\rm n}}{\FormulaRefAuto{a\neq b\vdash b\neq a}[thm]{6}}
 \proofstepwidestar[]{\mathsf c_0\neq\mathsf n_0}{\FormulaRefAuto{B48TaggedLeafNe}[thm]{4,5,7}}
\closeproofpartwideindR

\proofpartwideindR[Trennung der Quantorenmarke]{i\in\mathbb N\vdash \mathsf q_i\neq\mathsf n_0}{i\in\mathbb N\vdash \mathsf q_i\neq\mathsf n_0}[B48QNeNeg]
 \proofstepwidestar[1]{i\in\mathbb N}{\rA}
 \proofstepwidestar[]{0\in\mathbb N}{\FormulaRefAuto{PeanoZero}[ax]}
 \proofstepwidestar[]{t_{\rm q}\in\mathbb N}{\FormulaRefAuto{B48TagType5}[thm]}
 \proofstepwidestar[]{t_{\rm n}\in\mathbb N}{\FormulaRefAuto{B48TagType3}[thm]}
 \proofstepwidestar[1]{(t_{\rm q},(i,0))\in\Sigma}{\FormulaRefAuto{B48LabelType}[thm]{3,1,2}}
 \proofstepwidestar[]{(t_{\rm n},(0,0))\in\Sigma}{\FormulaRefAuto{B48LabelType}[thm]{4,2,2}}
 \proofstepwidestar[]{t_{\rm n}\neq t_{\rm q}}{\FormulaRefAuto{B48TagNe24}[thm]}
 \proofstepwidestar[]{t_{\rm q}\neq t_{\rm n}}{\FormulaRefAuto{a\neq b\vdash b\neq a}[thm]{7}}
 \proofstepwidestar[1]{\mathsf q_i\neq\mathsf n_0}{\FormulaRefAuto{B48TaggedLeafNe}[thm]{5,6,8}}
\closeproofpartwideindR

\proofpartwideindR[Trennung der Quantorenmarke]{i\in\mathbb N\vdash \mathsf q_i\neq\mathsf c_0}{i\in\mathbb N\vdash \mathsf q_i\neq\mathsf c_0}[B48QNeAnd]
 \proofstepwidestar[1]{i\in\mathbb N}{\rA}
 \proofstepwidestar[]{0\in\mathbb N}{\FormulaRefAuto{PeanoZero}[ax]}
 \proofstepwidestar[]{t_{\rm q}\in\mathbb N}{\FormulaRefAuto{B48TagType5}[thm]}
 \proofstepwidestar[]{t_{\rm c}\in\mathbb N}{\FormulaRefAuto{B48TagType4}[thm]}
 \proofstepwidestar[1]{(t_{\rm q},(i,0))\in\Sigma}{\FormulaRefAuto{B48LabelType}[thm]{3,1,2}}
 \proofstepwidestar[]{(t_{\rm c},(0,0))\in\Sigma}{\FormulaRefAuto{B48LabelType}[thm]{4,2,2}}
 \proofstepwidestar[]{t_{\rm c}\neq t_{\rm q}}{\FormulaRefAuto{B48TagNe34}[thm]}
 \proofstepwidestar[]{t_{\rm q}\neq t_{\rm c}}{\FormulaRefAuto{a\neq b\vdash b\neq a}[thm]{7}}
 \proofstepwidestar[1]{\mathsf q_i\neq\mathsf c_0}{\FormulaRefAuto{B48TaggedLeafNe}[thm]{5,6,8}}
\closeproofpartwideindR

\proofpartwideindR[Quantorenmarken bestimmen den Variablenindex]{i\in\mathbb N,\ j\in\mathbb N,\ \mathsf q_i=\mathsf q_j\vdash i=j}{i\in\mathbb N,\ j\in\mathbb N,\ \mathsf q_i=\mathsf q_j\vdash i=j}[B48QuantifierRootInjective]
 \proofstepwidestar[1]{i\in\mathbb N}{\rA}
 \proofstepwidestar[2]{j\in\mathbb N}{\rA}
 \proofstepwidestar[3]{\mathsf q_i=\mathsf q_j}{\rA}
 \proofstepwidestar[]{t_{\rm q}\in\mathbb N}{\FormulaRefAuto{B48TagType5}[thm]}
 \proofstepwidestar[]{0\in\mathbb N}{\FormulaRefAuto{PeanoZero}[ax]}
 \proofstepwidestar[1]{(t_{\rm q},(i,0))\in\Sigma}{\FormulaRefAuto{B48LabelType}[thm]{4,1,5}}
 \proofstepwidestar[2]{(t_{\rm q},(j,0))\in\Sigma}{\FormulaRefAuto{B48LabelType}[thm]{4,2,5}}
 \proofstepwidestar[1,2,3]{(t_{\rm q},(i,0))=(t_{\rm q},(j,0))}{\FormulaRefAuto{B48LeafInjective}[thm]{6,7,3}}
 \proofstepwidestar[1,2,3]{i=j\land0=0}{\FormulaRefAuto{TaggedNestedPairInjective}[thm]{8}}
 \proofstepwidestar[1,2,3]{i=j}{\rAEa{9}}
\closeproofpartwideindR

\proofpartwideindR[Ein Blatt ist kein Konjunktions-Zwischencode]{c\in\Sigma\vdash\neg J(\mathsf l(c))}{c\in\Sigma\vdash\neg J(\mathsf l(c))}[B48LeafNotAndSlot]
 \proofstepwidestar[1]{c\in\Sigma}{\rA}
 \proofstepwidestar[2]{J(\mathsf l(c))}{\rA}
 \proofstepwidestar[2]{\exists p\in\mathcal T\;(\mathsf l(c)=\mathsf n(\mathsf c_0,p))}{\FormulaRefAuto{B48EvaluationOperators}[def]{2}}
 \proofstepwidestar[4]{p\in\mathcal T\land\mathsf l(c)=\mathsf n(\mathsf c_0,p)}{\rA}
 \proofstepwidestar[4]{p\in\mathcal T}{\rAEa{4}}
 \proofstepwidestar[]{\mathsf c_0\in\mathcal T}{\FormulaRefAuto{B48AndRootType}[thm]}
 \proofstepwidestar[1,4]{\mathsf l(c)\neq\mathsf n(\mathsf c_0,p)}{\FormulaRefAuto{B48LeafNotNode}[thm]{1,6,5}}
 \proofstepwidestar[4]{\mathsf l(c)=\mathsf n(\mathsf c_0,p)}{\rAEb{4}}
 \proofstepwidestar[1,4]{\bot}{\rBI{8,7}}
 \proofstepwidestar[1,2]{\bot}{\rEE{3,4:9}}
 \proofstepwidestar[1]{\neg J(\mathsf l(c))}{\rCI{2,10}}
\closeproofpartwideindR

\proofpartwideindR[Konjunktions-Zwischencodes werden erkannt]{p\in\mathcal T\vdash J(\mathsf n(\mathsf c_0,p))}{p\in\mathcal T\vdash J(\mathsf n(\mathsf c_0,p))}[B48AndSlotSelector]
 \proofstepwidestar[1]{p\in\mathcal T}{\rA}
 \proofstepwidestar[]{\mathsf n(\mathsf c_0,p)=\mathsf n(\mathsf c_0,p)}{\rII}
 \proofstepwidestar[1]{\exists q\in\mathcal T\;(\mathsf n(\mathsf c_0,p)=\mathsf n(\mathsf c_0,q))}{\rEI{\rAI{1,2}}}
 \proofstepwidestar[1]{J(\mathsf n(\mathsf c_0,p))}{\FormulaRefAuto{B48EvaluationOperators}[def]{3}}
\closeproofpartwideindR

\proofpartwideindR[Ein Konjunktions-Zwischencode ist kein Blatt]{p\in\mathcal T\vdash\mathsf n(\mathsf c_0,p)\neq\mathsf n_0}{p\in\mathcal T\vdash\mathsf n(\mathsf c_0,p)\neq\mathsf n_0}[B48AndSlotNeNeg]
 \proofstepwidestar[1]{p\in\mathcal T}{\rA}
 \proofstepwidestar[]{0\in\mathbb N}{\FormulaRefAuto{PeanoZero}[ax]}
 \proofstepwidestar[]{t_{\rm n}\in\mathbb N}{\FormulaRefAuto{B48TagType3}[thm]}
 \proofstepwidestar[]{(t_{\rm n},(0,0))\in\Sigma}{\FormulaRefAuto{B48LabelType}[thm]{3,2,2}}
 \proofstepwidestar[]{\mathsf c_0\in\mathcal T}{\FormulaRefAuto{B48AndRootType}[thm]}
 \proofstepwidestar[1]{\mathsf n_0\neq\mathsf n(\mathsf c_0,p)}{\FormulaRefAuto{B48LeafNotNode}[thm]{4,5,1}}
 \proofstepwidestar[1]{\mathsf n(\mathsf c_0,p)\neq\mathsf n_0}{\FormulaRefAuto{a\neq b\vdash b\neq a}[thm]{6}}
\closeproofpartwideindR

\proofpartwideindR[Ein Konjunktions-Zwischencode ist kein Blatt]{p\in\mathcal T\vdash\mathsf n(\mathsf c_0,p)\neq\mathsf c_0}{p\in\mathcal T\vdash\mathsf n(\mathsf c_0,p)\neq\mathsf c_0}[B48AndSlotNeAnd]
 \proofstepwidestar[1]{p\in\mathcal T}{\rA}
 \proofstepwidestar[]{0\in\mathbb N}{\FormulaRefAuto{PeanoZero}[ax]}
 \proofstepwidestar[]{t_{\rm c}\in\mathbb N}{\FormulaRefAuto{B48TagType4}[thm]}
 \proofstepwidestar[]{(t_{\rm c},(0,0))\in\Sigma}{\FormulaRefAuto{B48LabelType}[thm]{3,2,2}}
 \proofstepwidestar[]{\mathsf c_0\in\mathcal T}{\FormulaRefAuto{B48AndRootType}[thm]}
 \proofstepwidestar[1]{\mathsf c_0\neq\mathsf n(\mathsf c_0,p)}{\FormulaRefAuto{B48LeafNotNode}[thm]{4,5,1}}
 \proofstepwidestar[1]{\mathsf n(\mathsf c_0,p)\neq\mathsf c_0}{\FormulaRefAuto{a\neq b\vdash b\neq a}[thm]{6}}
\closeproofpartwideindR

\proofpartwideindR[Einzeln referenzierter Auswertungszweig]{t=\mathsf n_0\vdash H_M(t,U,V)=A_D\setminus V}{t=\mathsf n_0\vdash H_M(t,U,V)=A_D\setminus V}[B48HNegSwitch]
 \proofstepwidestar[1]{t=\mathsf n_0}{\rA}
 \proofstepwidestar[]{H_M(t,U,V)=\BXLVIIIIf{t=\mathsf n_0}{A_D\setminus V}{H_{1,M}(t,U,V)}}{\FormulaRefAuto{B48EvaluationOperators}[def]}
 \proofstepwidestar[1]{\BXLVIIIIf{t=\mathsf n_0}{A_D\setminus V}{H_{1,M}(t,U,V)}=A_D\setminus V}{\FormulaRefAuto{B48IfTrue}[thm]{1}}
 \proofstepwidestar[1]{H_M(t,U,V)=A_D\setminus V}{\rIE{3,2}}
\closeproofpartwideindR

\proofpartwideindR[Einzeln referenzierter Auswertungszweig]{\neg(t=\mathsf n_0)\vdash H_M(t,U,V)=H_{1,M}(t,U,V)}{\neg(t=\mathsf n_0)\vdash H_M(t,U,V)=H_{1,M}(t,U,V)}[B48HSkipNeg]
 \proofstepwidestar[1]{\neg(t=\mathsf n_0)}{\rA}
 \proofstepwidestar[]{H_M(t,U,V)=\BXLVIIIIf{t=\mathsf n_0}{A_D\setminus V}{H_{1,M}(t,U,V)}}{\FormulaRefAuto{B48EvaluationOperators}[def]}
 \proofstepwidestar[1]{\BXLVIIIIf{t=\mathsf n_0}{A_D\setminus V}{H_{1,M}(t,U,V)}=H_{1,M}(t,U,V)}{\FormulaRefAuto{B48IfFalse}[thm]{1}}
 \proofstepwidestar[1]{H_M(t,U,V)=H_{1,M}(t,U,V)}{\rIE{3,2}}
\closeproofpartwideindR

\proofpartwideindR[Einzeln referenzierter Auswertungszweig]{t=\mathsf c_0\vdash H_{1,M}(t,U,V)=V}{t=\mathsf c_0\vdash H_{1,M}(t,U,V)=V}[B48HAndSwitch]
 \proofstepwidestar[1]{t=\mathsf c_0}{\rA}
 \proofstepwidestar[]{H_{1,M}(t,U,V)=\BXLVIIIIf{t=\mathsf c_0}{V}{H_{2,M}(t,U,V)}}{\FormulaRefAuto{B48EvaluationOperators}[def]}
 \proofstepwidestar[1]{\BXLVIIIIf{t=\mathsf c_0}{V}{H_{2,M}(t,U,V)}=V}{\FormulaRefAuto{B48IfTrue}[thm]{1}}
 \proofstepwidestar[1]{H_{1,M}(t,U,V)=V}{\rIE{3,2}}
\closeproofpartwideindR

\proofpartwideindR[Einzeln referenzierter Auswertungszweig]{\neg(t=\mathsf c_0)\vdash H_{1,M}(t,U,V)=H_{2,M}(t,U,V)}{\neg(t=\mathsf c_0)\vdash H_{1,M}(t,U,V)=H_{2,M}(t,U,V)}[B48HSkipAnd]
 \proofstepwidestar[1]{\neg(t=\mathsf c_0)}{\rA}
 \proofstepwidestar[]{H_{1,M}(t,U,V)=\BXLVIIIIf{t=\mathsf c_0}{V}{H_{2,M}(t,U,V)}}{\FormulaRefAuto{B48EvaluationOperators}[def]}
 \proofstepwidestar[1]{\BXLVIIIIf{t=\mathsf c_0}{V}{H_{2,M}(t,U,V)}=H_{2,M}(t,U,V)}{\FormulaRefAuto{B48IfFalse}[thm]{1}}
 \proofstepwidestar[1]{H_{1,M}(t,U,V)=H_{2,M}(t,U,V)}{\rIE{3,2}}
\closeproofpartwideindR

\proofpartwideindR[Einzeln referenzierter Auswertungszweig]{J(t)\vdash H_{2,M}(t,U,V)=U\cap V}{J(t)\vdash H_{2,M}(t,U,V)=U\cap V}[B48HSlotSwitch]
 \proofstepwidestar[1]{J(t)}{\rA}
 \proofstepwidestar[]{H_{2,M}(t,U,V)=\BXLVIIIIf{J(t)}{U\cap V}{Q(t,V)}}{\FormulaRefAuto{B48EvaluationOperators}[def]}
 \proofstepwidestar[1]{\BXLVIIIIf{J(t)}{U\cap V}{Q(t,V)}=U\cap V}{\FormulaRefAuto{B48IfTrue}[thm]{1}}
 \proofstepwidestar[1]{H_{2,M}(t,U,V)=U\cap V}{\rIE{3,2}}
\closeproofpartwideindR

\proofpartwideindR[Einzeln referenzierter Auswertungszweig]{\neg(J(t))\vdash H_{2,M}(t,U,V)=Q(t,V)}{\neg(J(t))\vdash H_{2,M}(t,U,V)=Q(t,V)}[B48HSkipSlot]
 \proofstepwidestar[1]{\neg(J(t))}{\rA}
 \proofstepwidestar[]{H_{2,M}(t,U,V)=\BXLVIIIIf{J(t)}{U\cap V}{Q(t,V)}}{\FormulaRefAuto{B48EvaluationOperators}[def]}
 \proofstepwidestar[1]{\BXLVIIIIf{J(t)}{U\cap V}{Q(t,V)}=Q(t,V)}{\FormulaRefAuto{B48IfFalse}[thm]{1}}
 \proofstepwidestar[1]{H_{2,M}(t,U,V)=Q(t,V)}{\rIE{3,2}}
\closeproofpartwideindR

\proofpartwideindR[Auswertung der Negationsmarke]{H_M(\mathsf n_0,U,V)=A_D\setminus V}{H_M(\mathsf n_0,U,V)=A_D\setminus V}[B48HNeg]
 \proofstepwidestar[]{\mathsf n_0=\mathsf n_0}{\rII}
 \proofstepwidestar[]{H_M(\mathsf n_0,U,V)=A_D\setminus V}{\FormulaRefAuto{B48HNegSwitch}[thm]{1}}
\closeproofpartwideindR

\proofpartwideindR[Auswertung des ersten Konjunktionsarguments]{H_M(\mathsf c_0,U,V)=V}{H_M(\mathsf c_0,U,V)=V}[B48HAndPrefix]
 \proofstepwidestar[]{\mathsf c_0\neq\mathsf n_0}{\FormulaRefAuto{B48AndNeNeg}[thm]}
 \proofstepwidestar[]{H_M(\mathsf c_0,U,V)=H_{1,M}(\mathsf c_0,U,V)}{\FormulaRefAuto{B48HSkipNeg}[thm]{1}}
 \proofstepwidestar[]{\mathsf c_0=\mathsf c_0}{\rII}
 \proofstepwidestar[]{H_{1,M}(\mathsf c_0,U,V)=V}{\FormulaRefAuto{B48HAndSwitch}[thm]{3}}
 \proofstepwidestar[]{H_M(\mathsf c_0,U,V)=V}{\rIE{4,2}}
\closeproofpartwideindR

\proofpartwideindR[Auswertung des zweiten Konjunktionsarguments]{p\in\mathcal T\vdash H_M(\mathsf n(\mathsf c_0,p),U,V)=U\cap V}{p\in\mathcal T\vdash H_M(\mathsf n(\mathsf c_0,p),U,V)=U\cap V}[B48HAnd]
 \proofstepwidestar[1]{p\in\mathcal T}{\rA}
 \proofstepwidestar[1]{\mathsf n(\mathsf c_0,p)\neq\mathsf n_0}{\FormulaRefAuto{B48AndSlotNeNeg}[thm]{1}}
 \proofstepwidestar[1]{\mathsf n(\mathsf c_0,p)\neq\mathsf c_0}{\FormulaRefAuto{B48AndSlotNeAnd}[thm]{1}}
 \proofstepwidestar[1]{J(\mathsf n(\mathsf c_0,p))}{\FormulaRefAuto{B48AndSlotSelector}[thm]{1}}
 \proofstepwidestar[1]{H_M(\mathsf n(\mathsf c_0,p),U,V)=H_{1,M}(\mathsf n(\mathsf c_0,p),U,V)}{\FormulaRefAuto{B48HSkipNeg}[thm]{2}}
 \proofstepwidestar[1]{H_{1,M}(\mathsf n(\mathsf c_0,p),U,V)=H_{2,M}(\mathsf n(\mathsf c_0,p),U,V)}{\FormulaRefAuto{B48HSkipAnd}[thm]{3}}
 \proofstepwidestar[1]{H_{2,M}(\mathsf n(\mathsf c_0,p),U,V)=U\cap V}{\FormulaRefAuto{B48HSlotSwitch}[thm]{4}}
 \proofstepwidestar[1]{H_M(\mathsf n(\mathsf c_0,p),U,V)=U\cap V}{\rIE{7,\rIE{6,5}}}
\closeproofpartwideindR

\proofpartwideindR[Elementkriterium bei bekanntem Traeger]{s\in A\vdash\bigl(s\in\{x\in A\mid P(x)\}\leftrightarrow P(s)\bigr)}{s\in A\vdash\bigl(s\in\{x\in A\mid P(x)\}\leftrightarrow P(s)\bigr)}[B48SeparationWithin]
 \proofstepwidestar[1]{s\in A}{\rA}
 \proofstepwidestar[2]{s\in\{x\in A\mid P(x)\}}{\rA}
 \proofstepwidestar[2]{s\in A\land P(s)}{\FormulaRefAuto{x\in\{u\in A\mid P(u)\}\eqvdash x\in A\land P(x)}[thm]{2}}
 \proofstepwidestar[2]{P(s)}{\rAEb{3}}
 \proofstepwidestar[]{s\in\{x\in A\mid P(x)\}\to P(s)}{\rRI{2,4}}
 \proofstepwidestar[6]{P(s)}{\rA}
 \proofstepwidestar[1,6]{s\in A\land P(s)}{\rAI{1,6}}
 \proofstepwidestar[1,6]{s\in\{x\in A\mid P(x)\}}{\FormulaRefAuto{x\in\{u\in A\mid P(u)\}\eqvdash x\in A\land P(x)}[thm]{7}}
 \proofstepwidestar[1]{P(s)\to s\in\{x\in A\mid P(x)\}}{\rRI{6,8}}
 \proofstepwidestar[1]{s\in\{x\in A\mid P(x)\}\leftrightarrow P(s)}{\rLRI{5,9}}
\closeproofpartwideindR

\proofpartwideindR[Eine Quantorenmarke waehlt genau ihren Index]{i\in\mathbb N\vdash Q(\mathsf q_i,U)=Q_i(U)}{i\in\mathbb N\vdash Q(\mathsf q_i,U)=Q_i(U)}[B48QValue]
 \proofstepwidestar[1]{i\in\mathbb N}{\rA}
 \proofstepwidestar[2]{s\in Q(\mathsf q_i,U)}{\rA}
 \proofstepwidestar[2]{s\in A_D\land\exists j\in\mathbb N\;(\mathsf q_i=\mathsf q_j\land\exists a\in D\;(s[j\mapsto a]\in U))}{\FormulaRefAuto{B48SemanticTerms}[def]{2}}
 \proofstepwidestar[2]{s\in A_D}{\rAEa{3}}
 \proofstepwidestar[2]{\exists j\in\mathbb N\;(\mathsf q_i=\mathsf q_j\land\exists a\in D\;(s[j\mapsto a]\in U))}{\rAEb{3}}
 \proofstepwidestar[6]{j\in\mathbb N\land(\mathsf q_i=\mathsf q_j\land\exists a\in D\;(s[j\mapsto a]\in U))}{\rA}
 \proofstepwidestar[6]{j\in\mathbb N}{\rAEa{6}}
 \proofstepwidestar[6]{\mathsf q_i=\mathsf q_j}{\rAEa{\rAEb{6}}}
 \proofstepwidestar[1,6]{i=j}{\FormulaRefAuto{B48QuantifierRootInjective}[thm]{1,7,8}}
 \proofstepwidestar[1,6]{j=i}{\FormulaRefAuto{a=b\vdash b=a}[thm]{9}}
 \proofstepwidestar[6]{\exists a\in D\;(s[j\mapsto a]\in U)}{\rAEb{\rAEb{6}}}
 \proofstepwidestar[1,6]{\exists a\in D\;(s[i\mapsto a]\in U)}{\rIE{10,11}}
 \proofstepwidestar[1,2]{\exists a\in D\;(s[i\mapsto a]\in U)}{\rEE{5,6:12}}
 \proofstepwidestar[1,2]{s\in Q_i(U)}{\FormulaRefAuto{B48SemanticTerms}[def]{\rAI{4,13}}}
 \proofstepwidestar[1]{s\in Q(\mathsf q_i,U)\to s\in Q_i(U)}{\rRI{2,14}}
 \proofstepwidestar[16]{s\in Q_i(U)}{\rA}
 \proofstepwidestar[16]{s\in A_D\land\exists a\in D\;(s[i\mapsto a]\in U)}{\FormulaRefAuto{B48SemanticTerms}[def]{16}}
 \proofstepwidestar[]{\mathsf q_i=\mathsf q_i}{\rII}
 \proofstepwidestar[1,16]{\exists j\in\mathbb N\;(\mathsf q_i=\mathsf q_j\land\exists a\in D\;(s[j\mapsto a]\in U))}{\rEI{\rAI{1,\rAI{18,\rAEb{17}}}}}
 \proofstepwidestar[1,16]{s\in Q(\mathsf q_i,U)}{\FormulaRefAuto{B48SemanticTerms}[def]{\rAI{\rAEa{17},19}}}
 \proofstepwidestar[1]{s\in Q_i(U)\to s\in Q(\mathsf q_i,U)}{\rRI{16,20}}
 \proofstepwidestar[1]{s\in Q(\mathsf q_i,U)\leftrightarrow s\in Q_i(U)}{\rLRI{15,21}}
 \proofstepwidestar[1]{Q(\mathsf q_i,U)=Q_i(U)}{\FormulaRefAuto{\forall x(x\in A\leftrightarrow x\in B)\eqvdash A=B}[thm]{\rUI{22}}}
\closeproofpartwideindR

\proofpartwideindR[Auswertung einer Quantorenmarke]{i\in\mathbb N\vdash H_M(\mathsf q_i,U,V)=Q_i(V)}{i\in\mathbb N\vdash H_M(\mathsf q_i,U,V)=Q_i(V)}[B48HExists]
 \proofstepwidestar[1]{i\in\mathbb N}{\rA}
 \proofstepwidestar[1]{\mathsf q_i\neq\mathsf n_0}{\FormulaRefAuto{B48QNeNeg}[thm]{1}}
 \proofstepwidestar[1]{\mathsf q_i\neq\mathsf c_0}{\FormulaRefAuto{B48QNeAnd}[thm]{1}}
 \proofstepwidestar[]{t_{\rm q}\in\mathbb N}{\FormulaRefAuto{B48TagType5}[thm]}
 \proofstepwidestar[]{0\in\mathbb N}{\FormulaRefAuto{PeanoZero}[ax]}
 \proofstepwidestar[1]{(t_{\rm q},(i,0))\in\Sigma}{\FormulaRefAuto{B48LabelType}[thm]{4,1,5}}
 \proofstepwidestar[1]{\neg J(\mathsf q_i)}{\FormulaRefAuto{B48LeafNotAndSlot}[thm]{6}}
 \proofstepwidestar[1]{H_M(\mathsf q_i,U,V)=H_{1,M}(\mathsf q_i,U,V)}{\FormulaRefAuto{B48HSkipNeg}[thm]{2}}
 \proofstepwidestar[1]{H_{1,M}(\mathsf q_i,U,V)=H_{2,M}(\mathsf q_i,U,V)}{\FormulaRefAuto{B48HSkipAnd}[thm]{3}}
 \proofstepwidestar[1]{H_{2,M}(\mathsf q_i,U,V)=Q(\mathsf q_i,V)}{\FormulaRefAuto{B48HSkipSlot}[thm]{7}}
 \proofstepwidestar[1]{Q(\mathsf q_i,V)=Q_i(V)}{\FormulaRefAuto{B48QValue}[thm]{1}}
 \proofstepwidestar[1]{H_M(\mathsf q_i,U,V)=Q_i(V)}{\rIE{11,\rIE{10,\rIE{9,8}}}}
\closeproofpartwideindR

\par\medskip\noindent\textbf{D. Mengenbildung und Strukturrekursion}\par
\proofpartwideindR[Die leere Wertemenge]{\varnothing\in\powerset(A)}{\varnothing\in\powerset(A)}[B48EmptyPower]
 \proofstepwidestar[]{\varnothing\subseteq A}{\FormulaRefAuto{\varnothing\subseteq A}[thm]}
 \proofstepwidestar[]{\varnothing\in\powerset(A)}{\FormulaRefAuto{x\in\powerset(A)\eqvdash x\subseteq A}[thm]{1}}
\closeproofpartwideindR

\proofpartwideindR[Schnitt zweier Wertemengen]{U\in\powerset(A),\ V\in\powerset(A)\vdash U\cap V\in\powerset(A)}{U\in\powerset(A),\ V\in\powerset(A)\vdash U\cap V\in\powerset(A)}[B48IntersectionPower]
 \proofstepwidestar[1]{U\in\powerset(A)}{\rA}
 \proofstepwidestar[2]{V\in\powerset(A)}{\rA}
 \proofstepwidestar[1]{U\subseteq A}{\FormulaRefAuto{x\in\powerset(A)\eqvdash x\subseteq A}[thm]{1}}
 \proofstepwidestar[]{U\cap V\subseteq U}{\FormulaRefAuto{A\cap B\subseteq A}[thm]}
 \proofstepwidestar[1]{U\cap V\subseteq A}{\FormulaRefAuto{A\subseteq B,\ B\subseteq C\vdash A\subseteq C}[thm]{4,3}}
 \proofstepwidestar[1]{U\cap V\in\powerset(A)}{\FormulaRefAuto{x\in\powerset(A)\eqvdash x\subseteq A}[thm]{5}}
\closeproofpartwideindR

\proofpartwideindR[Eindeutig zerlegbare Alphabetdaten]{c\in\Sigma\vdash\exists k\in\mathbb N\,\exists i\in\mathbb N\,\exists j\in\mathbb N\;(c=(k,(i,j)))}{c\in\Sigma\vdash\exists k\in\mathbb N\,\exists i\in\mathbb N\,\exists j\in\mathbb N\;(c=(k,(i,j)))}[B48LabelDecomposition]
 \proofstepwidestar[1]{c\in\Sigma}{\rA}
 \proofstepwidestar[1]{\exists k\in\mathbb N\,\exists u\in\mathbb N\times\mathbb N\;(c=(k,u))}{\FormulaRefAuto{x\in A\times B\eqvdash\exists a\in A\,\exists b\in B\,x=(a,b)}[thm]{1}}
 \proofstepwidestar[3]{k\in\mathbb N\land(u\in\mathbb N\times\mathbb N\land c=(k,u))}{\rA}
 \proofstepwidestar[3]{u\in\mathbb N\times\mathbb N}{\rAEa{\rAEb{3}}}
 \proofstepwidestar[3]{\exists i\in\mathbb N\,\exists j\in\mathbb N\;(u=(i,j))}{\FormulaRefAuto{x\in A\times B\eqvdash\exists a\in A\,\exists b\in B\,x=(a,b)}[thm]{4}}
 \proofstepwidestar[6]{i\in\mathbb N\land(j\in\mathbb N\land u=(i,j))}{\rA}
 \proofstepwidestar[3,6]{c=(k,(i,j))}{\rIE{\rAEb{\rAEb{6}},\rAEb{\rAEb{3}}}}
 \proofstepwidestar[3,6]{\exists a\in\mathbb N\,\exists b\in\mathbb N\,\exists d\in\mathbb N\;(c=(a,(b,d)))}{\rEI{\rAI{\rAEa{3},\rEI{\rAI{\rAEa{6},\rEI{\rAI{\rAEa{\rAEb{6}},7}}}}}}}
 \proofstepwidestar[3]{\exists a\in\mathbb N\,\exists b\in\mathbb N\,\exists d\in\mathbb N\;(c=(a,(b,d)))}{\rEEStar{5,6:8}}
 \proofstepwidestar[1]{\exists a\in\mathbb N\,\exists b\in\mathbb N\,\exists d\in\mathbb N\;(c=(a,(b,d)))}{\rEEStar{2,3:9}}
\closeproofpartwideindR

\proofpartwideindR[Atomare Wertemengen sind Mengen von Belegungen]{c\in\Sigma\vdash B_M(c)\in\powerset(A_D)}{c\in\Sigma\vdash B_M(c)\in\powerset(A_D)}[B48BType]
 \proofstepwidestar[1]{c\in\Sigma}{\rA}
 \proofstepwidestar[1]{\exists k\in\mathbb N\,\exists i\in\mathbb N\,\exists j\in\mathbb N\;(c=(k,(i,j)))}{\FormulaRefAuto{B48LabelDecomposition}[thm]{1}}
 \proofstepwidestar[3]{k\in\mathbb N\land(i\in\mathbb N\land(j\in\mathbb N\land c=(k,(i,j))))}{\rA}
 \proofstepwidestar[]{\{s\in A_D\mid s(i)=s(j)\}\in\powerset(A_D)}{\FormulaRefAuto{B48SeparationPower}[thm]}
 \proofstepwidestar[]{\{s\in A_D\mid(s(i),s(j))\in E\}\in\powerset(A_D)}{\FormulaRefAuto{B48SeparationPower}[thm]}
 \proofstepwidestar[]{\varnothing\in\powerset(A_D)}{\FormulaRefAuto{B48EmptyPower}[thm]}
 \proofstepwidestar[]{\BXLVIIIIf{k=t_{\rm m}}{\{s\in A_D\mid(s(i),s(j))\in E\}}{\varnothing}\in\powerset(A_D)}{\FormulaRefAuto{B48IfTyped}[thm]{5,6}}
 \proofstepwidestar[]{\BXLVIIIIf{k=t_{\rm e}}{\{s\in A_D\mid s(i)=s(j)\}}{\BXLVIIIIf{k=t_{\rm m}}{\{s\in A_D\mid(s(i),s(j))\in E\}}{\varnothing}}\in\powerset(A_D)}{\FormulaRefAuto{B48IfTyped}[thm]{4,7}}
 \proofstepwidestar[3]{B_M((k,(i,j)))\in\powerset(A_D)}{\FormulaRefAuto{B48SemanticTerms}[def]{8}}
 \proofstepwidestar[3]{c=(k,(i,j))}{\rAEb{\rAEb{\rAEb{3}}}}
 \proofstepwidestar[3]{(k,(i,j))=c}{\FormulaRefAuto{a=b\vdash b=a}[thm]{10}}
 \proofstepwidestar[3]{B_M(c)\in\powerset(A_D)}{\rIE{11,9}}
 \proofstepwidestar[1]{B_M(c)\in\powerset(A_D)}{\rEEStar{2,3:12}}
\closeproofpartwideindR

\proofpartwideindR[Der Knotenoperator bleibt im Wertemengenbereich]{U\in\powerset(A_D),\ V\in\powerset(A_D)\vdash H_M(t,U,V)\in\powerset(A_D)}{U\in\powerset(A_D),\ V\in\powerset(A_D)\vdash H_M(t,U,V)\in\powerset(A_D)}[B48HType]
 \proofstepwidestar[1]{U\in\powerset(A_D)}{\rA}
 \proofstepwidestar[2]{V\in\powerset(A_D)}{\rA}
 \proofstepwidestar[]{A_D\setminus V\in\powerset(A_D)}{\FormulaRefAuto{B48SeparationPower}[thm]}
 \proofstepwidestar[1,2]{U\cap V\in\powerset(A_D)}{\FormulaRefAuto{B48IntersectionPower}[thm]{1,2}}
 \proofstepwidestar[]{Q(t,V)\in\powerset(A_D)}{\FormulaRefAuto{B48SeparationPower}[thm]}
 \proofstepwidestar[1,2]{\BXLVIIIIf{J(t)}{U\cap V}{Q(t,V)}\in\powerset(A_D)}{\FormulaRefAuto{B48IfTyped}[thm]{4,5}}
 \proofstepwidestar[1,2]{H_{2,M}(t,U,V)\in\powerset(A_D)}{\FormulaRefAuto{B48EvaluationOperators}[def]{6}}
 \proofstepwidestar[1,2]{\BXLVIIIIf{t=\mathsf c_0}{V}{H_{2,M}(t,U,V)}\in\powerset(A_D)}{\FormulaRefAuto{B48IfTyped}[thm]{2,7}}
 \proofstepwidestar[1,2]{H_{1,M}(t,U,V)\in\powerset(A_D)}{\FormulaRefAuto{B48EvaluationOperators}[def]{8}}
 \proofstepwidestar[1,2]{\BXLVIIIIf{t=\mathsf n_0}{A_D\setminus V}{H_{1,M}(t,U,V)}\in\powerset(A_D)}{\FormulaRefAuto{B48IfTyped}[thm]{3,9}}
 \proofstepwidestar[1,2]{H_M(t,U,V)\in\powerset(A_D)}{\FormulaRefAuto{B48EvaluationOperators}[def]{10}}
\closeproofpartwideindR

\proofpartwideindR[Typisierung des Blattauswertungswerts]{c\in\Sigma\vdash(\mathsf l(c),B_M(c))\in Y_D}{c\in\Sigma\vdash(\mathsf l(c),B_M(c))\in Y_D}[B48LeafGraphTerm]
 \proofstepwidestar[1]{c\in\Sigma}{\rA}
 \proofstepwidestar[1]{\mathsf l(c)\in\mathcal T}{\FormulaRefAuto{B48LeafType}[thm]{1}}
 \proofstepwidestar[1]{B_M(c)\in\powerset(A_D)}{\FormulaRefAuto{B48BType}[thm]{1}}
 \proofstepwidestar[1]{(\mathsf l(c),B_M(c))\in Y_D}{\FormulaRefAuto{(a,b)\in A\times B\eqvdash a\in A\land b\in B}[thm]{\rAI{2,3}}}
\closeproofpartwideindR

\proofpartwideindR[Der Blattgraph ist eine Funktion]{f_M\colon\Sigma\to Y_D}{f_M\colon\Sigma\to Y_D}[B48LeafGraphType]
 \proofstepwidestar[1]{c\in\Sigma}{\rA}
 \proofstepwidestar[1]{(\mathsf l(c),B_M(c))\in Y_D}{\FormulaRefAuto{B48LeafGraphTerm}[thm]{1}}
 \proofstepwidestar[]{\forall c\in\Sigma\;((\mathsf l(c),B_M(c))\in Y_D)}{\rUI{\rRI{1,2}}}
 \proofstepwidestar[]{f_M\colon\Sigma\to Y_D}{\FormulaRefAuto{TypedTermGraphFunction}[thm]{3}}
\closeproofpartwideindR

\proofpartwideindR[Werte der Blattauswertungsfunktion]{c\in\Sigma\vdash f_M(c)=(\mathsf l(c),B_M(c))}{c\in\Sigma\vdash f_M(c)=(\mathsf l(c),B_M(c))}[B48LeafGraphValue]
 \proofstepwidestar[1]{c\in\Sigma}{\rA}
 \proofstepwidestar[]{f_M\colon\Sigma\to Y_D}{\FormulaRefAuto{B48LeafGraphType}[thm]}
 \proofstepwidestar[1]{(c,f_M(c))\in f_M}{\FormulaRefAuto{F\colon A\to B,\ x\in A\vdash(x,F(x))\in F}[thm]{2,1}}
 \proofstepwidestar[1]{(c,f_M(c))\in\Sigma\times Y_D\land f_M(c)=(\mathsf l(c),B_M(c))}{\FormulaRefAuto{B48EvaluationOperators}[def]{3}}
 \proofstepwidestar[1]{f_M(c)=(\mathsf l(c),B_M(c))}{\rAEb{4}}
\closeproofpartwideindR

\proofpartwideindR[Typisierung des Knotenauswertungswerts]{(p,U)\in Y_D,\ (q,V)\in Y_D\vdash(\mathsf n(p,q),H_M(p,U,V))\in Y_D}{(p,U)\in Y_D,\ (q,V)\in Y_D\vdash(\mathsf n(p,q),H_M(p,U,V))\in Y_D}[B48NodeGraphTerm]
 \proofstepwidestar[1]{(p,U)\in Y_D}{\rA}
 \proofstepwidestar[2]{(q,V)\in Y_D}{\rA}
 \proofstepwidestar[1]{p\in\mathcal T\land U\in\powerset(A_D)}{\FormulaRefAuto{(a,b)\in A\times B\eqvdash a\in A\land b\in B}[thm]{1}}
 \proofstepwidestar[2]{q\in\mathcal T\land V\in\powerset(A_D)}{\FormulaRefAuto{(a,b)\in A\times B\eqvdash a\in A\land b\in B}[thm]{2}}
 \proofstepwidestar[1,2]{\mathsf n(p,q)\in\mathcal T}{\FormulaRefAuto{B48NodeType}[thm]{\rAEa{3},\rAEa{4}}}
 \proofstepwidestar[1,2]{H_M(p,U,V)\in\powerset(A_D)}{\FormulaRefAuto{B48HType}[thm]{\rAEb{3},\rAEb{4}}}
 \proofstepwidestar[1,2]{(\mathsf n(p,q),H_M(p,U,V))\in Y_D}{\FormulaRefAuto{(a,b)\in A\times B\eqvdash a\in A\land b\in B}[thm]{\rAI{5,6}}}
\closeproofpartwideindR

\proofpartwideindR[Der Knotengraph ist eine Funktion]{g_M\colon Y_D\times Y_D\to Y_D}{g_M\colon Y_D\times Y_D\to Y_D}[B48NodeGraphType]
 \proofstepwidestar[1]{(p,U)\in Y_D}{\rA}
 \proofstepwidestar[2]{(q,V)\in Y_D}{\rA}
 \proofstepwidestar[1,2]{(\mathsf n(p,q),H_M(p,U,V))\in Y_D}{\FormulaRefAuto{B48NodeGraphTerm}[thm]{1,2}}
 \proofstepwidestar[]{\forall(p,U)\in Y_D\,\forall(q,V)\in Y_D\;((\mathsf n(p,q),H_M(p,U,V))\in Y_D)}{\rUI{\rRI{1,\rUI{\rRI{2,3}}}}}
 \proofstepwidestar[]{g_M\colon Y_D\times Y_D\to Y_D}{\FormulaRefAuto{TypedTermGraphFunction}[thm]{4}}
\closeproofpartwideindR

\proofpartwideindR[Werte der Knotenauswertungsfunktion]{(p,U)\in Y_D,\ (q,V)\in Y_D\vdash g_M((p,U),(q,V))=(\mathsf n(p,q),H_M(p,U,V))}{(p,U)\in Y_D,\ (q,V)\in Y_D\vdash g_M((p,U),(q,V))=(\mathsf n(p,q),H_M(p,U,V))}[B48NodeGraphValue]
 \proofstepwidestar[1]{(p,U)\in Y_D}{\rA}
 \proofstepwidestar[2]{(q,V)\in Y_D}{\rA}
 \proofstepwidestar[]{g_M\colon Y_D\times Y_D\to Y_D}{\FormulaRefAuto{B48NodeGraphType}[thm]}
 \proofstepwidestar[1,2]{((p,U),(q,V))\in Y_D\times Y_D}{\FormulaRefAuto{(a,b)\in A\times B\eqvdash a\in A\land b\in B}[thm]{\rAI{1,2}}}
 \proofstepwidestar[1,2]{(((p,U),(q,V)),g_M((p,U),(q,V)))\in g_M}{\FormulaRefAuto{F\colon A\to B,\ x\in A\vdash(x,F(x))\in F}[thm]{3,4}}
 \proofstepwidestar[1,2]{\begin{gathered}(((p,U),(q,V)),g_M((p,U),(q,V)))\\{}\in(Y_D\times Y_D)\times Y_D\\{}\land g_M((p,U),(q,V))=(\mathsf n(p,q),H_M(p,U,V))\end{gathered}}{\FormulaRefAuto{B48EvaluationOperators}[def]{5}}
 \proofstepwidestar[1,2]{g_M((p,U),(q,V))=(\mathsf n(p,q),H_M(p,U,V))}{\rAEb{6}}
\closeproofpartwideindR

\proofpartwideindR[Anwendung des bereits bewiesenen Rekursionssatzes]{\exists e\;\mathsf{Rec}(e,M)}{\exists e\;\mathsf{Rec}(e,M)}[B48TreeEvaluationExists]
 \proofstepwidestar[]{f_M\colon\Sigma\to Y_D}{\FormulaRefAuto{B48LeafGraphType}[thm]}
 \proofstepwidestar[]{g_M\colon Y_D\times Y_D\to Y_D}{\FormulaRefAuto{B48NodeGraphType}[thm]}
 \proofstepwidestar[]{\begin{gathered}\exists!e\colon\mathcal T\to Y_D\;\bigl((\forall c\in\Sigma\;e(\mathsf l(c))=f_M(c))\\{}\land(\forall p\in\mathcal T\,\forall q\in\mathcal T\;e(\mathsf n(p,q))=g_M(e(p),e(q)))\bigr)\end{gathered}}{\FormulaRefAuto{FullPlanarBinaryTreeRecursion}[thm]{1,2}}
 \proofstepwidestar[]{\begin{gathered}\exists e\colon\mathcal T\to Y_D\;\bigl((\forall c\in\Sigma\;e(\mathsf l(c))=f_M(c))\\{}\land(\forall p\in\mathcal T\,\forall q\in\mathcal T\;e(\mathsf n(p,q))=g_M(e(p),e(q)))\bigr)\end{gathered}}{\FormulaRefAuto{\exists! x\,P(x)\vdash \exists x\,P(x)}[thm]{3}}
 \proofstepwidestar[5]{\begin{gathered}e\colon\mathcal T\to Y_D\\{}\land(\forall c\in\Sigma\;e(\mathsf l(c))=f_M(c))\\{}\land(\forall p\in\mathcal T\,\forall q\in\mathcal T\;e(\mathsf n(p,q))=g_M(e(p),e(q)))\end{gathered}}{\rA}
 \proofstepwidestar[5]{e\colon\mathcal T\to Y_D}{\rAEa{5}}
 \proofstepwidestar[5]{\forall c\in\Sigma\;(e(\mathsf l(c))=f_M(c))}{\rAEa{\rAEb{5}}}
 \proofstepwidestar[5]{\forall p\in\mathcal T\,\forall q\in\mathcal T\;(e(\mathsf n(p,q))=g_M(e(p),e(q)))}{\rAEb{\rAEb{5}}}
 \proofstepwidestar[5]{\mathsf{Rec}(e,M)}{\FormulaRefAuto{B48RecursiveEvaluationDef}[def]{6,7,8}}
 \proofstepwidestar[5]{\exists e\;\mathsf{Rec}(e,M)}{\rEI{9}}
 \proofstepwidestar[]{\exists e\;\mathsf{Rec}(e,M)}{\rEE{4,5,10}}
\closeproofpartwideindR

\proofpartwideindR[Rekursionsklausel 1]{\mathsf{Rec}(e,M)\vdash e\colon\mathcal T\to Y_D}{\mathsf{Rec}(e,M)\vdash e\colon\mathcal T\to Y_D}[B48RecType]
 \proofstepwidestar[1]{\mathsf{Rec}(e,M)}{\rA}
 \proofstepwidestar[1]{e\colon\mathcal T\to Y_D}{\FormulaRefAuto{B48RecTypeAxiom}[ax]{1}}
\closeproofpartwideindR

\proofpartwideindR[Rekursionsklausel 2]{\mathsf{Rec}(e,M)\vdash \forall c\in\Sigma\;(e(\mathsf l(c))=f_M(c))}{\mathsf{Rec}(e,M)\vdash \forall c\in\Sigma\;(e(\mathsf l(c))=f_M(c))}[B48RecLeaf]
 \proofstepwidestar[1]{\mathsf{Rec}(e,M)}{\rA}
 \proofstepwidestar[1]{\forall c\in\Sigma\;(e(\mathsf l(c))=f_M(c))}{\FormulaRefAuto{B48RecLeafAxiom}[ax]{1}}
\closeproofpartwideindR

\proofpartwideindR[Rekursionsklausel 3]{\mathsf{Rec}(e,M)\vdash \forall p\in\mathcal T\,\forall q\in\mathcal T\;(e(\mathsf n(p,q))=g_M(e(p),e(q)))}{\mathsf{Rec}(e,M)\vdash \forall p\in\mathcal T\,\forall q\in\mathcal T\;(e(\mathsf n(p,q))=g_M(e(p),e(q)))}[B48RecNode]
 \proofstepwidestar[1]{\mathsf{Rec}(e,M)}{\rA}
 \proofstepwidestar[1]{\forall p\in\mathcal T\,\forall q\in\mathcal T\;(e(\mathsf n(p,q))=g_M(e(p),e(q)))}{\FormulaRefAuto{B48RecNodeAxiom}[ax]{1}}
\closeproofpartwideindR

\proofpartwideindR[Auslesen einer eindeutig bestimmten Wertemenge]{U\in\powerset(A),\ z=(p,U)\vdash U=V_A(z,p)}{U\in\powerset(A),\ z=(p,U)\vdash U=V_A(z,p)}[B48ReadValue]
 \proofstepwidestar[1]{U\in\powerset(A)}{\rA}
 \proofstepwidestar[2]{z=(p,U)}{\rA}
 \proofstepwidestar[1]{U\subseteq A}{\FormulaRefAuto{x\in\powerset(A)\eqvdash x\subseteq A}[thm]{1}}
 \proofstepwidestar[4]{s\in U}{\rA}
 \proofstepwidestar[1,4]{s\in A}{\FormulaRefAuto{A\subseteq B,\ x\in A\vdash x\in B}[thm]{3,4}}
 \proofstepwidestar[1,2,4]{\exists W\in\powerset(A)\;(z=(p,W)\land s\in W)}{\rEI{\rAI{1,\rAI{2,4}}}}
 \proofstepwidestar[1,2,4]{s\in V_A(z,p)}{\FormulaRefAuto{B48EvaluationCandidate}[def]{\rAI{5,6}}}
 \proofstepwidestar[1,2]{s\in U\to s\in V_A(z,p)}{\rRI{4,7}}
 \proofstepwidestar[9]{s\in V_A(z,p)}{\rA}
 \proofstepwidestar[9]{s\in A\land\exists W\in\powerset(A)\;(z=(p,W)\land s\in W)}{\FormulaRefAuto{B48EvaluationCandidate}[def]{9}}
 \proofstepwidestar[9]{\exists W\in\powerset(A)\;(z=(p,W)\land s\in W)}{\rAEb{10}}
 \proofstepwidestar[12]{W\in\powerset(A)\land(z=(p,W)\land s\in W)}{\rA}
 \proofstepwidestar[12]{z=(p,W)}{\rAEa{\rAEb{12}}}
 \proofstepwidestar[2,12]{(p,U)=(p,W)}{\rIE{2,13}}
 \proofstepwidestar[2,12]{U=W}{\FormulaRefAuto{TaggedPairRightInjective}[thm]{14}}
 \proofstepwidestar[2,12]{W=U}{\FormulaRefAuto{a=b\vdash b=a}[thm]{15}}
 \proofstepwidestar[12]{s\in W}{\rAEb{\rAEb{12}}}
 \proofstepwidestar[2,12]{s\in U}{\rIE{16,17}}
 \proofstepwidestar[2,9]{s\in U}{\rEE{11,12:18}}
 \proofstepwidestar[2]{s\in V_A(z,p)\to s\in U}{\rRI{9,19}}
 \proofstepwidestar[1,2]{s\in U\leftrightarrow s\in V_A(z,p)}{\rLRI{8,20}}
 \proofstepwidestar[1,2]{U=V_A(z,p)}{\FormulaRefAuto{\forall x(x\in A\leftrightarrow x\in B)\eqvdash A=B}[thm]{\rUI{21}}}
\closeproofpartwideindR

\proofpartwideindR[Typisierung jeder ausgelesenen Wertemenge]{\beta_e(p)\in\powerset(A_D)}{\beta_e(p)\in\powerset(A_D)}[B48BetaType]
 \proofstepwidestar[]{\beta_e(p)\in\powerset(A_D)}{\FormulaRefAuto{B48SeparationPower}[thm]}
\closeproofpartwideindR

\proofpartwideindR[Blattauswertung mit sichtbarem Baumcode]{\mathsf{Rec}(e,M),\ c\in\Sigma\vdash e(\mathsf l(c))=(\mathsf l(c),B_M(c))}{\mathsf{Rec}(e,M),\ c\in\Sigma\vdash e(\mathsf l(c))=(\mathsf l(c),B_M(c))}[B48EvalLeafEquation]
 \proofstepwidestar[1]{\mathsf{Rec}(e,M)}{\rA}
 \proofstepwidestar[2]{c\in\Sigma}{\rA}
 \proofstepwidestar[1]{\forall d\in\Sigma\;(e(\mathsf l(d))=f_M(d))}{\FormulaRefAuto{B48RecLeaf}[thm]{1}}
 \proofstepwidestar[1,2]{e(\mathsf l(c))=f_M(c)}{\rRE{\rUE{3},2}}
 \proofstepwidestar[2]{f_M(c)=(\mathsf l(c),B_M(c))}{\FormulaRefAuto{B48LeafGraphValue}[thm]{2}}
 \proofstepwidestar[1,2]{e(\mathsf l(c))=(\mathsf l(c),B_M(c))}{\rIE{5,4}}
\closeproofpartwideindR

\proofpartwideindR[Wertemenge eines Blattes]{\mathsf{Rec}(e,M),\ c\in\Sigma\vdash\beta_e(\mathsf l(c))=B_M(c)}{\mathsf{Rec}(e,M),\ c\in\Sigma\vdash\beta_e(\mathsf l(c))=B_M(c)}[B48EvalLeafValue]
 \proofstepwidestar[1]{\mathsf{Rec}(e,M)}{\rA}
 \proofstepwidestar[2]{c\in\Sigma}{\rA}
 \proofstepwidestar[1,2]{e(\mathsf l(c))=(\mathsf l(c),B_M(c))}{\FormulaRefAuto{B48EvalLeafEquation}[thm]{1,2}}
 \proofstepwidestar[2]{B_M(c)\in\powerset(A_D)}{\FormulaRefAuto{B48BType}[thm]{2}}
 \proofstepwidestar[1,2]{B_M(c)=\beta_e(\mathsf l(c))}{\FormulaRefAuto{B48ReadValue}[thm]{4,3}}
 \proofstepwidestar[1,2]{\beta_e(\mathsf l(c))=B_M(c)}{\FormulaRefAuto{a=b\vdash b=a}[thm]{5}}
\closeproofpartwideindR

\proofpartwideindR[Erhaltung des Blattcodes]{\mathsf{Rec}(e,M),\ c\in\Sigma\vdash e(\mathsf l(c))=(\mathsf l(c),\beta_e(\mathsf l(c)))}{\mathsf{Rec}(e,M),\ c\in\Sigma\vdash e(\mathsf l(c))=(\mathsf l(c),\beta_e(\mathsf l(c)))}[B48MetadataLeaf]
 \proofstepwidestar[1]{\mathsf{Rec}(e,M)}{\rA}
 \proofstepwidestar[2]{c\in\Sigma}{\rA}
 \proofstepwidestar[1,2]{e(\mathsf l(c))=(\mathsf l(c),B_M(c))}{\FormulaRefAuto{B48EvalLeafEquation}[thm]{1,2}}
 \proofstepwidestar[1,2]{\beta_e(\mathsf l(c))=B_M(c)}{\FormulaRefAuto{B48EvalLeafValue}[thm]{1,2}}
 \proofstepwidestar[1,2]{B_M(c)=\beta_e(\mathsf l(c))}{\FormulaRefAuto{a=b\vdash b=a}[thm]{4}}
 \proofstepwidestar[1,2]{e(\mathsf l(c))=(\mathsf l(c),\beta_e(\mathsf l(c)))}{\rIE{5,3}}
\closeproofpartwideindR

\proofpartwideindR[Erhaltung des Knotencodes]{\begin{gathered}\mathsf{Rec}(e,M),\ p\in\mathcal T,\ q\in\mathcal T,\\ e(p)=(p,\beta_e(p))\land e(q)=(q,\beta_e(q))\\ \vdash e(\mathsf n(p,q))=(\mathsf n(p,q),\beta_e(\mathsf n(p,q)))\end{gathered}}{\begin{gathered}\mathsf{Rec}(e,M),\ p\in\mathcal T,\ q\in\mathcal T,\\ e(p)=(p,\beta_e(p))\land e(q)=(q,\beta_e(q))\\ \vdash e(\mathsf n(p,q))=(\mathsf n(p,q),\beta_e(\mathsf n(p,q)))\end{gathered}}[B48MetadataNode]
 \proofstepwidestar[1]{\mathsf{Rec}(e,M)}{\rA}
 \proofstepwidestar[2]{p\in\mathcal T}{\rA}
 \proofstepwidestar[3]{q\in\mathcal T}{\rA}
 \proofstepwidestar[4]{e(p)=(p,\beta_e(p))\land e(q)=(q,\beta_e(q))}{\rA}
 \proofstepwidestar[1]{\forall u\in\mathcal T\,\forall v\in\mathcal T\;(e(\mathsf n(u,v))=g_M(e(u),e(v)))}{\FormulaRefAuto{B48RecNode}[thm]{1}}
 \proofstepwidestar[1,2,3]{e(\mathsf n(p,q))=g_M(e(p),e(q))}{\rRE{\rUE{\rRE{\rUE{5},2}},3}}
 \proofstepwidestar[]{\beta_e(p)\in\powerset(A_D)}{\FormulaRefAuto{B48BetaType}[thm]}
 \proofstepwidestar[]{\beta_e(q)\in\powerset(A_D)}{\FormulaRefAuto{B48BetaType}[thm]}
 \proofstepwidestar[2]{(p,\beta_e(p))\in Y_D}{\FormulaRefAuto{(a,b)\in A\times B\eqvdash a\in A\land b\in B}[thm]{\rAI{2,7}}}
 \proofstepwidestar[3]{(q,\beta_e(q))\in Y_D}{\FormulaRefAuto{(a,b)\in A\times B\eqvdash a\in A\land b\in B}[thm]{\rAI{3,8}}}
 \proofstepwidestar[2,3]{g_M((p,\beta_e(p)),(q,\beta_e(q)))=(\mathsf n(p,q),H_M(p,\beta_e(p),\beta_e(q)))}{\FormulaRefAuto{B48NodeGraphValue}[thm]{9,10}}
 \proofstepwidestar[1,2,3,4]{e(\mathsf n(p,q))=g_M((p,\beta_e(p)),(q,\beta_e(q)))}{\rIE{\rAEb{4},\rIE{\rAEa{4},6}}}
 \proofstepwidestar[1,2,3,4]{e(\mathsf n(p,q))=(\mathsf n(p,q),H_M(p,\beta_e(p),\beta_e(q)))}{\rIE{11,12}}
 \proofstepwidestar[]{H_M(p,\beta_e(p),\beta_e(q))\in\powerset(A_D)}{\FormulaRefAuto{B48HType}[thm]{7,8}}
 \proofstepwidestar[1,2,3,4]{H_M(p,\beta_e(p),\beta_e(q))=\beta_e(\mathsf n(p,q))}{\FormulaRefAuto{B48ReadValue}[thm]{14,13}}
 \proofstepwidestar[1,2,3,4]{e(\mathsf n(p,q))=(\mathsf n(p,q),\beta_e(\mathsf n(p,q)))}{\rIE{15,13}}
\closeproofpartwideindR

\proofpartwideindR[Strukturinduktion ueber alle Hilfsbaeume]{\mathsf{Rec}(e,M)\vdash\forall p\in\mathcal T\;(e(p)=(p,\beta_e(p)))}{\mathsf{Rec}(e,M)\vdash\forall p\in\mathcal T\;(e(p)=(p,\beta_e(p)))}[B48MetadataAll]
 \proofstepwidestar[1]{\mathsf{Rec}(e,M)}{\rA}
 \proofstepwidestar[2]{c\in\Sigma}{\rA}
 \proofstepwidestar[1,2]{e(\mathsf l(c))=(\mathsf l(c),\beta_e(\mathsf l(c)))}{\FormulaRefAuto{B48MetadataLeaf}[thm]{1,2}}
 \proofstepwidestar[1]{\forall c\in\Sigma\;(e(\mathsf l(c))=(\mathsf l(c),\beta_e(\mathsf l(c))))}{\rUI{\rRI{2,3}}}
 \proofstepwidestar[5]{p\in\mathcal T}{\rA}
 \proofstepwidestar[6]{q\in\mathcal T}{\rA}
 \proofstepwidestar[7]{e(p)=(p,\beta_e(p))\land e(q)=(q,\beta_e(q))}{\rA}
 \proofstepwidestar[1,5,6,7]{e(\mathsf n(p,q))=(\mathsf n(p,q),\beta_e(\mathsf n(p,q)))}{\FormulaRefAuto{B48MetadataNode}[thm]{1,5,6,7}}
 \proofstepwidestar[1]{\begin{gathered}\forall p\in\mathcal T\,\forall q\in\mathcal T\;\\
 ((e(p)=(p,\beta_e(p))\land e(q)=(q,\beta_e(q)))\\
 \to e(\mathsf n(p,q))=(\mathsf n(p,q),\beta_e(\mathsf n(p,q))))\end{gathered}}{\rUI{\rRI{5,\rUI{\rRI{6,\rRI{7,8}}}}}}
 \proofstepwidestar[1]{\forall p\in\mathcal T\;(e(p)=(p,\beta_e(p)))}{\FormulaRefAuto{FullPlanarBinaryTreeStructuralInduction}[thm]{4,9}}
\closeproofpartwideindR

\proofpartwideindR[Wertemenge eines Knotens]{\mathsf{Rec}(e,M),\ p\in\mathcal T,\ q\in\mathcal T\vdash\beta_e(\mathsf n(p,q))=H_M(p,\beta_e(p),\beta_e(q))}{\mathsf{Rec}(e,M),\ p\in\mathcal T,\ q\in\mathcal T\vdash\beta_e(\mathsf n(p,q))=H_M(p,\beta_e(p),\beta_e(q))}[B48EvalNodeValue]
 \proofstepwidestar[1]{\mathsf{Rec}(e,M)}{\rA}
 \proofstepwidestar[2]{p\in\mathcal T}{\rA}
 \proofstepwidestar[3]{q\in\mathcal T}{\rA}
 \proofstepwidestar[1]{\forall t\in\mathcal T\;(e(t)=(t,\beta_e(t)))}{\FormulaRefAuto{B48MetadataAll}[thm]{1}}
 \proofstepwidestar[1,2]{e(p)=(p,\beta_e(p))}{\rRE{\rUE{4},2}}
 \proofstepwidestar[1,3]{e(q)=(q,\beta_e(q))}{\rRE{\rUE{4},3}}
 \proofstepwidestar[1]{\forall u\in\mathcal T\,\forall v\in\mathcal T\;(e(\mathsf n(u,v))=g_M(e(u),e(v)))}{\FormulaRefAuto{B48RecNode}[thm]{1}}
 \proofstepwidestar[1,2,3]{e(\mathsf n(p,q))=g_M(e(p),e(q))}{\rRE{\rUE{\rRE{\rUE{7},2}},3}}
 \proofstepwidestar[]{\beta_e(p)\in\powerset(A_D)}{\FormulaRefAuto{B48BetaType}[thm]}
 \proofstepwidestar[]{\beta_e(q)\in\powerset(A_D)}{\FormulaRefAuto{B48BetaType}[thm]}
 \proofstepwidestar[2]{(p,\beta_e(p))\in Y_D}{\FormulaRefAuto{(a,b)\in A\times B\eqvdash a\in A\land b\in B}[thm]{\rAI{2,9}}}
 \proofstepwidestar[3]{(q,\beta_e(q))\in Y_D}{\FormulaRefAuto{(a,b)\in A\times B\eqvdash a\in A\land b\in B}[thm]{\rAI{3,10}}}
 \proofstepwidestar[2,3]{g_M((p,\beta_e(p)),(q,\beta_e(q)))=(\mathsf n(p,q),H_M(p,\beta_e(p),\beta_e(q)))}{\FormulaRefAuto{B48NodeGraphValue}[thm]{11,12}}
 \proofstepwidestar[1,2,3]{e(\mathsf n(p,q))=(\mathsf n(p,q),H_M(p,\beta_e(p),\beta_e(q)))}{\rIE{13,\rIE{6,\rIE{5,8}}}}
 \proofstepwidestar[]{H_M(p,\beta_e(p),\beta_e(q))\in\powerset(A_D)}{\FormulaRefAuto{B48HType}[thm]{9,10}}
 \proofstepwidestar[1,2,3]{H_M(p,\beta_e(p),\beta_e(q))=\beta_e(\mathsf n(p,q))}{\FormulaRefAuto{B48ReadValue}[thm]{15,14}}
 \proofstepwidestar[1,2,3]{\beta_e(\mathsf n(p,q))=H_M(p,\beta_e(p),\beta_e(q))}{\FormulaRefAuto{a=b\vdash b=a}[thm]{16}}
\closeproofpartwideindR

\par\medskip\noindent\textbf{E. Nachweis der Erfuellungsbedingungen}\par
\proofpartwideindR[Negationsgleichung der konstruierten Wertemengen]{\mathsf{Rec}(e,M),\ p\in\mathcal T\vdash\beta_e(\mathsf{neg}(p))=A_D\setminus\beta_e(p)}{\mathsf{Rec}(e,M),\ p\in\mathcal T\vdash\beta_e(\mathsf{neg}(p))=A_D\setminus\beta_e(p)}[B48BetaNeg]
 \proofstepwidestar[1]{\mathsf{Rec}(e,M)}{\rA}
 \proofstepwidestar[2]{p\in\mathcal T}{\rA}
 \proofstepwidestar[]{\mathsf n_0\in\mathcal T}{\FormulaRefAuto{B48NegRootType}[thm]}
 \proofstepwidestar[1,2]{\beta_e(\mathsf n(\mathsf n_0,p))=H_M(\mathsf n_0,\beta_e(\mathsf n_0),\beta_e(p))}{\FormulaRefAuto{B48EvalNodeValue}[thm]{1,3,2}}
 \proofstepwidestar[]{H_M(\mathsf n_0,\beta_e(\mathsf n_0),\beta_e(p))=A_D\setminus\beta_e(p)}{\FormulaRefAuto{B48HNeg}[thm]}
 \proofstepwidestar[1,2]{\beta_e(\mathsf{neg}(p))=A_D\setminus\beta_e(p)}{\rIE{5,4}}
\closeproofpartwideindR

\proofpartwideindR[Konjunktionsgleichung der konstruierten Wertemengen]{\mathsf{Rec}(e,M),\ p\in\mathcal T,\ q\in\mathcal T\vdash\beta_e(\mathsf{and}(p,q))=\beta_e(p)\cap\beta_e(q)}{\mathsf{Rec}(e,M),\ p\in\mathcal T,\ q\in\mathcal T\vdash\beta_e(\mathsf{and}(p,q))=\beta_e(p)\cap\beta_e(q)}[B48BetaAnd]
 \proofstepwidestar[1]{\mathsf{Rec}(e,M)}{\rA}
 \proofstepwidestar[2]{p\in\mathcal T}{\rA}
 \proofstepwidestar[3]{q\in\mathcal T}{\rA}
 \proofstepwidestar[]{\mathsf c_0\in\mathcal T}{\FormulaRefAuto{B48AndRootType}[thm]}
 \proofstepwidestar[2]{\mathsf n(\mathsf c_0,p)\in\mathcal T}{\FormulaRefAuto{B48NodeType}[thm]{4,2}}
 \proofstepwidestar[1,2]{\beta_e(\mathsf n(\mathsf c_0,p))=H_M(\mathsf c_0,\beta_e(\mathsf c_0),\beta_e(p))}{\FormulaRefAuto{B48EvalNodeValue}[thm]{1,4,2}}
 \proofstepwidestar[]{H_M(\mathsf c_0,\beta_e(\mathsf c_0),\beta_e(p))=\beta_e(p)}{\FormulaRefAuto{B48HAndPrefix}[thm]}
 \proofstepwidestar[1,2]{\beta_e(\mathsf n(\mathsf c_0,p))=\beta_e(p)}{\rIE{7,6}}
 \proofstepwidestar[1,2,3]{\beta_e(\mathsf{and}(p,q))=H_M(\mathsf n(\mathsf c_0,p),\beta_e(\mathsf n(\mathsf c_0,p)),\beta_e(q))}{\FormulaRefAuto{B48EvalNodeValue}[thm]{1,5,3}}
 \proofstepwidestar[2]{H_M(\mathsf n(\mathsf c_0,p),\beta_e(\mathsf n(\mathsf c_0,p)),\beta_e(q))=\beta_e(\mathsf n(\mathsf c_0,p))\cap\beta_e(q)}{\FormulaRefAuto{B48HAnd}[thm]{2}}
 \proofstepwidestar[1,2,3]{\beta_e(\mathsf{and}(p,q))=\beta_e(\mathsf n(\mathsf c_0,p))\cap\beta_e(q)}{\rIE{10,9}}
 \proofstepwidestar[1,2,3]{\beta_e(\mathsf{and}(p,q))=\beta_e(p)\cap\beta_e(q)}{\rIE{8,11}}
\closeproofpartwideindR

\proofpartwideindR[Existenzgleichung der konstruierten Wertemengen]{\mathsf{Rec}(e,M),\ i\in\mathbb N,\ p\in\mathcal T\vdash\beta_e(\mathsf{ex}_i(p))=Q_i(\beta_e(p))}{\mathsf{Rec}(e,M),\ i\in\mathbb N,\ p\in\mathcal T\vdash\beta_e(\mathsf{ex}_i(p))=Q_i(\beta_e(p))}[B48BetaExists]
 \proofstepwidestar[1]{\mathsf{Rec}(e,M)}{\rA}
 \proofstepwidestar[2]{i\in\mathbb N}{\rA}
 \proofstepwidestar[3]{p\in\mathcal T}{\rA}
 \proofstepwidestar[2]{\mathsf q_i\in\mathcal T}{\FormulaRefAuto{B48ExistsRootType}[thm]{2}}
 \proofstepwidestar[1,2,3]{\beta_e(\mathsf{ex}_i(p))=H_M(\mathsf q_i,\beta_e(\mathsf q_i),\beta_e(p))}{\FormulaRefAuto{B48EvalNodeValue}[thm]{1,4,3}}
 \proofstepwidestar[2]{H_M(\mathsf q_i,\beta_e(\mathsf q_i),\beta_e(p))=Q_i(\beta_e(p))}{\FormulaRefAuto{B48HExists}[thm]{2}}
 \proofstepwidestar[1,2,3]{\beta_e(\mathsf{ex}_i(p))=Q_i(\beta_e(p))}{\rIE{6,5}}
\closeproofpartwideindR

\proofpartwideindR[Der Kandidat als Formel-Belegungsrelation]{S_e\subseteq\mathcal F\times A_D}{S_e\subseteq\mathcal F\times A_D}[B48CandidateSubset]
 \proofstepwidestar[]{S_e\subseteq\mathcal F\times A_D}{\FormulaRefAuto{\{x\in A\mid P(x)\}\subseteq A}[thm]}
\closeproofpartwideindR

\proofpartwideindR[Fasern des Kandidaten]{p\in\mathcal F\vdash R_{S_e}(p)=\beta_e(p)}{p\in\mathcal F\vdash R_{S_e}(p)=\beta_e(p)}[B48CandidateRow]
 \proofstepwidestar[1]{p\in\mathcal F}{\rA}
 \proofstepwidestar[2]{s\in R_{S_e}(p)}{\rA}
 \proofstepwidestar[2]{s\in A_D\land(p,s)\in S_e}{\FormulaRefAuto{B48SemanticTerms}[def]{2}}
 \proofstepwidestar[2]{(p,s)\in S_e}{\rAEb{3}}
 \proofstepwidestar[2]{(p,s)\in\mathcal F\times A_D\land s\in\beta_e(p)}{\FormulaRefAuto{B48EvaluationCandidate}[def]{4}}
 \proofstepwidestar[2]{s\in\beta_e(p)}{\rAEb{5}}
 \proofstepwidestar[]{s\in R_{S_e}(p)\to s\in\beta_e(p)}{\rRI{2,6}}
 \proofstepwidestar[8]{s\in\beta_e(p)}{\rA}
 \proofstepwidestar[]{\beta_e(p)\in\powerset(A_D)}{\FormulaRefAuto{B48BetaType}[thm]}
 \proofstepwidestar[]{\beta_e(p)\subseteq A_D}{\FormulaRefAuto{x\in\powerset(A)\eqvdash x\subseteq A}[thm]{9}}
 \proofstepwidestar[8]{s\in A_D}{\FormulaRefAuto{A\subseteq B,\ x\in A\vdash x\in B}[thm]{10,8}}
 \proofstepwidestar[1,8]{(p,s)\in\mathcal F\times A_D}{\FormulaRefAuto{(a,b)\in A\times B\eqvdash a\in A\land b\in B}[thm]{\rAI{1,11}}}
 \proofstepwidestar[1,8]{(p,s)\in S_e}{\FormulaRefAuto{B48EvaluationCandidate}[def]{\rAI{12,8}}}
 \proofstepwidestar[1,8]{s\in R_{S_e}(p)}{\FormulaRefAuto{B48SemanticTerms}[def]{\rAI{11,13}}}
 \proofstepwidestar[1]{s\in\beta_e(p)\to s\in R_{S_e}(p)}{\rRI{8,14}}
 \proofstepwidestar[1]{s\in R_{S_e}(p)\leftrightarrow s\in\beta_e(p)}{\rLRI{7,15}}
 \proofstepwidestar[1]{R_{S_e}(p)=\beta_e(p)}{\FormulaRefAuto{\forall x(x\in A\leftrightarrow x\in B)\eqvdash A=B}[thm]{\rUI{16}}}
\closeproofpartwideindR

\proofpartwideindR[Der Kandidat erfuellt die atomare Klausel]{\mathsf{Rec}(e,M)\vdash C_{\rm at}(M,S_e)}{\mathsf{Rec}(e,M)\vdash C_{\rm at}(M,S_e)}[B48CandidateAtoms]
 \proofstepwidestar[1]{\mathsf{Rec}(e,M)}{\rA}
 \proofstepwidestar[2]{c\in\Sigma_{\rm at}}{\rA}
 \proofstepwidestar[2]{c\in\Sigma}{\FormulaRefAuto{B48AtomicLabelType}[thm]{2}}
 \proofstepwidestar[2]{\mathsf l(c)\in\mathcal F}{\FormulaRefAuto{B48FormulaAtom}[thm]{2}}
 \proofstepwidestar[2]{R_{S_e}(\mathsf l(c))=\beta_e(\mathsf l(c))}{\FormulaRefAuto{B48CandidateRow}[thm]{4}}
 \proofstepwidestar[1,2]{\beta_e(\mathsf l(c))=B_M(c)}{\FormulaRefAuto{B48EvalLeafValue}[thm]{1,3}}
 \proofstepwidestar[1,2]{R_{S_e}(\mathsf l(c))=B_M(c)}{\rIE{6,5}}
 \proofstepwidestar[1]{\forall c\in\Sigma_{\rm at}\;(R_{S_e}(\mathsf l(c))=B_M(c))}{\rUI{\rRI{2,7}}}
 \proofstepwidestar[1]{C_{\rm at}(M,S_e)}{\FormulaRefAuto{B48SatisfactionAtomsDef}[def]{8}}
\closeproofpartwideindR

\proofpartwideindR[Der Kandidat erfuellt die Negationsklausel]{\mathsf{Rec}(e,M)\vdash C_{\rm n}(M,S_e)}{\mathsf{Rec}(e,M)\vdash C_{\rm n}(M,S_e)}[B48CandidateNeg]
 \proofstepwidestar[1]{\mathsf{Rec}(e,M)}{\rA}
 \proofstepwidestar[2]{p\in\mathcal F}{\rA}
 \proofstepwidestar[2]{p\in\mathcal T}{\FormulaRefAuto{B48FormulaInTree}[thm]{2}}
 \proofstepwidestar[2]{\mathsf{neg}(p)\in\mathcal F}{\FormulaRefAuto{B48FormulaNeg}[thm]{2}}
 \proofstepwidestar[2]{R_{S_e}(\mathsf{neg}(p))=\beta_e(\mathsf{neg}(p))}{\FormulaRefAuto{B48CandidateRow}[thm]{4}}
 \proofstepwidestar[1,2]{\beta_e(\mathsf{neg}(p))=A_D\setminus\beta_e(p)}{\FormulaRefAuto{B48BetaNeg}[thm]{1,3}}
 \proofstepwidestar[2]{R_{S_e}(p)=\beta_e(p)}{\FormulaRefAuto{B48CandidateRow}[thm]{2}}
 \proofstepwidestar[2]{\beta_e(p)=R_{S_e}(p)}{\FormulaRefAuto{a=b\vdash b=a}[thm]{7}}
 \proofstepwidestar[1,2]{R_{S_e}(\mathsf{neg}(p))=A_D\setminus R_{S_e}(p)}{\rIE{8,\rIE{6,5}}}
 \proofstepwidestar[1]{\forall p\in\mathcal F\;(R_{S_e}(\mathsf{neg}(p))=A_D\setminus R_{S_e}(p))}{\rUI{\rRI{2,9}}}
 \proofstepwidestar[1]{C_{\rm n}(M,S_e)}{\FormulaRefAuto{B48SatisfactionNegDef}[def]{10}}
\closeproofpartwideindR

\proofpartwideindR[Der Kandidat erfuellt die Konjunktionsklausel]{\mathsf{Rec}(e,M)\vdash C_{\rm c}(M,S_e)}{\mathsf{Rec}(e,M)\vdash C_{\rm c}(M,S_e)}[B48CandidateAnd]
 \proofstepwidestar[1]{\mathsf{Rec}(e,M)}{\rA}
 \proofstepwidestar[2]{p\in\mathcal F}{\rA}
 \proofstepwidestar[3]{q\in\mathcal F}{\rA}
 \proofstepwidestar[2]{p\in\mathcal T}{\FormulaRefAuto{B48FormulaInTree}[thm]{2}}
 \proofstepwidestar[3]{q\in\mathcal T}{\FormulaRefAuto{B48FormulaInTree}[thm]{3}}
 \proofstepwidestar[2,3]{\mathsf{and}(p,q)\in\mathcal F}{\FormulaRefAuto{B48FormulaAnd}[thm]{2,3}}
 \proofstepwidestar[2,3]{R_{S_e}(\mathsf{and}(p,q))=\beta_e(\mathsf{and}(p,q))}{\FormulaRefAuto{B48CandidateRow}[thm]{6}}
 \proofstepwidestar[1,2,3]{\beta_e(\mathsf{and}(p,q))=\beta_e(p)\cap\beta_e(q)}{\FormulaRefAuto{B48BetaAnd}[thm]{1,4,5}}
 \proofstepwidestar[2]{R_{S_e}(p)=\beta_e(p)}{\FormulaRefAuto{B48CandidateRow}[thm]{2}}
 \proofstepwidestar[3]{R_{S_e}(q)=\beta_e(q)}{\FormulaRefAuto{B48CandidateRow}[thm]{3}}
 \proofstepwidestar[2]{\beta_e(p)=R_{S_e}(p)}{\FormulaRefAuto{a=b\vdash b=a}[thm]{9}}
 \proofstepwidestar[3]{\beta_e(q)=R_{S_e}(q)}{\FormulaRefAuto{a=b\vdash b=a}[thm]{10}}
 \proofstepwidestar[1,2,3]{R_{S_e}(\mathsf{and}(p,q))=R_{S_e}(p)\cap R_{S_e}(q)}{\rIE{12,\rIE{11,\rIE{8,7}}}}
 \proofstepwidestar[1]{\forall p\in\mathcal F\,\forall q\in\mathcal F\;(R_{S_e}(\mathsf{and}(p,q))=R_{S_e}(p)\cap R_{S_e}(q))}{\rUI{\rRI{2,\rUI{\rRI{3,13}}}}}
 \proofstepwidestar[1]{C_{\rm c}(M,S_e)}{\FormulaRefAuto{B48SatisfactionAndDef}[def]{14}}
\closeproofpartwideindR

\proofpartwideindR[Der Kandidat erfuellt die Quantorenklausel]{\mathsf{Rec}(e,M)\vdash C_{\rm q}(M,S_e)}{\mathsf{Rec}(e,M)\vdash C_{\rm q}(M,S_e)}[B48CandidateExists]
 \proofstepwidestar[1]{\mathsf{Rec}(e,M)}{\rA}
 \proofstepwidestar[2]{i\in\mathbb N}{\rA}
 \proofstepwidestar[3]{p\in\mathcal F}{\rA}
 \proofstepwidestar[3]{p\in\mathcal T}{\FormulaRefAuto{B48FormulaInTree}[thm]{3}}
 \proofstepwidestar[2,3]{\mathsf{ex}_i(p)\in\mathcal F}{\FormulaRefAuto{B48FormulaExists}[thm]{2,3}}
 \proofstepwidestar[2,3]{R_{S_e}(\mathsf{ex}_i(p))=\beta_e(\mathsf{ex}_i(p))}{\FormulaRefAuto{B48CandidateRow}[thm]{5}}
 \proofstepwidestar[1,2,3]{\beta_e(\mathsf{ex}_i(p))=Q_i(\beta_e(p))}{\FormulaRefAuto{B48BetaExists}[thm]{1,2,4}}
 \proofstepwidestar[3]{R_{S_e}(p)=\beta_e(p)}{\FormulaRefAuto{B48CandidateRow}[thm]{3}}
 \proofstepwidestar[3]{\beta_e(p)=R_{S_e}(p)}{\FormulaRefAuto{a=b\vdash b=a}[thm]{8}}
 \proofstepwidestar[1,2,3]{R_{S_e}(\mathsf{ex}_i(p))=Q_i(R_{S_e}(p))}{\rIE{9,\rIE{7,6}}}
 \proofstepwidestar[1]{\forall i\in\mathbb N\,\forall p\in\mathcal F\;(R_{S_e}(\mathsf{ex}_i(p))=Q_i(R_{S_e}(p)))}{\rUI{\rRI{2,\rUI{\rRI{3,10}}}}}
 \proofstepwidestar[1]{C_{\rm q}(M,S_e)}{\FormulaRefAuto{B48SatisfactionExistsDef}[def]{11}}
\closeproofpartwideindR

\proofpartwideindR[Der konstruierte Kandidat ist eine Erfuellungsrelation]{\mathsf{Str}(M),\ \mathsf{Rec}(e,M)\vdash\mathsf{Erf}(M,S_e)}{\mathsf{Str}(M),\ \mathsf{Rec}(e,M)\vdash\mathsf{Erf}(M,S_e)}[B48CandidateSatisfies]
 \proofstepwidestar[1]{\mathsf{Str}(M)}{\rA}
 \proofstepwidestar[2]{\mathsf{Rec}(e,M)}{\rA}
 \proofstepwidestar[]{S_e\subseteq\mathcal F\times A_D}{\FormulaRefAuto{B48CandidateSubset}[thm]}
 \proofstepwidestar[2]{C_{\rm at}(M,S_e)}{\FormulaRefAuto{B48CandidateAtoms}[thm]{2}}
 \proofstepwidestar[2]{C_{\rm n}(M,S_e)}{\FormulaRefAuto{B48CandidateNeg}[thm]{2}}
 \proofstepwidestar[2]{C_{\rm c}(M,S_e)}{\FormulaRefAuto{B48CandidateAnd}[thm]{2}}
 \proofstepwidestar[2]{C_{\rm q}(M,S_e)}{\FormulaRefAuto{B48CandidateExists}[thm]{2}}
 \proofstepwidestar[1,2]{\mathsf{Erf}(M,S_e)}{\FormulaRefAuto{B48SatisfactionStructure}[def]{1,3,4,5,6,7}}
\closeproofpartwideindR

\par\medskip\noindent\textbf{F. Eindeutigkeit aus den Erfuellungsbedingungen}\par
\proofpartwideindR[Bedingung 1 einer Erfuellungsrelation]{\mathsf{Erf}(M,S)\vdash \mathsf{Str}(M)}{\mathsf{Erf}(M,S)\vdash \mathsf{Str}(M)}[B48ErfStructure]
 \proofstepwidestar[1]{\mathsf{Erf}(M,S)}{\rA}
 \proofstepwidestar[1]{\mathsf{Str}(M)}{\FormulaRefAuto{B48ErfStructureAxiom}[ax]{1}}
\closeproofpartwideindR

\proofpartwideindR[Bedingung 2 einer Erfuellungsrelation]{\mathsf{Erf}(M,S)\vdash S\subseteq\mathcal F\times A_D}{\mathsf{Erf}(M,S)\vdash S\subseteq\mathcal F\times A_D}[B48ErfSubset]
 \proofstepwidestar[1]{\mathsf{Erf}(M,S)}{\rA}
 \proofstepwidestar[1]{S\subseteq\mathcal F\times A_D}{\FormulaRefAuto{B48ErfSubsetAxiom}[ax]{1}}
\closeproofpartwideindR

\proofpartwideindR[Bedingung 3 einer Erfuellungsrelation]{\mathsf{Erf}(M,S)\vdash C_{\rm at}(M,S)}{\mathsf{Erf}(M,S)\vdash C_{\rm at}(M,S)}[B48ErfAtoms]
 \proofstepwidestar[1]{\mathsf{Erf}(M,S)}{\rA}
 \proofstepwidestar[1]{\forall c\in\Sigma_{\rm at}\;(R_S(\mathsf l(c))=B_M(c))}{\FormulaRefAuto{B48ErfAtomsAxiom}[ax]{1}}
 \proofstepwidestar[1]{C_{\rm at}(M,S)}{\FormulaRefAuto{B48SatisfactionAtomsDef}[def]{2}}
\closeproofpartwideindR

\proofpartwideindR[Bedingung 4 einer Erfuellungsrelation]{\mathsf{Erf}(M,S)\vdash C_{\rm n}(M,S)}{\mathsf{Erf}(M,S)\vdash C_{\rm n}(M,S)}[B48ErfNeg]
 \proofstepwidestar[1]{\mathsf{Erf}(M,S)}{\rA}
 \proofstepwidestar[1]{\forall p\in\mathcal F\;(R_S(\mathsf{neg}(p))=A_D\setminus R_S(p))}{\FormulaRefAuto{B48ErfNegAxiom}[ax]{1}}
 \proofstepwidestar[1]{C_{\rm n}(M,S)}{\FormulaRefAuto{B48SatisfactionNegDef}[def]{2}}
\closeproofpartwideindR

\proofpartwideindR[Bedingung 5 einer Erfuellungsrelation]{\mathsf{Erf}(M,S)\vdash C_{\rm c}(M,S)}{\mathsf{Erf}(M,S)\vdash C_{\rm c}(M,S)}[B48ErfAnd]
 \proofstepwidestar[1]{\mathsf{Erf}(M,S)}{\rA}
 \proofstepwidestar[1]{\forall p\in\mathcal F\,\forall q\in\mathcal F\;(R_S(\mathsf{and}(p,q))=R_S(p)\cap R_S(q))}{\FormulaRefAuto{B48ErfAndAxiom}[ax]{1}}
 \proofstepwidestar[1]{C_{\rm c}(M,S)}{\FormulaRefAuto{B48SatisfactionAndDef}[def]{2}}
\closeproofpartwideindR

\proofpartwideindR[Bedingung 6 einer Erfuellungsrelation]{\mathsf{Erf}(M,S)\vdash C_{\rm q}(M,S)}{\mathsf{Erf}(M,S)\vdash C_{\rm q}(M,S)}[B48ErfExists]
 \proofstepwidestar[1]{\mathsf{Erf}(M,S)}{\rA}
 \proofstepwidestar[1]{\forall i\in\mathbb N\,\forall p\in\mathcal F\;(R_S(\mathsf{ex}_i(p))=Q_i(R_S(p)))}{\FormulaRefAuto{B48ErfExistsAxiom}[ax]{1}}
 \proofstepwidestar[1]{C_{\rm q}(M,S)}{\FormulaRefAuto{B48SatisfactionExistsDef}[def]{2}}
\closeproofpartwideindR

\proofpartwideindR[Fasergleichung der Atome]{\mathsf{Erf}(M,S),\ c\in\Sigma_{\rm at}\vdash R_S(\mathsf l(c))=B_M(c)}{\mathsf{Erf}(M,S),\ c\in\Sigma_{\rm at}\vdash R_S(\mathsf l(c))=B_M(c)}[B48ErfAtomRow]
 \proofstepwidestar[1]{\mathsf{Erf}(M,S)}{\rA}
 \proofstepwidestar[2]{c\in\Sigma_{\rm at}}{\rA}
 \proofstepwidestar[1]{C_{\rm at}(M,S)}{\FormulaRefAuto{B48ErfAtoms}[thm]{1}}
 \proofstepwidestar[1]{\forall c\in\Sigma_{\rm at}\;(R_S(\mathsf l(c))=B_M(c))}{\FormulaRefAuto{B48SatisfactionAtomsDef}[def]{3}}
 \proofstepwidestar[1,2]{R_S(\mathsf l(c))=B_M(c)}{\rRE{\rUE{4},2}}
\closeproofpartwideindR

\proofpartwideindR[Fasergleichung der Negation]{\mathsf{Erf}(M,S),\ p\in\mathcal F\vdash R_S(\mathsf{neg}(p))=A_D\setminus R_S(p)}{\mathsf{Erf}(M,S),\ p\in\mathcal F\vdash R_S(\mathsf{neg}(p))=A_D\setminus R_S(p)}[B48ErfNegRow]
 \proofstepwidestar[1]{\mathsf{Erf}(M,S)}{\rA}
 \proofstepwidestar[2]{p\in\mathcal F}{\rA}
 \proofstepwidestar[1]{C_{\rm n}(M,S)}{\FormulaRefAuto{B48ErfNeg}[thm]{1}}
 \proofstepwidestar[1]{\forall p\in\mathcal F\;(R_S(\mathsf{neg}(p))=A_D\setminus R_S(p))}{\FormulaRefAuto{B48SatisfactionNegDef}[def]{3}}
 \proofstepwidestar[1,2]{R_S(\mathsf{neg}(p))=A_D\setminus R_S(p)}{\rRE{\rUE{4},2}}
\closeproofpartwideindR

\proofpartwideindR[Fasergleichung der Konjunktion]{\mathsf{Erf}(M,S),\ p\in\mathcal F,\ q\in\mathcal F\vdash R_S(\mathsf{and}(p,q))=R_S(p)\cap R_S(q)}{\mathsf{Erf}(M,S),\ p\in\mathcal F,\ q\in\mathcal F\vdash R_S(\mathsf{and}(p,q))=R_S(p)\cap R_S(q)}[B48ErfAndRow]
 \proofstepwidestar[1]{\mathsf{Erf}(M,S)}{\rA}
 \proofstepwidestar[2]{p\in\mathcal F}{\rA}
 \proofstepwidestar[3]{q\in\mathcal F}{\rA}
 \proofstepwidestar[1]{C_{\rm c}(M,S)}{\FormulaRefAuto{B48ErfAnd}[thm]{1}}
 \proofstepwidestar[1]{\forall u\in\mathcal F\,\forall v\in\mathcal F\;(R_S(\mathsf{and}(u,v))=R_S(u)\cap R_S(v))}{\FormulaRefAuto{B48SatisfactionAndDef}[def]{4}}
 \proofstepwidestar[1,2,3]{R_S(\mathsf{and}(p,q))=R_S(p)\cap R_S(q)}{\rRE{\rUE{\rRE{\rUE{5},2}},3}}
\closeproofpartwideindR

\proofpartwideindR[Fasergleichung des Existenzquantors]{\mathsf{Erf}(M,S),\ i\in\mathbb N,\ p\in\mathcal F\vdash R_S(\mathsf{ex}_i(p))=Q_i(R_S(p))}{\mathsf{Erf}(M,S),\ i\in\mathbb N,\ p\in\mathcal F\vdash R_S(\mathsf{ex}_i(p))=Q_i(R_S(p))}[B48ErfExistsRow]
 \proofstepwidestar[1]{\mathsf{Erf}(M,S)}{\rA}
 \proofstepwidestar[2]{i\in\mathbb N}{\rA}
 \proofstepwidestar[3]{p\in\mathcal F}{\rA}
 \proofstepwidestar[1]{C_{\rm q}(M,S)}{\FormulaRefAuto{B48ErfExists}[thm]{1}}
 \proofstepwidestar[1]{\forall j\in\mathbb N\,\forall q\in\mathcal F\;(R_S(\mathsf{ex}_j(q))=Q_j(R_S(q)))}{\FormulaRefAuto{B48SatisfactionExistsDef}[def]{4}}
 \proofstepwidestar[1,2,3]{R_S(\mathsf{ex}_i(p))=Q_i(R_S(p))}{\rRE{\rUE{\rRE{\rUE{5},2}},3}}
\closeproofpartwideindR

\proofpartwideindR[Elementkriterium des Gleichheitstraegers]{p\in K(S,T)\eqvdash p\in\mathcal F\land R_S(p)=R_T(p)}{p\in K(S,T)\eqvdash p\in\mathcal F\land R_S(p)=R_T(p)}[B48EqualityCarrierMember]
 \proofstepwidestar[]{p\in K(S,T)\leftrightarrow p\in\mathcal F\land R_S(p)=R_T(p)}{\FormulaRefAuto{B48EqualityCarrier}[def]}
\closeproofpartwideindR

\proofpartwideindR[Gleichheitsinduktion bei Atomen]{\mathsf{Erf}(M,S),\ \mathsf{Erf}(M,T),\ c\in\Sigma_{\rm at}\vdash\mathsf l(c)\in K(S,T)}{\mathsf{Erf}(M,S),\ \mathsf{Erf}(M,T),\ c\in\Sigma_{\rm at}\vdash\mathsf l(c)\in K(S,T)}[B48EqualityAtom]
 \proofstepwidestar[1]{\mathsf{Erf}(M,S)}{\rA}
 \proofstepwidestar[2]{\mathsf{Erf}(M,T)}{\rA}
 \proofstepwidestar[3]{c\in\Sigma_{\rm at}}{\rA}
 \proofstepwidestar[3]{\mathsf l(c)\in\mathcal F}{\FormulaRefAuto{B48FormulaAtom}[thm]{3}}
 \proofstepwidestar[1,3]{R_S(\mathsf l(c))=B_M(c)}{\FormulaRefAuto{B48ErfAtomRow}[thm]{1,3}}
 \proofstepwidestar[2,3]{R_T(\mathsf l(c))=B_M(c)}{\FormulaRefAuto{B48ErfAtomRow}[thm]{2,3}}
 \proofstepwidestar[1,2,3]{R_S(\mathsf l(c))=R_T(\mathsf l(c))}{\FormulaRefAuto{a=b,\ c=b\vdash a=c}[thm]{5,6}}
 \proofstepwidestar[1,2,3]{\mathsf l(c)\in K(S,T)}{\FormulaRefAuto{B48EqualityCarrierMember}[thm]{\rAI{4,7}}}
\closeproofpartwideindR

\proofpartwideindR[Gleichheitsinduktion bei Negation]{\mathsf{Erf}(M,S),\ \mathsf{Erf}(M,T),\ p\in K(S,T)\vdash\mathsf{neg}(p)\in K(S,T)}{\mathsf{Erf}(M,S),\ \mathsf{Erf}(M,T),\ p\in K(S,T)\vdash\mathsf{neg}(p)\in K(S,T)}[B48EqualityNeg]
 \proofstepwidestar[1]{\mathsf{Erf}(M,S)}{\rA}
 \proofstepwidestar[2]{\mathsf{Erf}(M,T)}{\rA}
 \proofstepwidestar[3]{p\in K(S,T)}{\rA}
 \proofstepwidestar[3]{p\in\mathcal F\land R_S(p)=R_T(p)}{\FormulaRefAuto{B48EqualityCarrierMember}[thm]{3}}
 \proofstepwidestar[3]{p\in\mathcal F}{\rAEa{4}}
 \proofstepwidestar[3]{R_S(p)=R_T(p)}{\rAEb{4}}
 \proofstepwidestar[3]{\mathsf{neg}(p)\in\mathcal F}{\FormulaRefAuto{B48FormulaNeg}[thm]{5}}
 \proofstepwidestar[1,3]{R_S(\mathsf{neg}(p))=A_D\setminus R_S(p)}{\FormulaRefAuto{B48ErfNegRow}[thm]{1,5}}
 \proofstepwidestar[2,3]{R_T(\mathsf{neg}(p))=A_D\setminus R_T(p)}{\FormulaRefAuto{B48ErfNegRow}[thm]{2,5}}
 \proofstepwidestar[1,3]{R_S(\mathsf{neg}(p))=A_D\setminus R_T(p)}{\rIE{6,8}}
 \proofstepwidestar[1,2,3]{R_S(\mathsf{neg}(p))=R_T(\mathsf{neg}(p))}{\FormulaRefAuto{a=b,\ c=b\vdash a=c}[thm]{10,9}}
 \proofstepwidestar[1,2,3]{\mathsf{neg}(p)\in K(S,T)}{\FormulaRefAuto{B48EqualityCarrierMember}[thm]{\rAI{7,11}}}
\closeproofpartwideindR

\proofpartwideindR[Gleichheitsinduktion bei Konjunktion]{\mathsf{Erf}(M,S),\ \mathsf{Erf}(M,T),\ p\in K(S,T),\ q\in K(S,T)\vdash\mathsf{and}(p,q)\in K(S,T)}{\mathsf{Erf}(M,S),\ \mathsf{Erf}(M,T),\ p\in K(S,T),\ q\in K(S,T)\vdash\mathsf{and}(p,q)\in K(S,T)}[B48EqualityAnd]
 \proofstepwidestar[1]{\mathsf{Erf}(M,S)}{\rA}
 \proofstepwidestar[2]{\mathsf{Erf}(M,T)}{\rA}
 \proofstepwidestar[3]{p\in K(S,T)}{\rA}
 \proofstepwidestar[4]{q\in K(S,T)}{\rA}
 \proofstepwidestar[3]{p\in\mathcal F\land R_S(p)=R_T(p)}{\FormulaRefAuto{B48EqualityCarrierMember}[thm]{3}}
 \proofstepwidestar[4]{q\in\mathcal F\land R_S(q)=R_T(q)}{\FormulaRefAuto{B48EqualityCarrierMember}[thm]{4}}
 \proofstepwidestar[3,4]{\mathsf{and}(p,q)\in\mathcal F}{\FormulaRefAuto{B48FormulaAnd}[thm]{\rAEa{5},\rAEa{6}}}
 \proofstepwidestar[1,3,4]{R_S(\mathsf{and}(p,q))=R_S(p)\cap R_S(q)}{\FormulaRefAuto{B48ErfAndRow}[thm]{1,\rAEa{5},\rAEa{6}}}
 \proofstepwidestar[2,3,4]{R_T(\mathsf{and}(p,q))=R_T(p)\cap R_T(q)}{\FormulaRefAuto{B48ErfAndRow}[thm]{2,\rAEa{5},\rAEa{6}}}
 \proofstepwidestar[1,3,4]{R_S(\mathsf{and}(p,q))=R_T(p)\cap R_T(q)}{\rIE{\rAEb{6},\rIE{\rAEb{5},8}}}
 \proofstepwidestar[1,2,3,4]{R_S(\mathsf{and}(p,q))=R_T(\mathsf{and}(p,q))}{\FormulaRefAuto{a=b,\ c=b\vdash a=c}[thm]{10,9}}
 \proofstepwidestar[1,2,3,4]{\mathsf{and}(p,q)\in K(S,T)}{\FormulaRefAuto{B48EqualityCarrierMember}[thm]{\rAI{7,11}}}
\closeproofpartwideindR

\proofpartwideindR[Gleichheitsinduktion beim Existenzquantor]{\mathsf{Erf}(M,S),\ \mathsf{Erf}(M,T),\ i\in\mathbb N,\ p\in K(S,T)\vdash\mathsf{ex}_i(p)\in K(S,T)}{\mathsf{Erf}(M,S),\ \mathsf{Erf}(M,T),\ i\in\mathbb N,\ p\in K(S,T)\vdash\mathsf{ex}_i(p)\in K(S,T)}[B48EqualityExists]
 \proofstepwidestar[1]{\mathsf{Erf}(M,S)}{\rA}
 \proofstepwidestar[2]{\mathsf{Erf}(M,T)}{\rA}
 \proofstepwidestar[3]{i\in\mathbb N}{\rA}
 \proofstepwidestar[4]{p\in K(S,T)}{\rA}
 \proofstepwidestar[4]{p\in\mathcal F\land R_S(p)=R_T(p)}{\FormulaRefAuto{B48EqualityCarrierMember}[thm]{4}}
 \proofstepwidestar[3,4]{\mathsf{ex}_i(p)\in\mathcal F}{\FormulaRefAuto{B48FormulaExists}[thm]{3,\rAEa{5}}}
 \proofstepwidestar[1,3,4]{R_S(\mathsf{ex}_i(p))=Q_i(R_S(p))}{\FormulaRefAuto{B48ErfExistsRow}[thm]{1,3,\rAEa{5}}}
 \proofstepwidestar[2,3,4]{R_T(\mathsf{ex}_i(p))=Q_i(R_T(p))}{\FormulaRefAuto{B48ErfExistsRow}[thm]{2,3,\rAEa{5}}}
 \proofstepwidestar[1,3,4]{R_S(\mathsf{ex}_i(p))=Q_i(R_T(p))}{\rIE{\rAEb{5},7}}
 \proofstepwidestar[1,2,3,4]{R_S(\mathsf{ex}_i(p))=R_T(\mathsf{ex}_i(p))}{\FormulaRefAuto{a=b,\ c=b\vdash a=c}[thm]{9,8}}
 \proofstepwidestar[1,2,3,4]{\mathsf{ex}_i(p)\in K(S,T)}{\FormulaRefAuto{B48EqualityCarrierMember}[thm]{\rAI{6,10}}}
\closeproofpartwideindR

\proofpartwideindR[Der Gleichheitstraeger ist abgeschlossen]{\mathsf{Erf}(M,S),\ \mathsf{Erf}(M,T)\vdash\mathsf{Cl}(K(S,T))}{\mathsf{Erf}(M,S),\ \mathsf{Erf}(M,T)\vdash\mathsf{Cl}(K(S,T))}[B48EqualityClosed]
 \proofstepwidestar[1]{\mathsf{Erf}(M,S)}{\rA}
 \proofstepwidestar[2]{\mathsf{Erf}(M,T)}{\rA}
 \proofstepwidestar[]{K(S,T)\subseteq\mathcal F}{\FormulaRefAuto{\{x\in A\mid P(x)\}\subseteq A}[thm]}
 \proofstepwidestar[]{\mathcal F\subseteq\mathcal T}{\FormulaRefAuto{\{x\in A\mid P(x)\}\subseteq A}[thm]}
 \proofstepwidestar[]{K(S,T)\subseteq\mathcal T}{\FormulaRefAuto{A\subseteq B,\ B\subseteq C\vdash A\subseteq C}[thm]{3,4}}
 \proofstepwidestar[6]{c\in\Sigma_{\rm at}}{\rA}
 \proofstepwidestar[1,2,6]{\mathsf l(c)\in K(S,T)}{\FormulaRefAuto{B48EqualityAtom}[thm]{1,2,6}}
 \proofstepwidestar[1,2]{\forall c\in\Sigma_{\rm at}\;(\mathsf l(c)\in K(S,T))}{\rUI{\rRI{6,7}}}
 \proofstepwidestar[9]{p\in K(S,T)}{\rA}
 \proofstepwidestar[1,2,9]{\mathsf{neg}(p)\in K(S,T)}{\FormulaRefAuto{B48EqualityNeg}[thm]{1,2,9}}
 \proofstepwidestar[1,2]{\forall p\in K(S,T)\;(\mathsf{neg}(p)\in K(S,T))}{\rUI{\rRI{9,10}}}
 \proofstepwidestar[12]{p\in K(S,T)}{\rA}
 \proofstepwidestar[13]{q\in K(S,T)}{\rA}
 \proofstepwidestar[1,2,12,13]{\mathsf{and}(p,q)\in K(S,T)}{\FormulaRefAuto{B48EqualityAnd}[thm]{1,2,12,13}}
 \proofstepwidestar[1,2]{\forall p\in K(S,T)\,\forall q\in K(S,T)\;(\mathsf{and}(p,q)\in K(S,T))}{\rUI{\rRI{12,\rUI{\rRI{13,14}}}}}
 \proofstepwidestar[16]{i\in\mathbb N}{\rA}
 \proofstepwidestar[17]{p\in K(S,T)}{\rA}
 \proofstepwidestar[1,2,16,17]{\mathsf{ex}_i(p)\in K(S,T)}{\FormulaRefAuto{B48EqualityExists}[thm]{1,2,16,17}}
 \proofstepwidestar[1,2]{\forall i\in\mathbb N\,\forall p\in K(S,T)\;(\mathsf{ex}_i(p)\in K(S,T))}{\rUI{\rRI{16,\rUI{\rRI{17,18}}}}}
 \proofstepwidestar[1,2]{\mathsf{Cl}(K(S,T))}{\FormulaRefAuto{B48FormulaClosure}[def]{5,8,11,15,19}}
\closeproofpartwideindR

\proofpartwideindR[Eindeutigkeit aller Formelfasern]{\mathsf{Erf}(M,S),\ \mathsf{Erf}(M,T)\vdash\forall p\in\mathcal F\;(R_S(p)=R_T(p))}{\mathsf{Erf}(M,S),\ \mathsf{Erf}(M,T)\vdash\forall p\in\mathcal F\;(R_S(p)=R_T(p))}[B48EqualityAllRows]
 \proofstepwidestar[1]{\mathsf{Erf}(M,S)}{\rA}
 \proofstepwidestar[2]{\mathsf{Erf}(M,T)}{\rA}
 \proofstepwidestar[1,2]{\mathsf{Cl}(K(S,T))}{\FormulaRefAuto{B48EqualityClosed}[thm]{1,2}}
 \proofstepwidestar[4]{p\in\mathcal F}{\rA}
 \proofstepwidestar[1,2,4]{p\in K(S,T)}{\FormulaRefAuto{B48FormulaMinimalMember}[thm]{3,4}}
 \proofstepwidestar[1,2,4]{p\in\mathcal F\land R_S(p)=R_T(p)}{\FormulaRefAuto{B48EqualityCarrierMember}[thm]{5}}
 \proofstepwidestar[1,2,4]{R_S(p)=R_T(p)}{\rAEb{6}}
 \proofstepwidestar[1,2]{\forall p\in\mathcal F\;(R_S(p)=R_T(p))}{\rUI{\rRI{4,7}}}
\closeproofpartwideindR

\proofpartwideindR[Gleiche Fasern liefern eine Inklusion]{S\subseteq\mathcal F\times A_D,\ \forall p\in\mathcal F\;(R_S(p)=R_T(p))\vdash S\subseteq T}{S\subseteq\mathcal F\times A_D,\ \forall p\in\mathcal F\;(R_S(p)=R_T(p))\vdash S\subseteq T}[B48RowsSubset]
 \proofstepwidestar[1]{S\subseteq\mathcal F\times A_D}{\rA}
 \proofstepwidestar[2]{\forall p\in\mathcal F\;(R_S(p)=R_T(p))}{\rA}
 \proofstepwidestar[3]{z\in S}{\rA}
 \proofstepwidestar[1,3]{z\in\mathcal F\times A_D}{\FormulaRefAuto{A\subseteq B,\ x\in A\vdash x\in B}[thm]{1,3}}
 \proofstepwidestar[1,3]{\exists p\in\mathcal F\,\exists s\in A_D\;(z=(p,s))}{\FormulaRefAuto{x\in A\times B\eqvdash\exists a\in A\,\exists b\in B\,x=(a,b)}[thm]{4}}
 \proofstepwidestar[6]{p\in\mathcal F\land(s\in A_D\land z=(p,s))}{\rA}
 \proofstepwidestar[6]{p\in\mathcal F}{\rAEa{6}}
 \proofstepwidestar[6]{s\in A_D}{\rAEa{\rAEb{6}}}
 \proofstepwidestar[6]{z=(p,s)}{\rAEb{\rAEb{6}}}
 \proofstepwidestar[3,6]{(p,s)\in S}{\rIE{9,3}}
 \proofstepwidestar[3,6]{s\in R_S(p)}{\FormulaRefAuto{B48SemanticTerms}[def]{\rAI{8,10}}}
 \proofstepwidestar[2,6]{R_S(p)=R_T(p)}{\rRE{\rUE{2},7}}
 \proofstepwidestar[2,3,6]{s\in R_T(p)}{\rIE{12,11}}
 \proofstepwidestar[2,3,6]{s\in A_D\land(p,s)\in T}{\FormulaRefAuto{B48SemanticTerms}[def]{13}}
 \proofstepwidestar[2,3,6]{(p,s)\in T}{\rAEb{14}}
 \proofstepwidestar[6]{(p,s)=z}{\FormulaRefAuto{a=b\vdash b=a}[thm]{9}}
 \proofstepwidestar[2,3,6]{z\in T}{\rIE{16,15}}
 \proofstepwidestar[1,2,3]{z\in T}{\rEEStar{5,6:17}}
 \proofstepwidestar[1,2]{S\subseteq T}{\FormulaRefAuto{SubsetIntroductionRule}[thm]{[3]\vdots 18}}
\closeproofpartwideindR

\proofpartwideindR[Eindeutigkeit der Erfuellungsrelation]{\mathsf{Erf}(M,S),\ \mathsf{Erf}(M,T)\vdash S=T}{\mathsf{Erf}(M,S),\ \mathsf{Erf}(M,T)\vdash S=T}[B48SatisfactionUnique]
 \proofstepwidestar[1]{\mathsf{Erf}(M,S)}{\rA}
 \proofstepwidestar[2]{\mathsf{Erf}(M,T)}{\rA}
 \proofstepwidestar[1]{S\subseteq\mathcal F\times A_D}{\FormulaRefAuto{B48ErfSubset}[thm]{1}}
 \proofstepwidestar[2]{T\subseteq\mathcal F\times A_D}{\FormulaRefAuto{B48ErfSubset}[thm]{2}}
 \proofstepwidestar[1,2]{\forall p\in\mathcal F\;(R_S(p)=R_T(p))}{\FormulaRefAuto{B48EqualityAllRows}[thm]{1,2}}
 \proofstepwidestar[1,2]{\forall p\in\mathcal F\;(R_T(p)=R_S(p))}{\FormulaRefAuto{B48EqualityAllRows}[thm]{2,1}}
 \proofstepwidestar[1,2]{S\subseteq T}{\FormulaRefAuto{B48RowsSubset}[thm]{3,5}}
 \proofstepwidestar[1,2]{T\subseteq S}{\FormulaRefAuto{B48RowsSubset}[thm]{4,6}}
 \proofstepwidestar[1,2]{S=T}{\FormulaRefAuto{A\subseteq B,\ B\subseteq A\vdash A=B}[thm]{7,8}}
\closeproofpartwideindR

\par\medskip\noindent\textbf{G. Punktweise Wahrheitsklauseln}\par
\proofpartwideindR[Umkehrung einer Aequivalenz]{P\leftrightarrow Q\vdash Q\leftrightarrow P}{P\leftrightarrow Q\vdash Q\leftrightarrow P}[B48IffSymmetry]
 \proofstepwidestar[1]{P\leftrightarrow Q}{\rA}
 \proofstepwidestar[2]{Q}{\rA}
 \proofstepwidestar[1,2]{P}{\rLREb{1,2}}
 \proofstepwidestar[1]{Q\to P}{\rRI{2,3}}
 \proofstepwidestar[5]{P}{\rA}
 \proofstepwidestar[1,5]{Q}{\rLREa{1,5}}
 \proofstepwidestar[1]{P\to Q}{\rRI{5,6}}
 \proofstepwidestar[1]{Q\leftrightarrow P}{\rLRI{4,7}}
\closeproofpartwideindR

\proofpartwideindR[Aequivalente beschraenkte Existenzbehauptungen]{\forall x\in A\;(P(x)\leftrightarrow Q(x))\vdash((\exists x\in A\;P(x))\leftrightarrow(\exists x\in A\;Q(x)))}{\forall x\in A\;(P(x)\leftrightarrow Q(x))\vdash((\exists x\in A\;P(x))\leftrightarrow(\exists x\in A\;Q(x)))}[B48ExistsIff]
 \proofstepwidestar[1]{\forall x\in A\;(P(x)\leftrightarrow Q(x))}{\rA}
 \proofstepwidestar[2]{\exists x\in A\;P(x)}{\rA}
 \proofstepwidestar[3]{a\in A\land P(a)}{\rA}
 \proofstepwidestar[1,3]{P(a)\leftrightarrow Q(a)}{\rRE{\rUE{1},\rAEa{3}}}
 \proofstepwidestar[1,3]{Q(a)}{\rLREa{4,\rAEb{3}}}
 \proofstepwidestar[1,3]{\exists x\in A\;Q(x)}{\rEI{\rAI{\rAEa{3},5}}}
 \proofstepwidestar[1,2]{\exists x\in A\;Q(x)}{\rEE{2,3:6}}
 \proofstepwidestar[1]{(\exists x\in A\;P(x))\to(\exists x\in A\;Q(x))}{\rRI{2,7}}
 \proofstepwidestar[9]{\exists x\in A\;Q(x)}{\rA}
 \proofstepwidestar[10]{b\in A\land Q(b)}{\rA}
 \proofstepwidestar[1,10]{P(b)\leftrightarrow Q(b)}{\rRE{\rUE{1},\rAEa{10}}}
 \proofstepwidestar[1,10]{P(b)}{\rLREb{11,\rAEb{10}}}
 \proofstepwidestar[1,10]{\exists x\in A\;P(x)}{\rEI{\rAI{\rAEa{10},12}}}
 \proofstepwidestar[1,9]{\exists x\in A\;P(x)}{\rEE{9,10:13}}
 \proofstepwidestar[1]{(\exists x\in A\;Q(x))\to(\exists x\in A\;P(x))}{\rRI{9,14}}
 \proofstepwidestar[1]{(\exists x\in A\;P(x))\leftrightarrow(\exists x\in A\;Q(x))}{\rLRI{8,15}}
\closeproofpartwideindR

\proofpartwideindR[Faserzugehoerigkeit und Relationszugehoerigkeit]{s\in A_D\vdash(s\in R_S(p)\leftrightarrow(p,s)\in S)}{s\in A_D\vdash(s\in R_S(p)\leftrightarrow(p,s)\in S)}[B48RowMember]
 \proofstepwidestar[1]{s\in A_D}{\rA}
 \proofstepwidestar[1]{s\in R_S(p)\leftrightarrow(p,s)\in S}{\FormulaRefAuto{B48SeparationWithin}[thm]{1}}
\closeproofpartwideindR

\proofpartwideindR[Punktweise atomare Erfuellungsklausel]{\mathsf{Erf}(M,S),\ c\in\Sigma_{\rm at},\ s\in A_D\vdash((\mathsf l(c),s)\in S\leftrightarrow s\in B_M(c))}{\mathsf{Erf}(M,S),\ c\in\Sigma_{\rm at},\ s\in A_D\vdash((\mathsf l(c),s)\in S\leftrightarrow s\in B_M(c))}[B48PointAtom]
 \proofstepwidestar[1]{\mathsf{Erf}(M,S)}{\rA}
 \proofstepwidestar[2]{c\in\Sigma_{\rm at}}{\rA}
 \proofstepwidestar[3]{s\in A_D}{\rA}
 \proofstepwidestar[1,2]{R_S(\mathsf l(c))=B_M(c)}{\FormulaRefAuto{B48ErfAtomRow}[thm]{1,2}}
 \proofstepwidestar[3]{s\in R_S(\mathsf l(c))\leftrightarrow(\mathsf l(c),s)\in S}{\FormulaRefAuto{B48RowMember}[thm]{3}}
 \proofstepwidestar[3]{(\mathsf l(c),s)\in S\leftrightarrow s\in R_S(\mathsf l(c))}{\FormulaRefAuto{B48IffSymmetry}[thm]{5}}
 \proofstepwidestar[1,2,3]{(\mathsf l(c),s)\in S\leftrightarrow s\in B_M(c)}{\rIE{4,6}}
\closeproofpartwideindR

\proofpartwideindR[Punktweise Negationsklausel]{\mathsf{Erf}(M,S),\ p\in\mathcal F,\ s\in A_D\vdash((\mathsf{neg}(p),s)\in S\leftrightarrow\neg((p,s)\in S))}{\mathsf{Erf}(M,S),\ p\in\mathcal F,\ s\in A_D\vdash((\mathsf{neg}(p),s)\in S\leftrightarrow\neg((p,s)\in S))}[B48PointNeg]
 \proofstepwidestar[1]{\mathsf{Erf}(M,S)}{\rA}
 \proofstepwidestar[2]{p\in\mathcal F}{\rA}
 \proofstepwidestar[3]{s\in A_D}{\rA}
 \proofstepwidestar[1,2]{R_S(\mathsf{neg}(p))=A_D\setminus R_S(p)}{\FormulaRefAuto{B48ErfNegRow}[thm]{1,2}}
 \proofstepwidestar[3]{s\in R_S(\mathsf{neg}(p))\leftrightarrow(\mathsf{neg}(p),s)\in S}{\FormulaRefAuto{B48RowMember}[thm]{3}}
 \proofstepwidestar[3]{(\mathsf{neg}(p),s)\in S\leftrightarrow s\in R_S(\mathsf{neg}(p))}{\FormulaRefAuto{B48IffSymmetry}[thm]{5}}
 \proofstepwidestar[1,2,3]{(\mathsf{neg}(p),s)\in S\leftrightarrow s\in A_D\setminus R_S(p)}{\rIE{4,6}}
 \proofstepwidestar[3]{s\in A_D\setminus R_S(p)\leftrightarrow\neg(s\in R_S(p))}{\FormulaRefAuto{B48SeparationWithin}[thm]{3}}
 \proofstepwidestar[1,2,3]{(\mathsf{neg}(p),s)\in S\leftrightarrow\neg(s\in R_S(p))}{\rChain{7,8}}
 \proofstepwidestar[3]{s\in R_S(p)\leftrightarrow(p,s)\in S}{\FormulaRefAuto{B48RowMember}[thm]{3}}
 \proofstepwidestar[1,2,3]{(\mathsf{neg}(p),s)\in S\leftrightarrow\neg((p,s)\in S)}{\rLRS{10,9}}
\closeproofpartwideindR

\proofpartwideindR[Punktweise Konjunktionsklausel]{\begin{gathered}\mathsf{Erf}(M,S),\ p\in\mathcal F,\ q\in\mathcal F,\ s\in A_D\\
 \vdash((\mathsf{and}(p,q),s)\in S\leftrightarrow((p,s)\in S\land(q,s)\in S))\end{gathered}}{\begin{gathered}\mathsf{Erf}(M,S),\ p\in\mathcal F,\ q\in\mathcal F,\ s\in A_D\\
 \vdash((\mathsf{and}(p,q),s)\in S\leftrightarrow((p,s)\in S\land(q,s)\in S))\end{gathered}}[B48PointAnd]
 \proofstepwidestar[1]{\mathsf{Erf}(M,S)}{\rA}
 \proofstepwidestar[2]{p\in\mathcal F}{\rA}
 \proofstepwidestar[3]{q\in\mathcal F}{\rA}
 \proofstepwidestar[4]{s\in A_D}{\rA}
 \proofstepwidestar[1,2,3]{R_S(\mathsf{and}(p,q))=R_S(p)\cap R_S(q)}{\FormulaRefAuto{B48ErfAndRow}[thm]{1,2,3}}
 \proofstepwidestar[4]{s\in R_S(\mathsf{and}(p,q))\leftrightarrow(\mathsf{and}(p,q),s)\in S}{\FormulaRefAuto{B48RowMember}[thm]{4}}
 \proofstepwidestar[4]{(\mathsf{and}(p,q),s)\in S\leftrightarrow s\in R_S(\mathsf{and}(p,q))}{\FormulaRefAuto{B48IffSymmetry}[thm]{6}}
 \proofstepwidestar[1,2,3,4]{(\mathsf{and}(p,q),s)\in S\leftrightarrow s\in R_S(p)\cap R_S(q)}{\rIE{5,7}}
 \proofstepwidestar[]{s\in R_S(p)\cap R_S(q)\leftrightarrow(s\in R_S(p)\land s\in R_S(q))}{\FormulaRefAuto{x\in A\cap B\eqvdash x\in A\land x\in B}[thm]}
 \proofstepwidestar[1,2,3,4]{(\mathsf{and}(p,q),s)\in S\leftrightarrow(s\in R_S(p)\land s\in R_S(q))}{\rChain{8,9}}
 \proofstepwidestar[4]{s\in R_S(p)\leftrightarrow(p,s)\in S}{\FormulaRefAuto{B48RowMember}[thm]{4}}
 \proofstepwidestar[4]{s\in R_S(q)\leftrightarrow(q,s)\in S}{\FormulaRefAuto{B48RowMember}[thm]{4}}
 \proofstepwidestar[1,2,3,4]{(\mathsf{and}(p,q),s)\in S\leftrightarrow((p,s)\in S\land(q,s)\in S)}{\rLRS{12,\rLRS{11,10}}}
\closeproofpartwideindR

\proofpartwideindR[Punktweise Quantorenklausel]{\begin{gathered}\mathsf{Erf}(M,S),\ i\in\mathbb N,\ p\in\mathcal F,\ s\in A_D\\
 \vdash((\mathsf{ex}_i(p),s)\in S\leftrightarrow
      \exists a\in D\;((p,s[i\mapsto a])\in S))\end{gathered}}{\begin{gathered}\mathsf{Erf}(M,S),\ i\in\mathbb N,\ p\in\mathcal F,\ s\in A_D\\
 \vdash((\mathsf{ex}_i(p),s)\in S\leftrightarrow
      \exists a\in D\;((p,s[i\mapsto a])\in S))\end{gathered}}[B48PointExists]
 \proofstepwidestar[1]{\mathsf{Erf}(M,S)}{\rA}
 \proofstepwidestar[2]{i\in\mathbb N}{\rA}
 \proofstepwidestar[3]{p\in\mathcal F}{\rA}
 \proofstepwidestar[4]{s\in A_D}{\rA}
 \proofstepwidestar[1,2,3]{R_S(\mathsf{ex}_i(p))=Q_i(R_S(p))}{\FormulaRefAuto{B48ErfExistsRow}[thm]{1,2,3}}
 \proofstepwidestar[4]{s\in R_S(\mathsf{ex}_i(p))\leftrightarrow(\mathsf{ex}_i(p),s)\in S}{\FormulaRefAuto{B48RowMember}[thm]{4}}
 \proofstepwidestar[4]{(\mathsf{ex}_i(p),s)\in S\leftrightarrow s\in R_S(\mathsf{ex}_i(p))}{\FormulaRefAuto{B48IffSymmetry}[thm]{6}}
 \proofstepwidestar[1,2,3,4]{(\mathsf{ex}_i(p),s)\in S\leftrightarrow s\in Q_i(R_S(p))}{\rIE{5,7}}
 \proofstepwidestar[4]{s\in Q_i(R_S(p))\leftrightarrow\exists a\in D\;(s[i\mapsto a]\in R_S(p))}{\FormulaRefAuto{B48SeparationWithin}[thm]{4}}
 \proofstepwidestar[10]{a\in D}{\rA}
 \proofstepwidestar[2,4,10]{s[i\mapsto a]\in A_D}{\FormulaRefAuto{B48AssignmentUpdate}[thm]{4,2,10}}
 \proofstepwidestar[2,4,10]{s[i\mapsto a]\in R_S(p)\leftrightarrow(p,s[i\mapsto a])\in S}{\FormulaRefAuto{B48RowMember}[thm]{11}}
 \proofstepwidestar[2,4]{\forall a\in D\;(s[i\mapsto a]\in R_S(p)\leftrightarrow(p,s[i\mapsto a])\in S)}{\rUI{\rRI{10,12}}}
 \proofstepwidestar[2,4]{(\exists a\in D\;s[i\mapsto a]\in R_S(p))\leftrightarrow(\exists a\in D\;(p,s[i\mapsto a])\in S)}{\FormulaRefAuto{B48ExistsIff}[thm]{13}}
 \proofstepwidestar[1,2,3,4]{(\mathsf{ex}_i(p),s)\in S\leftrightarrow\exists a\in D\;((p,s[i\mapsto a])\in S)}{\rChain{8,9,14}}
\closeproofpartwideindR

\proofpartwideindR[Atomarer Code fuer Gleichheit]{i\in\mathbb N,\ j\in\mathbb N\vdash (t_{\rm e},(i,j))\in\Sigma_{\rm at}}{i\in\mathbb N,\ j\in\mathbb N\vdash (t_{\rm e},(i,j))\in\Sigma_{\rm at}}[B48EqAtomic]
 \proofstepwidestar[1]{i\in\mathbb N}{\rA}
 \proofstepwidestar[2]{j\in\mathbb N}{\rA}
 \proofstepwidestar[]{t_{\rm e}\in\mathbb N}{\FormulaRefAuto{B48TagType1}[thm]}
 \proofstepwidestar[1,2]{(t_{\rm e},(i,j))\in\Sigma}{\FormulaRefAuto{B48LabelType}[thm]{3,1,2}}
 \proofstepwidestar[]{t_{\rm e}=t_{\rm e}}{\rII}
 \proofstepwidestar[]{t_{\rm e}=t_{\rm e}\lor(t_{\rm e}=t_{\rm m}\lor(t_{\rm e}=t_{\rm b}\land i=0\land j=0))}{\rOIa{5}}
 \proofstepwidestar[1,2]{(t_{\rm e},(i,j))\in\Sigma_{\rm at}}{\FormulaRefAuto{B48FormulaConstructors}[def]{\rAI{4,6}}}
\closeproofpartwideindR

\proofpartwideindR[Atomarer Code fuer Mitgliedschaft]{i\in\mathbb N,\ j\in\mathbb N\vdash (t_{\rm m},(i,j))\in\Sigma_{\rm at}}{i\in\mathbb N,\ j\in\mathbb N\vdash (t_{\rm m},(i,j))\in\Sigma_{\rm at}}[B48MemAtomic]
 \proofstepwidestar[1]{i\in\mathbb N}{\rA}
 \proofstepwidestar[2]{j\in\mathbb N}{\rA}
 \proofstepwidestar[]{t_{\rm m}\in\mathbb N}{\FormulaRefAuto{B48TagType2}[thm]}
 \proofstepwidestar[1,2]{(t_{\rm m},(i,j))\in\Sigma}{\FormulaRefAuto{B48LabelType}[thm]{3,1,2}}
 \proofstepwidestar[]{t_{\rm m}=t_{\rm m}}{\rII}
 \proofstepwidestar[]{t_{\rm m}=t_{\rm m}\lor(t_{\rm m}=t_{\rm b}\land i=0\land j=0)}{\rOIa{5}}
 \proofstepwidestar[]{t_{\rm m}=t_{\rm e}\lor(t_{\rm m}=t_{\rm m}\lor(t_{\rm m}=t_{\rm b}\land i=0\land j=0))}{\rOIb{6}}
 \proofstepwidestar[1,2]{(t_{\rm m},(i,j))\in\Sigma_{\rm at}}{\FormulaRefAuto{B48FormulaConstructors}[def]{\rAI{4,7}}}
\closeproofpartwideindR

\proofpartwideindR[Atomarer Widerspruchscode]{(t_{\rm b},(0,0))\in\Sigma_{\rm at}}{(t_{\rm b},(0,0))\in\Sigma_{\rm at}}[B48BottomAtomic]
 \proofstepwidestar[]{t_{\rm b}\in\mathbb N}{\FormulaRefAuto{B48TagType6}[thm]}
 \proofstepwidestar[]{0\in\mathbb N}{\FormulaRefAuto{PeanoZero}[ax]}
 \proofstepwidestar[]{(t_{\rm b},(0,0))\in\Sigma}{\FormulaRefAuto{B48LabelType}[thm]{1,2,2}}
 \proofstepwidestar[]{t_{\rm b}=t_{\rm b}}{\rII}
 \proofstepwidestar[]{0=0}{\rII}
 \proofstepwidestar[]{t_{\rm b}=t_{\rm b}\land(0=0\land0=0)}{\rAI{4,\rAI{5,5}}}
 \proofstepwidestar[]{t_{\rm b}=t_{\rm m}\lor(t_{\rm b}=t_{\rm b}\land0=0\land0=0)}{\rOIb{6}}
 \proofstepwidestar[]{t_{\rm b}=t_{\rm e}\lor(t_{\rm b}=t_{\rm m}\lor(t_{\rm b}=t_{\rm b}\land0=0\land0=0))}{\rOIb{7}}
 \proofstepwidestar[]{(t_{\rm b},(0,0))\in\Sigma_{\rm at}}{\FormulaRefAuto{B48FormulaConstructors}[def]{\rAI{3,8}}}
\closeproofpartwideindR

\proofpartwideindR[Wertemenge eines Gleichheitsatoms]{i\in\mathbb N,\ j\in\mathbb N\vdash B_M((t_{\rm e},(i,j)))=\{s\in A_D\mid s(i)=s(j)\}}{i\in\mathbb N,\ j\in\mathbb N\vdash B_M((t_{\rm e},(i,j)))=\{s\in A_D\mid s(i)=s(j)\}}[B48EqBValue]
 \proofstepwidestar[1]{i\in\mathbb N}{\rA}
 \proofstepwidestar[2]{j\in\mathbb N}{\rA}
 \proofstepwidestar[]{t_{\rm e}=t_{\rm e}}{\rII}
 \proofstepwidestar[]{\begin{gathered}\mathsf{if}\bigl(t_{\rm e}=t_{\rm e},\{s\in A_D\mid s(i)=s(j)\},\\
 \BXLVIIIIf{t_{\rm e}=t_{\rm m}}{\{s\in A_D\mid(s(i),s(j))\in E\}}{\varnothing}\bigr)\\
 =\{s\in A_D\mid s(i)=s(j)\}\end{gathered}}{\FormulaRefAuto{B48IfTrue}[thm]{3}}
 \proofstepwidestar[1,2]{B_M((t_{\rm e},(i,j)))=\{s\in A_D\mid s(i)=s(j)\}}{\FormulaRefAuto{B48SemanticTerms}[def]{4}}
\closeproofpartwideindR

\proofpartwideindR[Wertemenge eines Mitgliedschaftsatoms]{i\in\mathbb N,\ j\in\mathbb N\vdash B_M((t_{\rm m},(i,j)))=\{s\in A_D\mid(s(i),s(j))\in E\}}{i\in\mathbb N,\ j\in\mathbb N\vdash B_M((t_{\rm m},(i,j)))=\{s\in A_D\mid(s(i),s(j))\in E\}}[B48MemBValue]
 \proofstepwidestar[1]{i\in\mathbb N}{\rA}
 \proofstepwidestar[2]{j\in\mathbb N}{\rA}
 \proofstepwidestar[]{t_{\rm e}\neq t_{\rm m}}{\FormulaRefAuto{B48TagNe01}[thm]}
 \proofstepwidestar[]{t_{\rm m}\neq t_{\rm e}}{\FormulaRefAuto{a\neq b\vdash b\neq a}[thm]{3}}
 \proofstepwidestar[]{t_{\rm m}=t_{\rm m}}{\rII}
 \proofstepwidestar[]{\begin{gathered}\BXLVIIIIf{t_{\rm m}=t_{\rm m}}{\{s\in A_D\mid(s(i),s(j))\in E\}}{\varnothing}\\
 =\{s\in A_D\mid(s(i),s(j))\in E\}\end{gathered}}{\FormulaRefAuto{B48IfTrue}[thm]{5}}
 \proofstepwidestar[]{\begin{gathered}\mathsf{if}\bigl(t_{\rm m}=t_{\rm e},\{s\in A_D\mid s(i)=s(j)\},\\
 \BXLVIIIIf{t_{\rm m}=t_{\rm m}}{\{s\in A_D\mid(s(i),s(j))\in E\}}{\varnothing}\bigr)\\
 =\BXLVIIIIf{t_{\rm m}=t_{\rm m}}{\{s\in A_D\mid(s(i),s(j))\in E\}}{\varnothing}\end{gathered}}{\FormulaRefAuto{B48IfFalse}[thm]{4}}
 \proofstepwidestar[1,2]{\begin{gathered}B_M((t_{\rm m},(i,j)))\\
 =\BXLVIIIIf{t_{\rm m}=t_{\rm m}}{\{s\in A_D\mid(s(i),s(j))\in E\}}{\varnothing}\end{gathered}}{\FormulaRefAuto{B48SemanticTerms}[def]{7}}
 \proofstepwidestar[1,2]{B_M((t_{\rm m},(i,j)))=\{s\in A_D\mid(s(i),s(j))\in E\}}{\rIE{6,8}}
\closeproofpartwideindR

\proofpartwideindR[Wertemenge des Widerspruchscodes]{B_M((t_{\rm b},(0,0)))=\varnothing}{B_M((t_{\rm b},(0,0)))=\varnothing}[B48BottomBValue]
 \proofstepwidestar[]{t_{\rm e}\neq t_{\rm b}}{\FormulaRefAuto{B48TagNe05}[thm]}
 \proofstepwidestar[]{t_{\rm m}\neq t_{\rm b}}{\FormulaRefAuto{B48TagNe15}[thm]}
 \proofstepwidestar[]{t_{\rm b}\neq t_{\rm e}}{\FormulaRefAuto{a\neq b\vdash b\neq a}[thm]{1}}
 \proofstepwidestar[]{t_{\rm b}\neq t_{\rm m}}{\FormulaRefAuto{a\neq b\vdash b\neq a}[thm]{2}}
 \proofstepwidestar[]{\BXLVIIIIf{t_{\rm b}=t_{\rm m}}{\{s\in A_D\mid(s(0),s(0))\in E\}}{\varnothing}=\varnothing}{\FormulaRefAuto{B48IfFalse}[thm]{4}}
 \proofstepwidestar[]{\begin{gathered}\mathsf{if}\bigl(t_{\rm b}=t_{\rm e},\{s\in A_D\mid s(0)=s(0)\},\\
 \BXLVIIIIf{t_{\rm b}=t_{\rm m}}{\{s\in A_D\mid(s(0),s(0))\in E\}}{\varnothing}\bigr)\\
 =\BXLVIIIIf{t_{\rm b}=t_{\rm m}}{\{s\in A_D\mid(s(0),s(0))\in E\}}{\varnothing}\end{gathered}}{\FormulaRefAuto{B48IfFalse}[thm]{3}}
 \proofstepwidestar[]{B_M((t_{\rm b},(0,0)))=\BXLVIIIIf{t_{\rm b}=t_{\rm m}}{\{s\in A_D\mid(s(0),s(0))\in E\}}{\varnothing}}{\FormulaRefAuto{B48SemanticTerms}[def]{6}}
 \proofstepwidestar[]{B_M((t_{\rm b},(0,0)))=\varnothing}{\rIE{5,7}}
\closeproofpartwideindR

\proofpartwideindR[Punktweise Klausel fuer Gleichheit]{\begin{gathered}\mathsf{Erf}(M,S),\ i\in\mathbb N,\ j\in\mathbb N,\ s\in A_D\\\vdash((\mathsf e(i,j),s)\in S\leftrightarrow s(i)=s(j))\end{gathered}}{\begin{gathered}\mathsf{Erf}(M,S),\ i\in\mathbb N,\ j\in\mathbb N,\ s\in A_D\\\vdash((\mathsf e(i,j),s)\in S\leftrightarrow s(i)=s(j))\end{gathered}}[B48PointEq]
 \proofstepwidestar[1]{\mathsf{Erf}(M,S)}{\rA}
 \proofstepwidestar[2]{i\in\mathbb N}{\rA}
 \proofstepwidestar[3]{j\in\mathbb N}{\rA}
 \proofstepwidestar[4]{s\in A_D}{\rA}
 \proofstepwidestar[2,3]{(t_{\rm e},(i,j))\in\Sigma_{\rm at}}{\FormulaRefAuto{B48EqAtomic}[thm]{2,3}}
 \proofstepwidestar[1,2,3,4]{(\mathsf e(i,j),s)\in S\leftrightarrow s\in B_M((t_{\rm e},(i,j)))}{\FormulaRefAuto{B48PointAtom}[thm]{1,5,4}}
 \proofstepwidestar[2,3]{B_M((t_{\rm e},(i,j)))=\{s\in A_D\mid s(i)=s(j)\}}{\FormulaRefAuto{B48EqBValue}[thm]{2,3}}
 \proofstepwidestar[1,2,3,4]{(\mathsf e(i,j),s)\in S\leftrightarrow s\in\{u\in A_D\mid u(i)=u(j)\}}{\rIE{7,6}}
 \proofstepwidestar[4]{s\in\{u\in A_D\mid u(i)=u(j)\}\leftrightarrow s(i)=s(j)}{\FormulaRefAuto{B48SeparationWithin}[thm]{4}}
 \proofstepwidestar[1,2,3,4]{(\mathsf e(i,j),s)\in S\leftrightarrow s(i)=s(j)}{\rChain{8,9}}
\closeproofpartwideindR

\proofpartwideindR[Punktweise Klausel fuer Mitgliedschaft]{\begin{gathered}\mathsf{Erf}(M,S),\ i\in\mathbb N,\ j\in\mathbb N,\ s\in A_D\\\vdash((\mathsf m(i,j),s)\in S\leftrightarrow (s(i),s(j))\in E)\end{gathered}}{\begin{gathered}\mathsf{Erf}(M,S),\ i\in\mathbb N,\ j\in\mathbb N,\ s\in A_D\\\vdash((\mathsf m(i,j),s)\in S\leftrightarrow (s(i),s(j))\in E)\end{gathered}}[B48PointMem]
 \proofstepwidestar[1]{\mathsf{Erf}(M,S)}{\rA}
 \proofstepwidestar[2]{i\in\mathbb N}{\rA}
 \proofstepwidestar[3]{j\in\mathbb N}{\rA}
 \proofstepwidestar[4]{s\in A_D}{\rA}
 \proofstepwidestar[2,3]{(t_{\rm m},(i,j))\in\Sigma_{\rm at}}{\FormulaRefAuto{B48MemAtomic}[thm]{2,3}}
 \proofstepwidestar[1,2,3,4]{(\mathsf m(i,j),s)\in S\leftrightarrow s\in B_M((t_{\rm m},(i,j)))}{\FormulaRefAuto{B48PointAtom}[thm]{1,5,4}}
 \proofstepwidestar[2,3]{B_M((t_{\rm m},(i,j)))=\{s\in A_D\mid (s(i),s(j))\in E\}}{\FormulaRefAuto{B48MemBValue}[thm]{2,3}}
 \proofstepwidestar[1,2,3,4]{(\mathsf m(i,j),s)\in S\leftrightarrow s\in\{u\in A_D\mid (u(i),u(j))\in E\}}{\rIE{7,6}}
 \proofstepwidestar[4]{s\in\{u\in A_D\mid (u(i),u(j))\in E\}\leftrightarrow (s(i),s(j))\in E}{\FormulaRefAuto{B48SeparationWithin}[thm]{4}}
 \proofstepwidestar[1,2,3,4]{(\mathsf m(i,j),s)\in S\leftrightarrow (s(i),s(j))\in E}{\rChain{8,9}}
\closeproofpartwideindR

\proofpartwideindR[Keine Belegung erfuellt den Widerspruchscode]{\mathsf{Erf}(M,S),\ s\in A_D\vdash\neg((\mathsf b,s)\in S)}{\mathsf{Erf}(M,S),\ s\in A_D\vdash\neg((\mathsf b,s)\in S)}[B48PointBottom]
 \proofstepwidestar[1]{\mathsf{Erf}(M,S)}{\rA}
 \proofstepwidestar[2]{s\in A_D}{\rA}
 \proofstepwidestar[]{(t_{\rm b},(0,0))\in\Sigma_{\rm at}}{\FormulaRefAuto{B48BottomAtomic}[thm]}
 \proofstepwidestar[1,2]{(\mathsf b,s)\in S\leftrightarrow s\in B_M((t_{\rm b},(0,0)))}{\FormulaRefAuto{B48PointAtom}[thm]{1,3,2}}
 \proofstepwidestar[]{B_M((t_{\rm b},(0,0)))=\varnothing}{\FormulaRefAuto{B48BottomBValue}[thm]}
 \proofstepwidestar[1,2]{(\mathsf b,s)\in S\leftrightarrow s\in\varnothing}{\rIE{5,4}}
 \proofstepwidestar[7]{(\mathsf b,s)\in S}{\rA}
 \proofstepwidestar[1,2,7]{s\in\varnothing}{\rLREa{6,7}}
 \proofstepwidestar[]{s\notin\varnothing}{\FormulaRefAuto{x \not\in \varnothing}[thm]}
 \proofstepwidestar[1,2,7]{\bot}{\rBI{8,9}}
 \proofstepwidestar[1,2]{\neg((\mathsf b,s)\in S)}{\rCI{7,10}}
\closeproofpartwideindR


\proofpartwide{Schluss}
 \proofstepwidestar[1]{\mathsf{Str}(M)}{\rA}
 \proofstepwidestar[]{\exists e\;\mathsf{Rec}(e,M)}
   {\FormulaRefAuto{B48TreeEvaluationExists}[thm]}
 \proofstepwidestar[3]{\mathsf{Rec}(e,M)}{\rA}
 \proofstepwidestar[1,3]{\mathsf{Erf}(M,S_e)}
   {\FormulaRefAuto{B48CandidateSatisfies}[thm]{1,3}}
 \proofstepwidestar[1,3]{\exists S\;\mathsf{Erf}(M,S)}{\rEI{4}}
 \proofstepwidestar[1]{\exists S\;\mathsf{Erf}(M,S)}{\rEE{2,3:5}}
 \proofstepwidestar[7]{\mathsf{Erf}(M,S)}{\rA}
 \proofstepwidestar[8]{\mathsf{Erf}(M,T)}{\rA}
 \proofstepwidestar[7,8]{S=T}
   {\FormulaRefAuto{B48SatisfactionUnique}[thm]{7,8}}
 \proofstepwidestar[1]{\exists!S\;\mathsf{Erf}(M,S)}{\UEI{6,7,8,9}}
\closeproofpartwide
\end{tabproofsplitwide}



\subsection*{Die nun gerechtfertigte Schreibweise fuer Erfuellung}

Nach \FormulaRefAuto{B48SatisfactionExists}[thm] ist die folgende
Kennzeichnung wohldefiniert. Die rekursiven Klauseln werden hier als
gleichwertige definierende Gleichungen dieser eindeutigen Relation
mitgefuehrt, nicht als neue Existenzaxiome.
Die punktweisen Klauseln sind bereits hergeleitet.
Fuer Gleichheit und Elementrelation gelten\\
\FormulaRefAuto{B48PointEq}[thm] und \FormulaRefAuto{B48PointMem}[thm].

Fuer Negation und Konjunktion gelten\\
\FormulaRefAuto{B48PointNeg}[thm] und \FormulaRefAuto{B48PointAnd}[thm].

Fuer Existenzquantor und Widerspruch gelten\\
\FormulaRefAuto{B48PointExists}[thm] und \FormulaRefAuto{B48PointBottom}[thm].

\FormulaDefDeltaK[Erfuellung in einer Mengenstruktur]
{\begin{aligned}
 \operatorname{Sat}(M)&\coloneq\iota S\;\mathsf{Erf}(M,S),\\
 M,s\models p&\coloneq(p,s)\in\operatorname{Sat}(M),\\
 M,s\models v_i=v_j&\coloneq s(v_i)=s(v_j),\\
 M,s\models v_i\in v_j&\coloneq(s(v_i),s(v_j))\in E,\\
 M,s\models\neg p&\coloneq\neg(M,s\models p),\\
 M,s\models p\land q&\coloneq(M,s\models p)\land(M,s\models q),\\
 M,s\models\exists v_i\,p&\coloneq
          \exists a\in D\;(M,s[v_i\mapsto a]\models p),\\
 M,s\models\bot&\coloneq\bot.
\end{aligned}}
{B48Satisfaction}{
 \DeltaRow{Struktur}{M=(D,E),\quad D\neq\varnothing,\quad E\subseteq D\times D}
 \DeltaRow{Belegung}{s\colon\operatorname{Var}\to D}
 \DeltaRow{Formelcodes}{p,q\in\mathcal F}
 \DeltaRow{Variablencodes}{v_i=i,\quad v_j=j,\quad i,j\in\mathbb N}
}

Die punktweise Schreibweise ist damit auf eine konstruierte Menge
zurueckgefuehrt. Dies beweist noch nicht die Korrektheit beliebiger
syntaktischer Herleitungen: Dafuer werden spaeter zusaetzlich
Substitution, Koinzidenz und Herleitungsinduktion benoetigt.
